\documentclass[11pt]{article}
\usepackage[utf8]{inputenc}
\usepackage[T1]{fontenc}
\usepackage[a4paper,margin=1in]{geometry}
\usepackage{lmodern}
\usepackage{microtype}
\usepackage{amsmath,amssymb,mathtools}
\usepackage{mathrsfs}
\usepackage{tikz-cd}
\usepackage{enumitem}
\usepackage{longtable,booktabs,array}
\usepackage[hidelinks]{hyperref}
\setlength{\parindent}{0pt}
\setlength{\parskip}{0.55em}
\emergencystretch=3em
\hbadness=10000
\newcommand{\K}{K}
\newcommand{\Gr}{\operatorname{Gr}}
\newcommand{\Fil}{\operatorname{Fil}}
\newcommand{\ch}{\operatorname{ch}}
\newcommand{\Todd}{\operatorname{Todd}}
\newcommand{\cl}{\operatorname{cl}}
\newcommand{\Parf}{\operatorname{Parf}}
\newcommand{\Tor}{\operatorname{Tor}}
\newcommand{\Spec}{\operatorname{Spec}}
\newcommand{\RRR}{\mathrm{RRR}}
\providecommand{\codim}{\operatorname{codim}}
\providecommand{\Inv}{\operatorname{Inv}}
\providecommand{\Pic}{\operatorname{Pic}}
\providecommand{\Z}{\mathbb Z}
\providecommand{\Q}{\mathbb Q}
\providecommand{\Ahat}{\widehat A}
\providecommand{\Ahp}{\widehat A^{+}}
\providecommand{\That}{\widehat T}
\providecommand{\rad}{\operatorname{rad}}
\providecommand{\Gl}{\operatorname{Gl}}
\providecommand{\GL}{\operatorname{GL}}
\providecommand{\ch}{\operatorname{ch}}
\providecommand{\Todd}{\operatorname{Todd}}
\providecommand{\Gr}{\operatorname{Gr}}

% Batch current macro additions
\providecommand{\Atil}{\widetilde A}
\providecommand{\Ctil}{\widetilde C}
\providecommand{\K}{\mathrm K}
\providecommand{\cC}{\mathcal C}
\providecommand{\cF}{\mathscr F}
\providecommand{\eps}{\varepsilon}
\providecommand{\ellc}{\ell}
\providecommand{\Tor}{\operatorname{Tor}}
\providecommand{\Coker}{\operatorname{Coker}}
\providecommand{\Ker}{\operatorname{Ker}}
\providecommand{\Supp}{\operatorname{Supp}}
\providecommand{\codim}{\operatorname{codim}}
\providecommand{\Gr}{\operatorname{G}}
\providecommand{\rank}{\operatorname{rank}}
\providecommand{\rang}{\operatorname{rang}}
\providecommand{\cO}{\mathcal O}

\begin{document}
\begin{center}
{\Large Séminaire de Géométrie Algébrique du Bois Marie}\par
{\large 1966--67}\par
\vspace{1.5em}
{\Large Théorie des intersections et théorème de Riemann--Roch}\par
{\large (SGA 6)}\par
\vspace{1.5em}
un Séminaire dirigé par\par
P. Berthelot, A. Grothendieck, L. Illusie\par
avec la collaboration de\par
D. Ferrand, J.-P. Jouanolou, O. Jussila, S. Kleiman, M. Raynaud, J.-P. Serre
\end{center}

\section*{Préface}
La présente édition du Séminaire de Géométrie Algébrique 6 reproduit à peu de chose près la première édition, distribuée par l'I.H.E.S. sous forme d'exposés multigraphiés. Les corrections effectuées portent essentiellement sur des fautes d'impression, sauf dans les exposés I et III, où de légères modifications ont été apportées par L. Illusie à quelques énoncés incorrects dans la version primitive. D'autre part, des notes ont été rajoutées en bas de page, en particulier dans l'exposé XIV par A. Grothendieck, afin de signaler des progrès récents dans des questions encore ouvertes lors du Séminaire.

Signalons enfin que l'exposé XI, de D. Ferrand, concernant la théorie du déterminant pour les complexes parfaits, n'a pas été rédigé, et ne sera donc pas publié; pour éviter des erreurs de référence, nous avons gardé telle quelle la numérotation des exposés suivants.

\begin{flushright}Juin 1971\end{flushright}

\section*{Introduction}
Le but du présent Séminaire est de développer une théorie globale des intersections, et une formule de Riemann--Roch, pour des schémas quelconques. Le lecteur trouvera une description du programme du Séminaire dans l'exposé 0. La possibilité en principe d'une démonstration d'une formule de Riemann--Roch pour les schémas réguliers généraux, suivant les lignes du rapport bien connu de Borel--Serre en 1958, était claire dès ce moment, du moins pour un morphisme projectif. L'extension au cas général est moins évidente; le programme dans lequel elle s'insère, et partiellement réalisé dans notre séminaire, remonte à 1960. Comme c'était également le cas pour la théorie de dualité des faisceaux cohérents, un outil essentiel pour formuler une théorie satisfaisante est la théorie des catégories dérivées de Verdier, dont la connaissance est indispensable pour la compréhension du Séminaire.

À part cette théorie, l'étude du Séminaire n'exige guère qu'une connaissance générale des fondements de la géométrie algébrique, tels qu'ils sont exposés dans EGA, chapitres I, II, III; nous n'aurons en plus que des références occasionnelles à faire à EGA IV, pour certains faits concernant les morphismes plats ou lisses, qui pour l'essentiel figurent également dans les premiers exposés de SGA 1. Enfin, pour développer certains résultats avec toute la généralité souhaitable pour les applications, nous faisons usage parfois de la notion de site annelé et de topos annelé, pour laquelle nous renvoyons à SGA 4, exposés I à IV. Le lecteur qui ignorerait le langage des sites et topos pourra remplacer partout lesdits objets par des espaces topologiques ordinaires, les objets du topos étant alors remplacés par des ouverts de ces espaces; mais nous lui conseillons néanmoins, de préférence, de s'assimiler le langage des topos, qui fournit un principe d'unification extrêmement commode.

La notion fondamentale pour la théorie présentée ici est celle de complexe de Modules parfait, et ses diverses variantes, développées dans les exposés I, II, III. Il semble clair que ces notions, importantes également dans d'autres contextes en géométrie algébrique, notamment pour les formules de type Lefschetz--Verdier en cohomologie à coefficients discrets (SGA 5) ou cohérents, auront aussi leur rôle à jouer en dehors de la géométrie algébrique, notamment pour la formulation d'une variante analytique complexe du théorème de Riemann--Roch présenté ici, ou de variantes convenables du théorème de l'index d'Atiyah--Singer. Quelques indications dans ce sens seront données dans l'exposé II. Malheureusement, il manque encore à l'heure actuelle un énoncé, même heuristique, qui engloberait ces deux théorèmes, dont le premier pour l'instant reste conjectural. Il manque apparemment une idée nouvelle, comme il en manque aussi en géométrie algébrique pour parvenir à une démonstration du théorème de Riemann--Roch en dehors d'hypothèses projectives (cf. Exp. XIV, no 2).

Nous n'aurons à faire nul usage dans ce séminaire de la théorie locale des intersections, exposée dans J.-P. Serre, \emph{Algèbre locale. Multiplicités} (Lecture Notes in Mathematics no 11, Springer), et nous contenterons simplement de signaler ici que ce cours de Serre a eu une influence évidente sur la genèse de la théorie en 1957. Nous n'utiliserons pas non plus la théorie de l'anneau de Chow développée dans le Séminaire Chevalley 1958, exposés 2 et 3 (École Normale Supérieure). Cette théorie est liée de façon essentielle à des hypothèses de non singularité, alors que le but de notre séminaire est au contraire de développer une théorie des intersections sur les schémas généraux, et même les topos localement annelés généraux; voir à ce sujet les commentaires dans Exp. XIV nos 4 et 8, donnant les relations entre notre théorie et celle de l'anneau de Chow, et posant la question d'une généralisation de cette dernière. Nous pouvons dire que le Séminaire présente une théorie des intersections cohérente et se suffisant à elle-même, mais nullement exhaustive des différents points de vue utiles, voire indispensables, en théorie des intersections, et qu'il convient par suite de le compléter par les exposés cités de Serre et de Chevalley.

Nous avons joint au Séminaire, en appendice à l'exposé 0, le rapport de Grothendieck de 1957, «Classes de faisceaux et théorème de Riemann--Roch», qui avait servi de base au séminaire de Borel et Serre à Princeton la même année, ainsi qu'à leur article déjà cité. Ce rapport esquisse deux démonstrations du théorème de Riemann--Roch pour les variétés quasi-projectives non singulières, dont la première, valable pour le moment seulement en caractéristique nulle, mais donnant en revanche un résultat un peu plus précis dans le cas d'une immersion, ne figure pas encore dans la littérature, sans exclure le travail de Borel--Serre ni le présent séminaire. La lecture de ce rapport ne présuppose bien entendu celle d'aucun autre exposé du séminaire, et peut même servir d'introduction à l'étude de ce dernier au même titre que l'exposé 0, surtout pour un lecteur qui ne serait pas encore familier avec Borel--Serre. La démonstration à laquelle nous venons de faire allusion s'applique d'ailleurs essentiellement telle quelle au cas d'un morphisme projectif d'espaces analytiques complexes, et s'appliquera sans doute également en caractéristique quelconque, une fois résolu le problème des opérations «puissances extérieures» dans la catégorie dérivée (cf. Exp. XIV, nos 1, 2, 3). C'est l'absence d'une étude de cette question qui constitue sans doute la lacune la plus sérieuse dans ce Séminaire, dans l'optique même où nous nous y sommes placés.

On remarquera l'absence, dans la table des matières du présent Séminaire, de deux des participants mentionnés sur la page de garde. L'un, J.-P. Jouanolou, avait pris une part active à l'élaboration technique de la première partie du séminaire, mais avait été empêché de prendre part aux exposés oraux; on lui doit notamment l'assertion d'indépendance linéaire contenue dans l'important énoncé VI 1.1 donnant la structure de l'anneau $K^\circ$ des classes de Modules localement libres sur un fibré projectif, qui n'était prouvé auparavant que lorsqu'on supposait que le schéma de base admet un Module inversible ample. L'autre, J.-P. Serre, avait fait deux exposés oraux, logiquement indépendants du reste du Séminaire, et qu'il a par suite préféré publier sous forme d'article séparé (cf. J.-P. Serre, «Groupes de Grothendieck des schémas en groupes réductifs déployés», à paraître dans \emph{Publications Mathématiques}, no 34).

Comme dans les autres parties du Séminaire de Géométrie Algébrique du Bois Marie, les sigles SGA 1 à SGA 6 renvoient aux différentes parties dudit Séminaire, le chiffre romain suivant ce sigle indiquant le numéro de l'exposé, et les chiffres arabes qui le suivent correspondant à la numérotation intérieure de l'exposé en question; pour les références intérieures au présent Séminaire, on omet le sigle SGA 6, en utilisant des références du style I 4.9, signifiant: exposé I, énoncé 4.9. Le sigle EGA réfère aux \emph{Éléments de Géométrie Algébrique} de J. Dieudonné et A. Grothendieck.

\section*{Table des matières}
\begin{longtable}{@{}p{0.16\linewidth}p{0.64\linewidth}r@{}}
\toprule
Exposé & Titre et auteur & Page \\
\midrule
0 & Esquisse d'un programme pour une théorie des intersections sur les schémas généraux, par A. Grothendieck & 1 \\
Appendice & Classes de faisceaux et théorème de Riemann--Roch, par A. Grothendieck & 20 \\
I & Généralités sur les conditions de finitude dans les catégories dérivées, par L. Illusie & 56 \\
II & Existence de résolutions globales, par L. Illusie & 160 \\
III & Conditions de finitude relatives, par L. Illusie & 190 \\
IV & Groupes de Grothendieck des topos annelés, par L. Illusie & 222 \\
V & Généralités sur les $\lambda$-anneaux, par P. Berthelot & 297 \\
VI & Le $K^\circ$ d'un fibré projectif: calculs et conséquences, par P. Berthelot & 365 \\
VII & Immersions régulières et calcul de $K^\circ$ d'un éclatement, par P. Berthelot & 416 \\
VIII & Le théorème de Riemann--Roch, par P. Berthelot & 466 \\
IX & Quelques calculs de groupes $K_\bullet$, par P. Berthelot & 498 \\
X & Formalisme des intersections sur les schémas algébriques propres, par O. Jussila, avec un appendice par A. Grothendieck, «Spécialisation en théorie des intersections» & 519 \\
XI & Non rédigé & -- \\
XII & Un théorème de représentabilité relative sur le foncteur de Picard, par M. Raynaud, rédigé par S. Kleiman & 595 \\
XIII & Les théorèmes de finitude pour le foncteur de Picard, par S. Kleiman & 616 \\
XIV & Problèmes ouverts en théorie des intersections, par A. Grothendieck & 667 \\
\bottomrule
\end{longtable}

\section*{Exposé 0}
\begin{center}
{\large Esquisse d'un programme pour une théorie des intersections sur les schémas généraux}\par
par A. Grothendieck
\end{center}

Le présent exposé est de nature introductive, et sa lecture n'est pas logiquement indispensable pour l'étude du Séminaire. Il s'adresse plus particulièrement aux lecteurs au courant de la version provisoire du théorème de Riemann--Roch, telle qu'elle est exposée dans le rapport de Borel--Serre ou dans le rapport de Grothendieck déjà mentionné dans l'introduction, cité [RRR] dans la suite, et qui est reproduit sous forme d'appendice à la fin du présent exposé.

\subsection*{1. Rappelons la formule de Riemann--Roch}
Rappelons la formule de Riemann--Roch pour un morphisme propre
\[
f:X\longrightarrow Y
\]
de schémas quasi-projectifs et lisses sur un corps $k$, et un faisceau cohérent $F$ sur $X$, dont la classe dans le groupe $K(X)$ des classes de faisceaux cohérents sur $X$ est notée $\cl(F)$:
\begin{equation}\tag{1.1}
\Todd(T_Y)\,\ch_Y\bigl(f_!(\cl(F))\bigr)=f_*\bigl(\Todd(T_X)\,\ch_X(F)\bigr).
\end{equation}
Cette formule est valable dans $A(Y)\otimes\mathbb Q$, où $A(Y)$ est l'anneau de Chow de $Y$, le $f_*$ du second membre étant déduit par tensorisation par $\mathbb Q$ de l'homomorphisme «image directe de cycles»
\[
f_*:A(X)\longrightarrow A(Y),
\]
celui du premier membre étant la caractéristique d'Euler--Poincaré de $F$ relativement à $f$:
\[
f_!(\cl(F))=\sum_i(-1)^i\cl\bigl(R^if_*(F)\bigr).
\]
$\ch_X$, $\ch_Y$ désignent les caractères de Chern sur $X$ et sur $Y$, et $T_X$, $T_Y$ les fibrés tangents à $X$ et à $Y$. Comme on sait, $\Todd(-)$ et $\ch(-)$ sont des polynômes universels à coefficients dans $\mathbb Q$ en les classes de Chern de l'argument. Le terme constant de $\Todd(-)$ étant $1$, c'est un élément inversible pour toute valeur de l'argument, de sorte que la formule (1.1) prend, après multiplication par $\Todd(T_Y)^{-1}$, la forme équivalente, plus commode pour notre propos,
\begin{equation}\tag{1.2}
\ch_Y\bigl(f_!(\cl(F))\bigr)=f_*\bigl(\Todd(T_f)\,\ch_X(F)\bigr),
\end{equation}
où on a posé
\begin{equation}\tag{1.3}
T_f=T_X-f^*T_Y\in K(X).
\end{equation}
Ainsi $T_f$ joue le rôle d'un fibré tangent relatif virtuel de $X$ sur $Y$. Dans le cas où le morphisme $f$ est lisse, on a simplement $T_f=T_{X/Y}$, tandis que dans le cas où $f:X\to Y$ est une immersion, on trouve
\[
T_f=-\check N_{X/Y},
\]
où $N_{X/Y}$ désigne le faisceau normal de $X$ dans $Y$.

Un des buts principaux du Séminaire est de généraliser (1.2) simultanément dans deux directions:
\begin{enumerate}[label=\alph*)]
\item se débarrasser de l'hypothèse de l'existence d'un corps de base $k$;
\item remplacer les hypothèses de régularité sur $Y,X$ par une hypothèse de «régularité locale» de $f$.
\end{enumerate}
Enfin, chemin faisant, nous nous occuperons également du problème suivant:
\begin{enumerate}[label=\alph*),start=3]
\item éliminer les hypothèses de quasi-projectivité qui, en l'absence d'un corps de base, s'expriment par l'existence de Modules inversibles amples sur $X$ et sur $Y$.
\end{enumerate}

\subsection*{2. La suppression du corps de base}
Examinons d'abord la généralisation a), en gardant cependant les hypothèses de régularité et d'existence de Modules inversibles amples sur $X$ et sur $Y$.

La définition de $K(X)$, $K(Y)$ et de l'homomorphisme $f_!:K(X)\to K(Y)$ n'offre alors pas de nouveau problème, grâce au fait que $X$ et $Y$ sont réguliers. La voie la plus naturelle pour donner un sens à (1.2) semble donc consister en la définition d'anneaux de Chow $A(X),A(Y)$ et d'un homomorphisme de groupes
\[
f_*:A(X)\longrightarrow A(Y),
\]
en l'établissement d'une théorie des classes de Chern, fournissant des applications
\[
c_i:K(X)\longrightarrow A(X),
\]
et de même pour $Y$, et enfin en la description d'un élément fibré tangent relatif virtuel
\[
T_f\in K(X).
\]

\subsubsection*{2.1}
Pour ce qui est de cette dernière, une définition s'offre de façon assez évidente. On utilise le fait que le morphisme $f$, grâce à l'hypothèse d'existence d'un Module inversible ample sur $X$ et sur $Y$, peut se factoriser en
\[
X\xrightarrow{i}X'\xrightarrow{f'}Y,
\]
où $i$ est une immersion fermée et $f'$ est lisse; on pourra prendre $X'=\mathbb P^r_Y$ pour $r$ convenable. On posera alors
\begin{equation}\tag{2.1}
T_f=i^*T_{X'/Y}-\check N_{X/X'}\in K(X),
\end{equation}
où $\Omega^1_{X'/Y}$ est le Module localement libre des $1$-différentielles relatives de $X'$ sur $Y$, ou Module cotangent relatif de $X'$ sur $Y$, et $N_{X/X'}$ le Module conormal de $X$ dans $X'$, égal par définition à $J/J^2$, où $J$ est l'idéal sur $X'$ définissant $X$. Les hypothèses de régularité faites sur $X$ et $X'$ impliquent d'ailleurs que $N_{X/X'}$ est également localement libre. Enfin, $\check N$ désigne le dual de $N$. On vérifie sans peine, dans l'exposé VIII, que l'élément $T_f$ défini par (2.1) ne dépend pas de la factorisation de $f$ choisie.

\subsubsection*{2.2}
Quant à une définition, sous les conditions envisagées, d'un anneau de Chow et d'une théorie des classes de Chern correspondante, bien qu'il soit plausible qu'une telle théorie doive pouvoir se développer sur le modèle des variétés lisses sur un corps, il semble que ce travail n'ait pas été fait à l'heure actuelle. Nous adopterons donc un point de vue différent, en définissant directement un anneau qui remplace l'anneau de Chow, et dont le produit tensoriel par $\mathbb Q$ est effectivement isomorphe à $A(X)\otimes\mathbb Q$ dans le cas classique.

Une première filtration naturelle de $K(X)$, qui est celle envisagée dans [RRR], est la filtration topologique par la codimension du support des faisceaux cohérents: $\Fil^iK(X)$ est engendré par les $\cl(F)$ avec $F$ cohérent et $\operatorname{codim}(\operatorname{Supp}F,X)\ge i$. Une première question est celle de la compatibilité de cette filtration à la structure multiplicative de $K(X)$. Dans [RRR], cette compatibilité est obtenue, pour $X$ lisse sur un corps, à l'aide du moving lemma de Chow. Moyennant cette assertion, on obtient un anneau gradué $\Gr_{\mathrm{top}}(X)$ et un homomorphisme canonique surjectif
\[
\varphi:A(X)\longrightarrow \Gr_{\mathrm{top}}(X),
\]
transformant la classe d'un cycle premier $Z$ en $\cl(\mathcal O_Z)$, dont il s'avère ensuite, Exp. XIV 4.2, que le noyau est un sous-groupe de torsion. De plus, on montre dans [RRR] comment les images dans $\Gr_{\mathrm{top}}(X)$ des classes de Chern $c_i(F)$ peuvent se calculer directement en termes de $\cl(F)\in K(X)$ et des opérations $\lambda^i$ de puissance extérieure dans $K(X)$.

Ceci montre que, pour $X$ schéma noethérien régulier muni d'un Module inversible ample, l'anneau $\Gr_{\mathrm{top}}(X)$ est un substitut commode pour l'anneau de Chow hypothétique $A(X)$; il a de plus l'avantage sur ce dernier que c'est directement dans cet anneau que se font certains calculs de [RRR] aboutissant à des formules importantes dans cet anneau, formules qui même pour $X$ projectif et lisse sur un corps ne sont pas établies à l'heure actuelle dans $A(X)$ lui-même (Exp. XIV, 4.3 et 6.4). Pour justifier entièrement ce point de vue, il faudrait cependant résoudre le problème signalé plus haut de la compatibilité de la filtration topologique de $K(X)$ avec sa structure d'anneau, problème qui semble essentiellement lié au moving lemma de Chow, qui lui-même est l'ingrédient technique essentiel de la théorie de l'anneau de Chow que nous avons prétendu pouvoir éviter.

\subsubsection*{2.3}
Pour éviter ce problème, on munit $K(X)$ d'une deuxième filtration, plus fine que la filtration topologique, et qui peut se décrire en termes purement algébriques à partir de la structure de $\lambda$-anneau augmenté de $K(X)$. La définition de cette filtration est suggérée par l'expression donnée dans [RRR] des classes de Chern d'un $x\in K(X)$ dans $\Gr_{\mathrm{top}}(X)$:
\begin{equation}\tag{2.3}
c_i(x)=\gamma^i(x-\epsilon(x))\mod \Fil^{i+1}K(X),
\end{equation}
où $\gamma^i$ désigne une certaine combinaison linéaire des $\lambda^j$ ($j\le i$), explicitée dans loc. cit., $\epsilon$ étant l'augmentation. On montre que $\gamma^i(x-\epsilon(x))$ est bien de filtration topologique $\ge i$, ce qui donne un sens à la formule (2.3). On définit alors la $\lambda$-filtration de $K(X)$ par les idéaux $\Fil^i_\lambda K(X)$, où $\Fil^i_\lambda K(X)$ est formé par les sommes d'expressions de la forme
\[
\gamma^{i_1}(x_1-\epsilon(x_1))\cdots \gamma^{i_r}(x_r-\epsilon(x_r)),\qquad i_1+\cdots+i_r\ge i.
\]
L'anneau gradué associé à cette filtration sera noté $\Gr^\bullet(X)$, et c'est lui qui remplacera dans le Séminaire l'anneau de Chow de $X$. On peut d'ailleurs montrer que, dans le cas envisagé ici, l'homomorphisme canonique
\begin{equation}\tag{2.4}
\Gr^\bullet(X)\longrightarrow \Gr_{\mathrm{top}}(X)
\end{equation}
est un isomorphisme modulo torsion, ce qui justifie le point de vue suivant lequel cet anneau tensorisé par $\mathbb Q$ peut jouer en tous points le même rôle que l'anneau de Chow tensorisé par $\mathbb Q$. D'autre part, on a fait exactement ce qu'il fallait pour que la formule (2.3), où on remplace $\Fil_{\mathrm{top}}$ par $\Fil_\lambda$, définisse une théorie des classes de Chern à valeurs dans $\Gr^\bullet(X)$; donc les classes de Chern dans $\Gr_{\mathrm{top}}(X)$ se déduisent à l'aide de l'homomorphisme (2.4). On montre, Exp. V, grâce au fait, Exp. VI, que $K(X)$ est un $\lambda$-anneau spécial au sens de [RRR], c'est-à-dire un $\lambda$-anneau au sens de l'exposé V, que cette notion de classe de Chern satisfait à toutes les propriétés formelles habituelles, passées en revue dans [RRR].

\subsubsection*{2.4}
On notera un avantage très appréciable de la définition de $\Gr^\bullet(X)$ sur celle de $A(X)$ et $\Gr_{\mathrm{top}}(X)$: c'est qu'elle garde un sens pour un schéma $X$ absolument quelconque, et plus généralement pour un topos localement annelé quelconque $X$, en prenant alors pour $K(X)$ le groupe de classes construit avec les Modules localement libres sur $X$, qui est encore un $\lambda$-anneau augmenté, donc muni d'une filtration correspondante, permettant de définir un gradué associé $\Gr^\bullet(X)$. Comme prix de cette généralité, il faut cependant remarquer que lorsqu'on n'est pas disposé à négliger les phénomènes de torsion, c'est-à-dire à tensoriser par $\mathbb Q$, les propriétés de $\Gr^\bullet(X)$ sont moins satisfaisantes à divers égards que celles de ses deux concurrents (Exp. XIV no 4). C'est ainsi que, supposant à nouveau que $f:X\to Y$ est un morphisme propre de schémas noethériens réguliers munis de Modules inversibles amples, il n'est pas vrai en général que l'homomorphisme
\[
f_*:K(X)\longrightarrow K(Y)
\]
applique $\Fil^i_\lambda K(X)$ dans $\Fil^{i-d}_\lambda K(Y)$, où $d$ est la dimension relative de $X$ sur $Y$, et définisse par suite, par passage aux gradués associés, un homomorphisme de degré $-d$ de $\Gr^\bullet(X)$ dans $\Gr^\bullet(Y)$. Cependant, on peut prouver, Exp. VII, qu'il en est ainsi à condition de tensoriser par $\mathbb Q$, de sorte qu'on obtient un homomorphisme de degré $-d$
\begin{equation}\tag{2.5}
f_*:\Gr^\bullet(X)\otimes\mathbb Q\longrightarrow \Gr^\bullet(Y)\otimes\mathbb Q,
\end{equation}
jouant le rôle de l'homomorphisme «image directe de cycles» $A(X)\to A(Y)$.

Ceci posé, nous voyons donc que nous sommes en possession de tous les éléments de structure pour donner un sens à la formule (1.2) lorsqu'on ne postule pas l'existence d'un corps de base, $f:X\to Y$ étant un morphisme propre de schémas noethériens réguliers admettant des Modules inversibles amples.

\subsection*{3. Le cas des schémas généraux}
Montrons maintenant comment on peut encore donner un sens à la formule (1.2) lorsqu'on abandonne l'hypothèse de régularité faite sur $X,Y$, en la remplaçant par une hypothèse de nature locale sur le morphisme $f$, que nous allons préciser.

\subsubsection*{3.1}
Une première difficulté provient du fait que l'observation qui avait servi de point de départ à la théorie [RRR], savoir que le groupe $K(X)$ des classes de faisceaux sur $X$ est le même, qu'on le définisse en termes de faisceaux cohérents ou en termes de faisceaux localement libres, est liée de façon essentielle à l'hypothèse de régularité sur $X$. Dans le cas où $X$ est un schéma, disons noethérien pour simplifier, général, il convient de distinguer entre ces deux groupes, que nous noterons respectivement $K_\bullet(X)$ et $K^\bullet(X)$, le premier ayant un caractère covariant en $X$ pour les morphismes propres, le deuxième un caractère contravariant en $X$ pour les morphismes quelconques. On notera que $K^\bullet(X)$ est muni de façon naturelle d'une structure d'anneau, et même de $\lambda$-anneau augmenté, alors que sur $K_\bullet(X)$ il n'y a plus en général une structure d'anneau, mais seulement une structure de module sur $K^\bullet(X)$. L'homomorphisme $x\mapsto x\cdot\cl(\mathcal O_X)$ de $K^\bullet(X)$ dans $K_\bullet(X)$ n'est plus en général un isomorphisme. Pour notre formulation du théorème de Riemann--Roch généralisé, nous travaillerons exclusivement avec $K^\bullet(X)$, $K^\bullet(Y)$, à l'exclusion de $K_\bullet(X)$, $K_\bullet(Y)$. Comme nous avons signalé dans 2.3, ces $\lambda$-anneaux augmentés, munis de leur $\lambda$-filtration naturelle, donnent naissance à des anneaux gradués $\Gr^\bullet(X)$, $\Gr^\bullet(Y)$, et à des homomorphismes «classes de Chern»
\[
c_i:K^\bullet(X)\longrightarrow \Gr^\bullet(X),
\]
et de même pour $Y$. La difficulté vient cependant dans la définition d'un homomorphisme image directe
\begin{equation}\tag{3.1}
f_*:K^\bullet(X)\longrightarrow K^\bullet(Y),
\end{equation}
qui ne peut plus être définie par la formule naïve
\[
f_*(\cl(F))=\sum_i(-1)^i\cl\bigl(R^if_*(F)\bigr),
\]
car même si $F$ est localement libre, les Modules $R^if_*(F)$ ne sont plus en général localement libres, ni même de tor-dimension finie, et la formule précédente perd son sens. Une solution de cette difficulté est fournie par le point de vue des catégories dérivées. En effet, alors que les $R^if_*(F)$ sont en général des «mauvais» faisceaux, de tor-dimension infinie disons, le complexe $Rf_*(F)$ sur $Y$, image directe de $F$ au sens des catégories dérivées, a tendance à être excellent. De façon précise, lorsque $F$ est de tor-dimension finie sur $Y$, par exemple lorsque $f$ est de tor-dimension finie et $F$ est localement libre sur $X$, alors $Rf_*(F)$ est parfait sur $Y$, c'est-à-dire à cohomologie cohérente, et de tor-dimension globale finie (Exp. III); d'où il résulte facilement, Exp. IV, que sa classe dans le groupe de Grothendieck des complexes parfaits a un sens dans $K^\bullet(Y)$.

\subsubsection*{3.2}
Pour donner un sens à (1.2), il reste encore à faire sur $f$ une hypothèse permettant de donner un sens à l'homomorphisme (2.5), c'est-à-dire assurant que l'homomorphisme (3.1) qu'on vient de définir est, modulo un décalage et modulo torsion, compatible aux filtrations de ces anneaux; et enfin à définir encore un élément «fibré tangent relatif virtuel»
\[
T_f\in K^\bullet(X).
\]
Pour ce dernier, reprenant la définition proposée dans 2.1, on voit qu'il y a lieu de faire une hypothèse assurant que, avec les notations de loc. cit., le Module conormal $N_{X/X'}$ soit localement libre. L'hypothèse la plus naturelle à cet effet est celle que $i$ soit une «immersion régulière», c'est-à-dire identifie $X$ à un sous-schéma de $X'$ défini par un idéal qui localement peut être défini par une suite $\mathcal O_{X'}$-régulière de sections de $\mathcal O_{X'}$. Cette hypothèse, de nature purement locale, ne dépend pas en fait du choix de la factorisation de $f$ en une immersion suivie d'un morphisme lisse (Exp. VIII), et comme plus haut il en est de même pour la classe $T_f$ définie par la formule (2.1). Lorsque l'hypothèse qu'on vient d'envisager est vérifiée, on dit que le morphisme $f$ est un morphisme «localement d'intersection complète». Un tel morphisme est d'ailleurs localement de tor-dimension finie.

On se permettra également de supprimer le vocable «localement» dans cette terminologie, ce qui ne risque d'entraîner aucune confusion.

Si $f$ est un morphisme localement d'intersection compl\`ete, on introduit une notion naturelle de ``dimension relative virtuelle'' $d(f)$ par la formule
\[
d(f)(x)=\dim_x\bigl(f^{-1}(f(x))\bigr)-\codim_x(X,X').
\]
C'est une fonction \`a valeurs enti\`eres localement constante en le point $x\in X$, donc constante si $X$ est connexe, qui ne d\'epend d'ailleurs pas non plus du choix de la factorisation (2.1). Lorsque cette dimension relative virtuelle $d$ est constante, que $f$ est propre et que $X$ et $Y$ admettent des modules inversibles amples, on peut montrer (Exp. VIII), encore que ce soit loin d'\^etre un r\'esultat trivial, que l'on a
\begin{equation}\tag{3.4}
f_*\bigl(\Fil^iK^\bullet(X)_{\mathbb Q}\bigr)\subset \Fil^{i-d}K^\bullet(Y)_{\mathbb Q},
\end{equation}
d'o\`u un homomorphisme (3.3) de degr\'e $-d$.

\subsubsection*{3.3}
Nous avons donc r\'euni tous les \'el\'ements de structure pour donner un sens \`a la formule (1.2) dans le cas d'un morphisme propre et localement d'intersection compl\`ete de sch\'emas noeth\'eriens $X,Y$ admettant tous deux des modules inversibles amples. Moyennant un peu de soin suppl\'ementaire, cette formulation garde d'ailleurs un sens ind\'ependamment de toute hypoth\`ese noeth\'erienne.

Sous cette forme g\'en\'erale, cette formule de Riemann--Roch sera prouv\'ee dans Exp. VIII. Le lecteur constatera que la d\'emonstration donn\'ee dans ce S\'eminaire combine des \'el\'ements des deux d\'emonstrations qui figurent dans [RRR] et dont il a \'et\'e question dans l'Introduction. La n\'ecessit\'e de recourir aux m\'ethodes de la premi\`ere de ces d\'emonstrations, impliquant plus de $K$-formalisme que la deuxi\`eme, suivie dans Borel--Serre, est due \`a la n\'ecessit\'e d'une d\'emonstration de l'inclusion (3.4), question qui ne se posait pas dans le cadre plus restreint de [RRR] et Borel--Serre [2], o\`u on pouvait travailler avec l'anneau de Chow ou l'anneau $\Gr^\bullet_{\mathrm{top}}$. De la deuxi\`eme d\'emonstration, nous conservons l'introduction du sch\'ema \'eclat\'e de $Y$ le long de $X$ dans le cas o\`u $f$ est une immersion. Cet artifice technique dispara\^itra sans doute, et la d\'emonstration du cas d'une immersion s'\'etendra au cas o\`u on ne postule plus l'existence d'un module inversible ample sur $Y$, une fois r\'esolue la question soulev\'ee dans Exp. XIV, no. 1, sur l'introduction d'op\'erations de puissances ext\'erieures dans la cat\'egorie d\'eriv\'ee $D^-(Y)$.

\subsection*{4.}
Il nous faut enfin donner quelques indications sur la fa\c{c}on dont on peut encore donner un sens \`a (1.2) en l'absence d'une hypoth\`ese d'existence de modules inversibles amples sur $X$ et $Y$.

\subsubsection*{4.1}
Notons d'abord que pour un sch\'ema $X$, ou plus g\'en\'eralement un topos localement annel\'e quelconque, il semble hors de doute que le ``bon'' invariant qui doit remplacer l'anneau $K(X)$ de [RRR] n'est pas l'anneau $K^\bullet(X)$ consid\'er\'e pr\'ec\'edemment, d\'efini en termes de modules localement libres sur $X$, mais bien le groupe $K(\Parf(X))$, envisag\'e dans 3.1, d\'efini en termes de complexes parfaits sur $X$, groupe que nous noterons ici $K^\bullet(X)$, et que nous distinguerons du groupe ``na\"if'' de 3.1, que nous notons $K^\bullet_{\mathrm{naif}}(X)$. Ce point de vue est impos\'e par la n\'ecessit\'e de d\'efinir un homomorphisme (3.1) image directe pour un morphisme
\[
f:X\longrightarrow Y
\]
propre et de Tor-dimension finie, de sch\'emas noeth\'eriens disons, le cas non noeth\'erien exigeant des propri\'et\'es de finitude locales un peu plus fortes sur $f$, cf. Exp. III. Avec la d\'efinition de $K^\bullet(X)$ adopt\'ee maintenant, cet homomorphisme $f_*$ se d\'efinit en effet de fa\c{c}on triviale par la formule (3.2), o\`u maintenant $F$ est un complexe parfait quelconque sur $X$, ce qui implique que $Rf_*(F)$ est \'egalement parfait.

\subsubsection*{4.2}
Cette nouvelle d\'efinition de $K^\bullet(X)$ soul\`eve cependant un nouveau probl\`eme d'alg\`ebre homologique, qui est celui d'une d\'efinition d'une $\lambda$-structure sur $K^\bullet(X)$, compatible avec sa structure d'anneau naturelle provenant du produit tensoriel dans la cat\'egorie d\'eriv\'ee. Si $L$ et $L'$ sont deux complexes parfaits, il en est de m\^eme de leur produit tensoriel total $L\otimes^{\mathbf L}L'$. Plus pr\'ecis\'ement, il y aurait lieu de d\'efinir des op\'erations ``puissances ext\'erieures'' dans la cat\'egorie d\'eriv\'ee $D^-(X)$, transformant complexes parfaits en complexes parfaits. Cette question n'a pas encore \'et\'e tir\'ee au clair \`a l'heure actuelle, et comme il a \'et\'e dit dans l'Introduction, elle constitue la lacune la plus s\'erieuse dans le S\'eminaire. Une fois cette question r\'esolue, et proc\'edant comme dans 2.3, on trouve une filtration naturelle sur $K^\bullet(X)$, permettant de former l'anneau gradu\'e associ\'e $\Gr^\bullet(X)$, qui jouera le r\^ole de l'anneau de Chow, et de d\'efinir une th\'eorie des classes de Chern
\[
c_i:K^\bullet(X)\longrightarrow \Gr^i(X).
\]

\subsubsection*{4.3}
Pour donner un sens \`a la formule (1.2) pour un morphisme propre localement d'intersection compl\`ete
\[
f:X\longrightarrow Y
\]
de sch\'emas noeth\'eriens, il faut encore inclure dans l'\'enonc\'e de cette formule les inclusions non triviales (3.4), o\`u $d$ est la dimension relative virtuelle, pour permettre de d\'efinir (3.3), et de d\'efinir enfin la classe tangente relative virtuelle
\[
T_f\in K^\bullet(X).
\]
Lorsque $f$ admet une factorisation (2.1) par une immersion dans un sch\'ema relatif lisse, la d\'efinition de $T_f$ n'offre pas de nouvelle difficult\'e. On peut d'ailleurs la reformuler dans ce cas, d'une fa\c{c}on plus conforme \`a l'esprit ``cat\'egories d\'eriv\'ees'', en consid\'erant l'homomorphisme canonique bien connu
\begin{equation}\tag{4.1}
N_{X/X'}\longrightarrow \Omega^1_{X'/Y}\otimes_{\mathcal O_{X'}}\mathcal O_X
\end{equation}
induit par la diff\'erentielle ext\'erieure, et le complexe de cha\^ines $L_{X/Y}$ d\'efini gr\^ace \`a (4.1) ``en prolongeant par des z\'eros'', complexe \'egal en degr\'e $1$, respectivement $0$, \`a la source, respectivement au but, de l'homomorphisme (4.1). La formule (2.2) prend alors la forme plus sympathique
\begin{equation}\tag{4.2}
T_f=\cl\bigl(L_{X/Y}^{\vee}\bigr),
\end{equation}
formule qui a un sens dans $K^\bullet(X)$ lorsque $f$ est un morphisme localement d'intersection compl\`ete, puisque dans ce cas le complexe $L_{X/Y}$ est manifestement parfait. L'assertion d'invariance de $T_f$ relativement \`a la factorisation choisie de $f$ peut alors se pr\'eciser facilement en une assertion d'invariance pour le complexe $L_{X/Y}$; modulo isomorphisme canonique dans la cat\'egorie d\'eriv\'ee $D(X)$, ce complexe est \'egalement ind\'ependant de la factorisation choisie (Exp. VIII).

\subsubsection*{4.4}
Lorsque $f:X\to Y$ est un morphisme quelconque de sch\'emas, on peut encore construire un objet canonique $L_{X/Y}$ de la cat\'egorie d\'eriv\'ee $D(X)$, d\'efini \`a isomorphisme unique pr\`es, qui sur chaque ouvert affine de $X$ au dessus d'un ouvert affine de $Y$ soit isomorphe au complexe construit suivant le principe expos\'e \`a l'alin\'ea pr\'ec\'edent, o\`u il convient cependant dans le cas g\'en\'eral de remplacer ``lisse'' par ``formellement lisse'', de sorte que toute alg\`ebre $B$ sur un anneau $A$ appara\^it comme quotient d'une alg\`ebre formellement lisse, par exemple une alg\`ebre de polyn\^omes. Pour une telle d\'efinition il n'est pas possible de proc\'eder par simple ``recollement'' \`a partir des constructions affines, car les objets de la cat\'egorie d\'eriv\'ee $D(X)$ sont de nature essentiellement non recollable.

Voici une construction g\'en\'erale de $L_{X/Y}$, valable plus g\'en\'eralement pour tout morphisme
\[
f:X\longrightarrow Y
\]
de topos commutativement annel\'es. Soit $B=\mathcal O_X$, $A=f^{-1}\mathcal O_Y$, de sorte que $A$ est un anneau sur l'espace $X$, et $B$ une $A$-alg\`ebre. Soit $T\to B$ un homomorphisme de faisceaux d'ensembles qui engendre $B$ comme $A$-alg\`ebre, i.e. tel que l'homomorphisme de $A$-alg\`ebres
\[
C=A[T]\longrightarrow B
\]
correspondant soit un \'epimorphisme de faisceaux d'ensembles, o\`u $A[T]$ d\'esigne la $A$-alg\`ebre commutative libre engendr\'ee par $T$, i.e. le faisceau associ\'e au pr\'efaisceau
\[
U\longmapsto A(U)[T(U)].
\]
Soit $J$ le noyau de $C\to B$, et consid\'erons comme dans 4.3 le complexe de cha\^ines de longueur $1$ d\'efini par l'homomorphisme naturel
\[
J/J^2\longrightarrow \Omega^1_{C/A}\otimes_C B.
\]
Il n'est pas difficile de voir que, \`a isomorphisme canonique pr\`es dans la cat\'egorie $D(X)$, le complexe ainsi construit est ind\'ependant du choix de $T$ et de $T\to B$, de sorte qu'il m\'erite la d\'esignation de complexe cotangent relatif $L_{X/Y}$ de $X$ sur $Y$. On v\'erifie \'egalement que pour tout ouvert affine de $X$ au dessus d'un ouvert affine de $Y$, il est donn\'e, \`a isomorphisme canonique pr\`es, par le proc\'ed\'e de 4.3.

Dans le cas o\`u $f$ est localement d'intersection compl\`ete, on voit de plus que l'on obtient m\^eme un complexe parfait, ce qui permet alors de d\'efinir encore un \'el\'ement $T_f$ dans $K^\bullet(X)$ par la formule (4.2). Signalons qu'ici encore, il n'aurait pas \'et\'e possible en g\'en\'eral de d\'efinir $T_f$ comme un \'el\'ement de $K^\bullet_{\mathrm{naif}}(X)$, et rien ne permet de supposer que $T_f$ soit globalement isomorphe dans $D(X)$ \`a un complexe born\'e \`a composantes localement libres de type fini. Cela est donc une justification suppl\'ementaire du point de vue avanc\'e dans 4.1 en faveur du $K^\bullet(X)$ ``sophistiqu\'e'', de pr\'ef\'erence \`a $K^\bullet_{\mathrm{naif}}(X)$.

\subsubsection*{4.5}
Sous r\'eserve de la question soulev\'ee dans 4.2, nous avons donc encore r\'euni tous les \'el\'ements de structure pour donner un sens \`a (1.2) pour un morphisme propre localement d'intersection compl\`ete de sch\'emas quelconques, l'hypoth\`ese noeth\'erienne faite implicitement dans les sections qui pr\'ec\`edent pouvant en fait s'\'eliminer sans difficult\'e essentielle. Cela ne signifie pas malheureusement que, m\^eme sous cette r\'eserve, la formule en question soit \'etablie, m\^eme dans le cas o\`u $Y$ est le spectre d'un corps alg\'ebriquement clos de caract\'eristique $p>0$, et que $X$ est un sch\'ema alg\'ebrique propre et lisse de dimension $3$ disons. Voir commentaires dans Exp. XIV, no. 2.

\subsection*{5.}
Sous la forme pr\'esent\'ee dans le no. 4, la formule de Riemann--Roch (1.2) a \'egalement un sens pour un morphisme propre localement d'intersection compl\`ete d'espaces analytiques complexes, ou d'espaces rigide-analytiques au sens de Tate. Mais bien s\^ur la d\'emonstration donn\'ee dans [RRR] ne s'applique \`a ce cas que lorsqu'on suppose de plus que le morphisme $f$ est projectif, et, comme dans le cas alg\'ebrique, il semble qu'il faille une id\'ee essentiellement nouvelle pour traiter le cas g\'en\'eral. La formule que nous envisageons ici est une formule \`a valeurs dans un ``anneau de Chow analytique'' $\Gr^\bullet(X)$ ou plut\^ot $\Gr^\bullet(X)_{\mathbb Q}$, construit suivant le principe expliqu\'e dans 4.2, la difficult\'e qui y est rencontr\'ee s'\'evanouissant d'ailleurs du fait qu'on est en caract\'eristique nulle, comme on le signale dans Exp. XIV, no. 1.

Lorsque $X$ et $Y$ sont non singuliers, il est possible de donner une autre version de la formule de Riemann--Roch en proc\'edant comme suit. Soit d'abord $X$ un espace analytique complexe quelconque, $X^{\mathrm{top}}$ le m\^eme espace muni du faisceau des fonctions complexes continues. On a donc un morphisme canonique d'espaces annel\'es
\[
X^{\mathrm{top}}\longrightarrow X,
\]
qui d\'efinit, gr\^ace au caract\`ere contravariant \'evident du foncteur $K^\bullet(X)$ en l'espace annel\'e variable $X$, un homomorphisme d'anneaux
\begin{equation}\tag{5.1}
K^\bullet(X)\longrightarrow K^\bullet(X^{\mathrm{top}}).
\end{equation}
Supposons pour simplifier $X^{\mathrm{top}}$ compact; alors, comme dans le cas d\'ej\`a signal\'e dans 3.1 d'un sch\'ema admettant un module inversible ample, on voit (Exp. II) que tout complexe parfait sur $X^{\mathrm{top}}$ est globalement isomorphe \`a un complexe born\'e localement libre en chaque degr\'e, i.e. \`a un complexe de fibr\'es vectoriels complexes sur $X^{\mathrm{top}}$, et que par suite $K^\bullet(X^{\mathrm{top}})$ est isomorphe canoniquement au groupe $K(X^{\mathrm{top}})$ bien connu des topologues [1], d\'efini en termes de fibr\'es vectoriels complexes sur l'espace compact $X^{\mathrm{top}}$. On peut donc interpr\'eter l'homomorphisme (5.1) comme \'etant un homomorphisme du $K^\bullet(X)$ d\'efini en termes de la structure analytique, dans l'anneau bien connu $K(X^{\mathrm{top}})$ des topologues. Si maintenant
\[
f:X\longrightarrow Y
\]
est un morphisme d'espaces analytiques non singuliers compacts, on sait lui associer gr\^ace \`a Atiyah--Hirzebruch [1] un homomorphisme de groupes
\begin{equation}\tag{5.2}
f^{\mathrm{top}}_*:K^\bullet(X^{\mathrm{top}})\longrightarrow K^\bullet(Y^{\mathrm{top}}).
\end{equation}
Utilisant de m\^eme l'homomorphisme $f_*$ de (3.1) d\'efini par la formule (3.2), qui utilise ici implicitement le th\'eor\`eme de finitude de Grauert [4] pour un morphisme propre d'espaces analytiques, on trouve un carr\'e d'homomorphismes naturels, ne faisant intervenir que des groupes du type $K^\bullet$, \`a l'exclusion d'anneaux du type anneaux de Chow ou de cohomologie:
\begin{equation}\tag{5.3}
\begin{array}{ccc}
K^\bullet(X) & \longrightarrow & K^\bullet(X^{\mathrm{top}}) \\
\downarrow f_* && \downarrow f^{\mathrm{top}}_* \\
K^\bullet(Y) & \longrightarrow & K^\bullet(Y^{\mathrm{top}}).
\end{array}
\end{equation}
La formule de type Riemann--Roch \`a laquelle nous faisions allusion consiste \`a affirmer la commutativit\'e de ce carr\'e. On notera la diff\'erence de nature entre cette assertion, qui ne n\'eglige pas les ph\'enom\`enes de torsion et se place dans un groupe $K^\bullet(Y^{\mathrm{top}})$ de nature essentiellement discr\`ete, et la formule de Riemann--Roch du type envisag\'e au no. 4, se pla\c{c}ant dans un groupe du type ``Chow'' qui n'est plus de nature discr\`ete, mais qu'on doit tensoriser en revanche par $\mathbb Q$. Notons que la formule (5.3) est d\'emontr\'ee en tous cas lorsque $X$ et $Y$ sont des vari\'et\'es projectives non singuli\`eres [7], comme la formule du type envisag\'e au no. 4, qui est d\'emontr\'ee dans ce cas dans [RRR]. Il s'agit l\`a d'une variante plausible du th\'eor\`eme de l'index de Atiyah--Singer, et de ses g\'en\'eralisations dues \`a Shih [7] et Atiyah, qui appelle \'evidemment une g\'en\'eralisation commune. Une fois prouv\'ee la commutativit\'e de (5.3), on en conclurait aussit\^ot, \`a l'aide du ``th\'eor\`eme de Riemann--Roch diff\'erentiable'' de Atiyah--Hirzebruch [1], une variante cohomologique de la formule de Riemann--Roch analytique complexe, savoir la commutativit\'e du diagramme
\begin{equation}\tag{5.4}
\begin{array}{ccc}
K^\bullet(X) & \xrightarrow{\;x\mapsto \ch(x)\Todd(T_X)\;} & H^{2\bullet}(X,\mathbb Q) \\
\downarrow f_* && \downarrow f_* \\
K^\bullet(Y) & \xrightarrow{\;y\mapsto \ch(y)\Todd(T_Y)\;} & H^{2\bullet}(Y,\mathbb Q),
\end{array}
\end{equation}
o\`u les fl\`eches horizontales sont d\'eduites des classes de Chern cohomologiques, d\'eduites de l'homomorphisme (5.1) et de l'homomorphisme classe de Chern ordinaire
\[
K^\bullet(X^{\mathrm{top}})\longrightarrow H^{2\bullet}(X,\mathbb Z),
\]
la premi\`ere fl\`eche verticale est la m\^eme que dans (5.3), et la deuxi\`eme est l'homomorphisme de Gysin en cohomologie rationnelle.

\subsection*{6.}
Dans tout ceci, il n'a \'et\'e gu\`ere question que de l'invariant contravariant $K^\bullet(X)$, \`a l'exclusion de l'invariant covariant $K_\bullet(X)$ dont l'existence a \'et\'e signal\'ee dans 3.1. Si on consid\`ere $K^\bullet(X)$ comme donnant un analogue de l'anneau de cohomologie enti\`ere d'un espace topologique, on doit consid\'erer $K_\bullet(X)$ comme donnant un analogue de l'homologie enti\`ere, la structure de module de $K_\bullet(X)$ sur $K^\bullet(X)$ correspondant alors \`a l'op\'eration de cap-produit. Le fait que lorsque $X$ est r\'egulier, l'homomorphisme naturel $K^\bullet(X)\to K_\bullet(X)$ soit un isomorphisme, doit \^etre regard\'e comme l'analogue du th\'eor\`eme de dualit\'e de Poincar\'e sur les vari\'et\'es topologiques orient\'ees, affirmant que le cap-produit avec la classe d'homologie fondamentale d\'efinit un isomorphisme de la cohomologie sur l'homologie. Ces analogies sugg\`erent d'ailleurs l'opportunit\'e de d\'efinir \'egalement, pour un poly\`edre fini $X$ mettons, un invariant $K_\bullet(X)$ de nature covariante, dont les relations \`a l'homologie enti\`ere (suite spectrale, homomorphisme de Chern, etc.) seraient analogues \`a celles existant entre $K^\bullet(X)$ et la cohomologie enti\`ere; et il ne semble pas non plus exclu qu'il puisse exister, en termes de ces groupes $K_\bullet(X)$, des variantes du th\'eor\`eme de Riemann--Roch diff\'erentiable, valable pour des espaces plus g\'en\'eraux que des vari\'et\'es diff\'erentiables compactes.

Revenant au cas o\`u $X$ est un sch\'ema noeth\'erien quelconque, il convient de munir $K_\bullet(X)$ d'une filtration croissante, $\Fil_iK_\bullet(X)$ \'etant engendr\'e par des classes $\cl_\bullet(F)$, avec $F$ faisceau coh\'erent sur $X$ tel que $\dim \operatorname{Supp}F\le i$. On prouvera (Exp. X) la compatibilit\'e de cette filtration avec la structure de module et la filtration de $K^\bullet(X)$, i.e. la formule
\begin{equation}\tag{6.1}
\Fil^iK^\bullet(X)\cdot \Fil_jK_\bullet(X)\subset \Fil_{j-i}K_\bullet(X),
\end{equation}
ce qui permet de munir le gradu\'e associ\'e au module filtr\'e $K_\bullet(X)$, not\'e $\Gr_\bullet(X)$, d'une structure de module sur $\Gr^\bullet(X)$. Ses \'el\'ements peuvent d'ailleurs s'interpr\'eter comme des classes de cycles alg\'ebriques sur $X$, pour une relation d'\'equivalence moins fine que l'\'equivalence rationnelle, lorsque $X$ est de type fini sur un corps de base $k$. Si
\[
f:X\longrightarrow Y
\]
est un morphisme propre de sch\'emas noeth\'eriens, alors
\begin{equation}\tag{6.2}
f_*:K_\bullet(X)\longrightarrow K_\bullet(Y)
\end{equation}
est compatible avec les filtrations, et d\'efinit un homomorphisme de groupes gradu\'es
\begin{equation}\tag{6.3}
f_*:\Gr_\bullet(X)\longrightarrow \Gr_\bullet(Y),
\end{equation}
donnant lieu \`a une ``formule de projection'', i.e. qui est $\Gr^\bullet(Y)$-lin\'eaire, formule qui se r\'eduit par passage aux gradu\'es de la formule de projection analogue pour l'homomorphisme (6.2), qui est $K^\bullet(Y)$-lin\'eaire.

Dans le point de vue propos\'e dans le S\'eminaire, une th\'eorie des intersections maniable, sur les sch\'emas noeth\'eriens $X$ pas n\'ecessairement r\'eguliers, doit consister en l'\'etude combin\'ee des invariants contravariants $K^\bullet(X),\Gr^\bullet(X)$ et des invariants covariants $K_\bullet(X),\Gr_\bullet(X)$, qui jouent des r\^oles de nature essentiellement diff\'erente, et dont les propri\'et\'es se compl\`etent mutuellement. C'est ainsi que les groupes covariants $K_\bullet(X),\Gr_\bullet(X)$ se calculent mieux pour certaines fibrations non propres, gr\^ace \`a la ``suite exacte d'homotopie'' qu'on peut \'etablir pour eux ([RRR] et Exp. IX); par exemple on prouve (Exp. IX) des isomorphismes canoniques tels que
\[
K_\bullet(X[t])\simeq K_\bullet(X),\qquad \Gr_\bullet(X[t])\simeq \Gr_\bullet(X),
\]
qui tombent en d\'efaut en g\'en\'eral pour $K^\bullet$ lorsque $X$ est un sch\'ema non r\'egulier. En revanche, $K^\bullet(X)$ et $\Gr^\bullet(X)$ se pr\^etent bien aux calculs gr\^ace \`a leur structure d'anneaux.

Lorsqu'on consid\`ere des sch\'emas $X$ propres sur un corps, la projection
\[
\pi_X:X\longrightarrow(\operatorname{point})
\]
d\'efinit un homomorphisme ``caract\'eristique d'Euler--Poincar\'e''
\begin{equation}\tag{6.4}
\chi_X=\pi_{X*}:K_\bullet(X)\longrightarrow K_\bullet(\operatorname{point})=\mathbb Z,
\end{equation}
induisant sur le sous-groupe $\Gr_0(X)=\Fil_0K_\bullet(X)$ l'homomorphisme ``degr\'e des $0$-cycles''
\begin{equation}\tag{6.5}
\deg_X:\Gr_0(X)\longrightarrow \mathbb Z.
\end{equation}
L'homomorphisme de multiplication
\[
\Gr^i(X)\times\Gr_i(X)\longrightarrow \Gr_0(X),
\]
compos\'e avec l'homomorphisme degr\'e (6.5), fournit alors un accouplement
\begin{equation}\tag{6.6}
\Gr^i(X)\times \Gr_i(X)\longrightarrow \mathbb Z,
\end{equation}
qui explique qu'\`a certains \'egards $\Gr^\bullet(X)$ et $\Gr_\bullet(X)$ jouent des r\^oles duaux, tout comme la cohomologie et l'homologie. Ainsi, si $f:X\to Y$ est un morphisme de sch\'emas alg\'ebriques propres, les deux homomorphismes
\[
f^*:\Gr^\bullet(Y)\longrightarrow \Gr^\bullet(X),\qquad f_*:\Gr_\bullet(X)\longrightarrow \Gr_\bullet(Y)
\]
sont formellement transpos\'es l'un de l'autre, au sens des accouplements (6.6).

Les homomorphismes (6.4) et (6.5) et l'accouplement (6.6) qu'on en d\'eduit, peuvent \^etre consid\'er\'es comme la source des relations num\'eriques en th\'eorie des intersections sur les sch\'emas alg\'ebriques propres sur un corps donn\'e $k$. Ce point de vue sera d\'evelopp\'e avec quelques d\'etails dans Exp. X, ind\'ependant pour l'essentiel des Exp. VI, VII, VIII, centr\'es sur la d\'emonstration du th\'eor\`eme de Riemann--Roch, et de Exp. IX. On y donnera \'egalement quelques notions sur le comportement de $K^\bullet(X),\Gr^\bullet(X)$ par sp\'ecialisation.

\subsection*{7.}
Comme dans toute th\'eorie d'intersection, il importe de pr\'eciser la place que prend dans la th\'eorie le groupe de Picard, ou groupe des classes de diviseurs dans la terminologie classique. Moyennant une l\'eg\`ere restriction, automatiquement satisfaite lorsque $X$ est par exemple quasi-compact, on \'etablit dans Exp. X un isomorphisme canonique
\begin{equation}\tag{7.1}
c_1:\Pic(X)\xrightarrow{\sim}\Gr^1(X),
\end{equation}
qui montre le r\^ole exact du groupe de Picard dans la th\'eorie que nous proposons. Ce r\^ole justifie une \'etude plus d\'etaill\'ee de ce groupe de Picard dans les expos\'es XII, XIII, qui sont logiquement ind\'ependants du texte du S\'eminaire. Une \'etude de l'\'equivalence num\'erique et de certains th\'eor\`emes de finitude pour le groupe de Picard, annonc\'es sans d\'emonstration dans [6], est faite par Kleiman dans Exp. XIII, en utilisant des techniques de nature purement num\'erique, n'utilisant pas la th\'eorie du foncteur de Picard et de sa repr\'esentabilit\'e, dues \`a lui-m\^eme et \`a Mumford, nettement plus simples que les d\'emonstrations initiales de Grothendieck. Le seul r\'esultat r\'efractaire aux m\'ethodes de Kleiman, et dont il est oblig\'e de faire usage, est [6, C-07 (i)], dont la d\'emonstration et certains raffinements, incluant certains th\'eor\`emes de repr\'esentabilit\'e pour le foncteur de Picard, occupent l'expos\'e XII. Ce dernier rel\`eve d'ailleurs d'un esprit et d'une technique enti\`erement \'etrangers au reste du S\'eminaire, se rattachant \`a [5]; cet expos\'e est logiquement ind\'ependant de tous les autres.

Enfin, dans Exp. XI, de nature nettement plus g\'en\'erale que XII et XIII et plus proche de l'esprit g\'en\'eral du S\'eminaire, on \'etudie, sur un topos localement annel\'e quelconque, un foncteur remarquable appel\'e ``foncteur d\'eterminant'':
\begin{equation}\tag{7.2}
\det^\bullet:\Parf_{\mathrm{is}}(X)\longrightarrow \Inv(X)
\end{equation}
de la cat\'egorie des complexes parfaits sur $X$, les morphismes \'etant les seuls isomorphismes au sens de la cat\'egorie d\'eriv\'ee $D(X)$, dans la cat\'egorie des modules inversibles sur $X$. Ce foncteur est d\'ecrit \`a peu de choses pr\`es par l'exigence d'\^etre de ``nature locale'', et d'\^etre donn\'e, pour un complexe $L$ born\'e et \`a composantes localement libres, par la formule
\begin{equation}\tag{7.3}
\det^\bullet(L)=\bigotimes_i \det(L_i)^{(-1)^i},
\end{equation}
o\`u $\det(L_i)$ d\'esigne la puissance ext\'erieure de degr\'e maximum du module localement libre $L_i$. Bien que cet expos\'e aussi soit logiquement ind\'ependant du reste du S\'eminaire \`a l'exception de l'expos\'e I, il constitue n\'eanmoins un \'el\'ement important du ``yoga'' g\'en\'eral propos\'e dans le S\'eminaire, et d\'evelopp\'e pour l'essentiel dans les expos\'es I \`a XI.

\section*{Bibliographie}
\begin{enumerate}[label={[\arabic*]},leftmargin=2.5em]
\item M. Atiyah et F. Hirzebruch, \emph{Vector bundles and homogeneous spaces}, Symposia in Pure Mathematics, vol. 3, Differential Geometry, p. 7--38, 1961.
\item A. Borel et J.-P. Serre, \emph{Le th\'eor\`eme de Riemann--Roch}, Bull. Soc. math. France, t. 86, p. 97--136 (1958).
\item H. Grauert, \emph{Ein Theorem der analytischen Garbentheorie}, Pub. Math. no. 5, p. 5--64 (1960).
\item A. Grothendieck, \emph{Classes de faisceaux et th\'eor\`eme de Riemann--Roch}, notes mim\'eographi\'ees, Princeton 1957, cit\'e [RRR] et reproduit en Appendice \`a l'Exp. 0 du S\'eminaire.
\item A. Grothendieck, \emph{Technique de descente et th\'eor\`emes d'existence en G\'eom\'etrie Alg\'ebrique}, dans \emph{Fondements de la G\'eom\'etrie Alg\'ebrique}, Extraits du S\'eminaire Bourbaki 1957--1962, Secr\'etariat Math\'ematique, Paris.
\item W. Shih, S\'eminaire de Topologie \`a l'I.H.E.S., 1966/67.
\item M. Atiyah et F. Hirzebruch, \emph{The Riemann--Roch theorem for analytic embeddings}, Topology, vol. 1, p. 151--166 (1962).
\end{enumerate}



\clearpage



\section*{Appendice à l'Exposé 0 : RRR}
\begin{center}
{\Large Classes de faisceaux et théorème de Riemann--Roch}\par
\vspace{0.5em}
par A. Grothendieck\footnote{Ceci est la reproduction textuelle du rapport cité dans l'Introduction du présent Séminaire. Les notes de bas de page indiquées par des signes (*), (**) ont été rajoutées en Octobre 1967.}
\end{center}

\begin{center}
Démonstration du théorème de Riemann--Roch (1er Novembre 1957).
\end{center}

\subsection*{CHAP. I -- $\lambda$-Anneaux (préliminaires formels)}
\begin{longtable}{@{}p{0.74\linewidth}r@{}}
1. Définitions & 1\\
2. Exemples & 4\\
3. Le $\lambda$-anneau défini par un anneau gradué & 8\\
4. Les opérations $\lambda^p(N,x)$ & 13\\
\end{longtable}

\subsection*{CHAP. II -- Classes de faisceaux algébriques cohérents et classes de Chern}
\begin{longtable}{@{}p{0.74\linewidth}r@{}}
1. La théorie de Chow & 18\\
2. Définition des classes de Chern des faisceaux algébriques cohérents & 22\\
3. Généralités fonctorielles sur $K(X)$ & 27\\
4. Quelques résultats techniques & 34\\
5. Définition faisceautique des classes de Chern. Application à l'étude des morphismes d'injection & 43\\
6. Le théorème de Riemann--Roch & 49
\end{longtable}

\section*{CHAPITRE I}
\begin{center}
$\lambda$-Anneaux (Préliminaires formels)
\end{center}

\subsection*{\S 1. Définitions}

On appelle \emph{$\lambda$-anneau}\footnote{Dans ce séminaire, nous disons plutôt \emph{pré-$\lambda$-anneau} (V 2.1), réservant le nom de $\lambda$-anneau à ceux satisfaisant la condition supplémentaire de la page 3 (cf. V 2.4).} un anneau commutatif $K$ avec unité, muni d'une famille d'applications
\begin{equation}\tag{1.1}
\lambda^i:K\longrightarrow K
\end{equation}
($i$ entier $\geq 0$) satisfaisant aux conditions suivantes :
\begin{equation}\tag{1.2}
\lambda_x^0=1,\qquad \lambda_x^1=x,\qquad
\lambda^n(x+y)=\sum_{i=0}^n\lambda_x^i\lambda_y^{n-i}.
\end{equation}
Posons, pour tout $x\in K$,
\begin{equation}\tag{1.3}
\lambda_t(x)=\sum_{n\geq 0}\lambda^n(x)t^n\in K[[t]] ;
\end{equation}
les relations précédentes expriment alors que $x\longmapsto\lambda_t(x)$ est un homomorphisme du groupe additif $K$ dans le groupe multiplicatif $1+K[[t]]^+$ des séries formelles dans $K$, d'augmentation $1$, relevant l'homomorphisme canonique
\[
1+\sum_{i>0}x_it^i\longmapsto x_1.
\]

Ainsi, sur l'anneau $\Z$ existe une unique $\lambda$-structure compatible avec sa structure d'anneau et telle que
\begin{equation}\tag{1.4}
\lambda_t(1)=1+t.
\end{equation}
On a alors
\begin{equation}\tag{1.5}
\lambda_t(n)=(1+t)^n,
\end{equation}
d'où
\[
\lambda^i(n)=\binom{n}{i}.
\]

La notion de $\lambda$-homomorphisme de $\lambda$-anneaux est claire ; un \emph{$\lambda$-anneau augmenté} est un $\lambda$-anneau muni d'un $\lambda$-homomorphisme dans $\Z$.

Pour définir la notion de $\lambda$-anneau \emph{spécial}, les considérations suivantes sont utiles. Soit $k$ un anneau commutatif avec unité, soit
\[
K=1+k[[t]]^+
\]
le groupe des séries formelles à coefficients dans $k$, d'augmentation $1$. On va y introduire une loi de $\lambda$-anneau dont la structure additive soit la multiplication des séries formelles. La multiplication de cet anneau sera notée $f\circ g$. Elle est entièrement déterminée par la condition que les coefficients $c_n=c_n(f\circ g)$ de $f\circ g$ soient donnés par des polynômes universels à coefficients entiers en les coefficients $a_i=c_i(f)$ et $b_i=c_i(g)$ de $f$ et $g$ respectivement :
\begin{equation}\tag{1.6}
c_n=P_n(a_1,\ldots,a_n;b_1,\ldots,b_n),
\end{equation}
qu'elle soit distributive par rapport à l'``addition'' $fg$, et soit donnée pour des facteurs linéaires par :
\begin{equation}\tag{1.7}
(1+at)\circ(1+bt)=1+abt.
\end{equation}
Les $P_n$ se calculent par des calculs de fonctions symétriques ; $P_n$ est isobare de poids $n$ en les $a_i$ et isobare de poids $n$ en les $b_i$ ($a_i$ et $b_i$ étant de poids $i$) et d'ailleurs symétrique en $a=(a_i)$ et $b=(b_i)$.

De même, les $\lambda^i(f)$ sont déterminés sans ambiguïté par la condition que les coefficients
\[
c_n(\lambda^i(f))=c_n\text{ de }\lambda^i(f)
\]
se calculent au moyen de polynômes universels à coefficients entiers en les coefficients $c_n(f)=a_n$ de $f$ :
\begin{equation}\tag{1.8}
c_n=P_{i,n}(a_1,\ldots,a_{in}),
\end{equation}
que les relations (1.2) soient vérifiées, et que, pour $i>1$, on ait :
\begin{equation}\tag{1.9}
\lambda^i(1+at)=1
\end{equation}
($1$ est l'élément nul de l'anneau $K=1+k[[t]]^+$). Les $P_{i,n}$ se calculent encore par des calculs de fonctions symétriques, et l'on constate que $P_{i,n}$ est un polynôme isobare de poids $in$ en les $a_i$. Muni des opérations précédentes, $K$ devient un $\lambda$-anneau ; son élément nul est $1$, son élément unité est $1+t$. Notons aussi la relation :
\begin{equation}\tag{1.7 bis}
(1+at)\circ f=f(at).
\end{equation}

Soit maintenant $K$ un $\lambda$-anneau quelconque. On dit que $K$ est \emph{spécial} si l'homomorphisme additif $\lambda_t$ de $K$ dans $1+K[[t]]^+$ est un $\lambda$-homomorphisme. Cela signifie donc qu'on a les formules :
\begin{equation}\tag{1.10}
\lambda_t(1)=1+t,\quad
\lambda_t(xy)=\lambda_t(x)\circ\lambda_t(y),\quad
\lambda_t(\lambda^i(x))=\lambda^i(\lambda_t(x)),
\end{equation}
ou encore, explicitement :
\begin{equation}\tag{1.11}
\begin{cases}
\lambda^i(1)=0 & \text{si }i>1,\\
\lambda^n(xy)=P_n(\lambda^1x,\ldots,\lambda^n x;\lambda^1y,\ldots,\lambda^n y),\\
\lambda^n(\lambda^i(x))=P_{i,n}(\lambda^1x,\ldots,\lambda^{in}x),
\end{cases}
\end{equation}
où les $P_n$ et les $P_{i,n}$ sont les polynômes universels envisagés plus haut. Il en résulte en particulier les formules (1.5). (On peut montrer que le $\lambda$-anneau $1+k[[t]]^+$ construit avec un anneau commutatif $k$ est spécial ; mais nous n'aurons pas besoin de ce fait.)

\subsection*{\S 2. Exemples}

1) Soient $G$ un groupe et $k$ un anneau commutatif. On considère le groupe $K(G)$, quotient du groupe libre engendré par les classes de représentations de $G$ dans des $k$-modules de type fini par le sous-groupe engendré par les éléments de la forme
\[
(\rho\oplus\rho')-(\rho)-(\rho')
\]
(on note $\oplus$ la somme directe). Le produit tensoriel et les puissances extérieures définissent dans $K(G)$ une structure de $\lambda$-anneau. On peut aussi passer au quotient par les éléments de la forme
\[
(\rho)-(\rho')-(\rho''),
\]
où $\rho$ est une représentation extension de $\rho'$ par $\rho''$, triviale comme extension de $k$-modules. On obtient le groupe $K_r(G)$ des classes réduites de représentations de $G$, et l'on constate que la loi de $\lambda$-anneau passe au quotient.

Lorsque $k$ est un corps algébriquement clos de caractéristique $0$\footnote{Il est inutile que $k$ soit algébriquement clos, cf. VI 3.3 pour un résultat encore plus général (englobant tous ceux qui seront utilisés dans le présent rapport).}, on peut montrer que $K(G)$, et a fortiori $K_r(G)$, est un $\lambda$-anneau spécial ; il n'en est plus de même si la caractéristique est $\neq 0$, mais cependant si $k$ est algébriquement clos, $K_r(G)$ est encore un $\lambda$-anneau spécial. Pour le prouver, on est ramené au cas où $G$ est un produit $\Gl(m,k)\times\Gl(n,k)$ et où l'on ne considère que les représentations rationnelles de $G$ ; notre assertion résulte de la détermination explicite de $K_r(G)$ qui sera faite plus bas (exemple 3) ; enfin, lorsque $k$ est de caractéristique $0$, la complète réductibilité des représentations rationnelles de $G$ montre que $K(G)=K_r(G)$.

Variantes de l'exemple précédent : $G$ est un groupe algébrique défini sur $k$ et l'on ne considère que les représentations rationnelles de $G$ définies sur $k$ ; on peut considérer le cas d'un groupe topologique, ou de Lie, et des représentations continues, etc.\footnote{On peut construire d'autres exemples de $\lambda$-anneaux, par exemple l'ensemble des classes d'équivalence de formes quadratiques, l'ensemble des classes de représentations d'une algèbre de Lie, etc. Nous promettons que nous n'aurons pas besoin de tous ces exemples pour traiter le théorème de Riemann--Roch !}

2) Soit $X$ une variété algébrique, irréductible ou non, définie sur le corps $k$. On pourra considérer le groupe quotient du groupe libre engendré par les classes de fibrés vectoriels sur $X$ (définis sur $k$), par le sous-groupe engendré par les éléments de la forme
\[
(E\oplus E')-(E)-(E'),
\]
sur lequel les opérations de $\lambda$-anneau seront définies par le produit tensoriel et les puissances extérieures. Si $X$ est complète, c'est un groupe libre engendré par les fibrés vectoriels indécomposables sur $X$ (grâce à Krull--Remak--Schmidt) ; la connaissance de cet anneau et de sa base équivaut à la classification complète des fibrés vectoriels sur $X$ et à la connaissance des opérations tensorielles élémentaires sur ces fibrés. Si $k$ est algébriquement clos de caractéristique $0$, toutes les autres opérations tensorielles sur des classes de fibrés sont alors connues, à l'aide de ce qui va suivre dans 3). On peut faire un passage au quotient plus draconien, par le groupe engendré par les éléments
\[
(E)-(E')-(E'')
\]
où $E$ est un fibré vectoriel extension de $E'$ par $E''$ ; on trouve alors un $\lambda$-anneau, noté $K(\xi(X))$, où $\xi(X)$ désigne la catégorie additive des fibrés vectoriels sur $X$, définis sur $k$. En utilisant les résultats de l'exemple 1), on montre que si $k$ est algébriquement clos\footnote{Condition inutile, cf. note de bas de page précédente.}, le $\lambda$-anneau $K(\xi(X))$ est spécial, et qu'il en est de même du premier $\lambda$-anneau défini ici, si la caractéristique de $k$ est $0$, mais non dans le cas général.

3) Détermination de $K(G)$ et de $K_r(G)$ dans certains cas importants. On suppose le corps $k$ algébriquement clos, le groupe $G$ algébrique affine connexe et ``globalement réductif'' (le radical de $G$ est isomorphe à $(k^*)^s$) et l'on ne considère que les représentations linéaires rationnelles. En caractéristique $0$, la complète réductibilité des représentations linéaires rationnelles de $G$ montre que $K(G)=K_r(G)$ est le groupe libre engendré par les classes des représentations rationnelles irréductibles. Or, si $T$ est un tore maximal de $G$, on a
\[
K(T)=K_r(T)=\Z(\That),
\]
algèbre du groupe $\That$ des caractères rationnels de $T$ ; on a $\That\simeq\Z^r$ si $r$ est le rang de $G$ (ceci est indépendant de la caractéristique). Le groupe de Weyl $W$ opère dans $T$, donc dans $K(T)$, d'une manière qui ne dépend que des opérations de $W$ dans le groupe libre $\That$. On a un homomorphisme évident de $\lambda$-anneaux, déduit de l'injection de $T$ dans $G$ :
\begin{equation}\tag{1.13}
K_r(G)\longrightarrow K_r(T)^W=K(T)^W,
\end{equation}
où $K(T)^W$ dénote l'ensemble des points fixes de $W$ dans $K(T)$. La théorie des représentations linéaires de $G$\footnote{Cf. Séminaire Chevalley 1956/58, Groupes de Lie algébriques (École Normale Supérieure).} démontre alors le théorème suivant :

\textbf{Théorème 1.1.} L'homomorphisme (1.13) est un isomorphisme.

On voit donc que $K_r(G)$, une fois connu le ``type'' de $G$ (caractérisé par $W$ opérant sur un groupe $\Z^r$), ne dépend pas de la caractéristique, et que c'est un $\lambda$-anneau spécial. On a une base canonique de $K_r(G)$ correspondant aux orbites de $W$ dans $\That$ ; une autre base s'obtient ainsi : soient $(\rho_i)_{1\leq i\leq r'}$ les classes de représentation de $G$ correspondant aux représentations fondamentales de $G/\rad G$, et $(\sigma_j)_{1\leq j\leq r''}$ une base du groupe des homomorphismes rationnels de $G$ dans $k^*$ ; on a donc $r'+r''=r$ (rang de $G$).

\textbf{Corollaire :} L'anneau $K_r(G)$ s'identifie au sous-anneau
\[
\Z[\rho_1,\ldots,\rho_{r'};\sigma_1,\ldots,\sigma_{r''},\sigma_1^{-1},\ldots,\sigma_{r''}^{-1}]
\]
du corps des fonctions rationnelles en les $\rho_i$ et les $\sigma_j$ à coefficients dans $\Q$.\footnote{Il faut supposer le groupe dérivé $G'$ de $G$ simplement connexe pour que cet énoncé soit correct.}

En particulier, $K_r(G)$ n'a pas de diviseurs de $0$ ; si $G=\prod_i\Gl(n_i,k)$, l'anneau $K_r(G)$ s'explicite ainsi : soit $\rho_i$ la représentation de $G$ déduite de la représentation identique de $\Gl(n_i,k)$ ; si l'on considère alors le corps des fonctions rationnelles sur $\Q$ en des indéterminées $\lambda^j(\rho_i)$ ($1\leq j\leq n_i$), alors $K_r(G)$ s'identifie au sous-anneau engendré par ces indéterminées et les inverses des $\lambda^{n_i}(\rho_i)$. Cela justifie notre assertion que, du moins en caractéristique $0$, les opérations $\lambda^i$ permettent de reconstituer toutes les opérations tensorielles sur les fibrés vectoriels ; ceci est encore vrai sur un corps algébriquement clos de caractéristique quelconque, pourvu qu'on prenne des classes réduites.

\textbf{Remarques :} a) Dans l'exemple 1), l'anneau $K(G)$ est augmenté, l'augmentation étant définie par le degré des représentations, et par passage au quotient, cette augmentation définit une augmentation sur $K_r(G)$. De même, les $\lambda$-anneaux envisagés dans l'exemple 2) sont augmentés lorsque $X$ est $k$-irréductible ; dans $K(\xi(X))$, l'augmentation est définie par le rang d'un fibré vectoriel, c'est-à-dire la dimension de ses fibres.

b) Lorsque le corps $k$ n'est pas algébriquement clos, il est plausible que $K(G)$ est un $\lambda$-anneau spécial en caractéristique $0$, et de même pour $K_r(G)$ en toutes caractéristiques.

c) En caractéristique non nulle, le théorème 1.1 m'a été gracieusement fourni par Chevalley.

\subsection*{\S 3. Le $\lambda$-anneau défini par un anneau gradué}

Soit $A$ un anneau gradué commutatif avec unité ; on supposera pour simplifier que les degrés de $A$ sont positifs et que $A^0=\Z$. On désigne par $\Ahat$ le produit des $A^i$ (c'est un anneau commutatif avec unité), par $\Ahp$ le noyau de l'augmentation évidente $\Ahat\to A^0=\Z$. Alors $1+\Ahp$ est un sous-groupe du groupe multiplicatif des éléments inversibles de $\Ahat$. Considérons le groupe abélien
\begin{equation}\tag{1.14}
\widetilde A=\Z\times(1+\Ahp),
\end{equation}
dont les éléments seront notés $[n,x]$ avec $n\in\Z$ et
\[
x=1+\sum_{i\geq 1}x^i\in1+\Ahp\qquad (x^i\in A^i).
\]
Nous écrirons ce groupe additivement, ainsi :
\begin{equation}\tag{1.15}
[n,x]+[n',x']=[n+n',xx'].
\end{equation}
L'élément nul de ce groupe est $[0,1]$. Nous allons y introduire une structure de $\lambda$-anneau augmenté, compatible avec sa structure additive, l'augmentation étant $\varepsilon:[n,x]\mapsto n$. Le produit $[m,x][n,y]$ est complètement caractérisé par le fait qu'il s'écrive
\[
[mn,x^n y^m(x*y)],
\]
où $x*y$ est un élément de $1+\Ahp$ dont les composantes s'expriment comme polynômes universels à coefficients entiers en les $x^i$ et les $y^i$ :
\begin{equation}\tag{1.16}
(x*y)^i=Q_i(x^1,\ldots,x^i;y^1,\ldots,y^i)\qquad (i\geq1),
\end{equation}
le polynôme $Q_i$ étant isobare de poids $i$. De plus, $x*y$ doit être distributif par rapport à l'``addition'' $xy$ et pour des facteurs ``linéaires'' doit se réduire à ce qui suit :
\begin{equation}\tag{1.17}
[1,1+x^1][1,1+y^1]=[1,1+x^1+y^1].
\end{equation}
Les polynômes $Q_i$ s'évaluent par des calculs de fonctions symétriques élémentaires (il s'agit du calcul des classes de Chern d'un produit tensoriel de fibrés vectoriels à l'aide des classes de Chern $x^i$ et $y^i$ des facteurs et de leurs dimensions respectives $m$ et $n$). De (1.17) on déduit :
\begin{equation}\tag{1.17 bis}
(1+x^1)*(1+y^1)=\frac{1+x^1+y^1}{(1+x^1)(1+y^1)},
\end{equation}
et un produit $x*y$ quelconque se calcule terme à terme en décomposant ``formellement''
\[
x=\prod_{i=1}^n(1+\alpha_i),\qquad y=\prod_{i=1}^n(1+\beta_i),
\]
où les $\alpha_i$ et les $\beta_i$ sont de degré 1.

De cette manière, $\widetilde A$ devient un anneau commutatif augmenté, avec une unité $[1,1]$, et on peut l'interpréter comme l'anneau obtenu par adjonction d'une unité à l'anneau $1+\Ahp$. De plus, sur $\widetilde A$, on définit une filtration en posant $\widetilde A^0=\widetilde A$ et en notant $\widetilde A^n$ l'ensemble des $[0,x]$ avec $x^i=0$ pour $1\leq i<n$ ; cette filtration est compatible avec le produit, car on a
\begin{equation}\tag{1.18}
\left(1+\sum_{i\geq m}x^i\right)*\left(1+\sum_{j\geq n}y^j\right)
=1-\frac{(m+n-1)!}{(m-1)!(n-1)!}x^m y^n+\cdots,
\end{equation}
les termes non écrits étant de degré $>m+n$. Cette formule montre que l'anneau gradué $G(\widetilde A)$ associé à $\widetilde A$ s'identifie à $A$ au point de vue additif, mais non au point de vue multiplicatif ; on peut dire seulement que l'application $\eta'$ de $A$ dans $G(\widetilde A)$ définie par
\begin{equation}\tag{1.19}
\eta'(x^i)=(-1)^{i-1}(i-1)!\,x^i\qquad (x^i\in A^i)
\end{equation}
est un homomorphisme d'anneaux.

On introduit dans $\widetilde A$ des opérations $\lambda^i$ vérifiant les axiomes des $\lambda$-anneaux (\S 1, formules (1.2)), bien déterminées par la condition que les composantes de $\lambda^i[0,x]$ se déduisent des composantes $x^i$ de $x$ par des polynômes universels à coefficients entiers :
\begin{equation}\tag{1.20}
\lambda^i[0,x]=[0,\lambda^i x],\qquad (\lambda^i_x)^{(n)}=Q_{i,n}(x^1,\ldots,x^n)
\end{equation}
($Q_{i,n}$ est isobare de poids $n$), et que l'on ait
\begin{equation}\tag{1.21}
\lambda^i[1,1+x^1]=0\qquad\text{pour }i>1.
\end{equation}
Les polynômes $Q_{i,n}$ s'évaluent par un calcul de fonctions symétriques élémentaires, qui est le calcul des classes de Chern des puissances extérieures d'un fibré vectoriel, à l'aide des classes de Chern et du rang de ce dernier. On constate alors que $\widetilde A$ est un $\lambda$-anneau spécial augmenté et que les $\widetilde A^i$ sont stables par les opérations $\lambda^j$.

Le $\lambda$-anneau augmenté $\widetilde A$ peut être considéré comme un foncteur en $A$ ; par suite l'automorphisme principal
\[
x=\sum_{i\geq0}x^i\longmapsto \check x=\sum_{i\geq0}(-1)^i x^i
\]
de $A$ définit un automorphisme de $\widetilde A$, noté par le même symbole $\check{\ }$.

\textbf{Proposition 1.2.} Pour $x\in\widetilde A$ de la forme
\[
x=[n,1+x^1+\cdots+x^n]
\]
(on a donc $x^i=0$ pour $i>n$), on a $\lambda^i(x)=0$ pour $i>n$, on a $\lambda^n(x)=[1,1+x^1]$ et enfin
\begin{equation}\tag{1.22}
\sum_{i=0}^n(-1)^i\lambda^i(x)=[0,1-(n-1)!x^n+\cdots].
\end{equation}
De cette proposition et du fait que les $\widetilde A^i$ sont stables par les opérations $\lambda^j$, on déduit que la formule (1.22) est valable chaque fois que $\varepsilon(x)=n$, que l'on ait ou non $x^i=0$ pour $i>n$.

La formule (1.22) admet les généralisations suivantes : soit $K$ un $\lambda$-anneau quelconque et $n$ par ailleurs arbitraire ; pour $x\in K$, on pose
\begin{equation}\tag{1.23}
\gamma^n(x)=(-1)^n\sum_{i=0}^n(-1)^i\lambda^i(x+n)=\lambda^n(x+n-1)
\end{equation}
(l'identité des deux derniers membres résulte d'une récurrence immédiate). On a en particulier $\gamma^0(x)=1$ et $\gamma^1(x)=x$ ; on pose ensuite :
\begin{equation}\tag{1.24}
\gamma_t(x)=\sum_{n\geq0}\gamma^n(x)t^n\in K[[t]]
\end{equation}
d'où l'on déduit la formule :
\begin{equation}\tag{1.25}
\gamma_t(x)=\lambda_{t/(1-t)}(x),
\end{equation}
ce qui équivaut à
\begin{equation}\tag{1.25 bis}
\lambda_s(x)=\gamma_{s/(1+s)}(x).
\end{equation}
Ceci prouve que les $\lambda^i$ s'expriment à l'aide des $\gamma^i$ et que
\[
\gamma_t(x+y)=\gamma_t(x)\gamma_t(y),
\]
et par suite que les applications $\gamma^i$ définissent une nouvelle structure de $\lambda$-anneau sur $K$. La formule (1.22) montre alors que pour $K=\widetilde A$ et tout $x\in\widetilde A$, on a
\begin{equation}\tag{1.22 bis}
\gamma^n(x-\varepsilon(x))=[0,1+(-1)^{n-1}(n-1)!x^n+\cdots]
\end{equation}
et par suite, la composante $x^n$ s'obtient, au facteur $(-1)^{n-1}(n-1)!$ près, en réduisant modulo $\widetilde A^{n+1}$ l'élément $\gamma^n(x-\varepsilon(x))$ de $\widetilde A^n$.

\paragraph{L'homomorphisme de Chern :}
Pour tout anneau gradué $A$ vérifiant les conditions précédentes, on va définir un homomorphisme d'anneaux augmentés :
\begin{equation}\tag{1.26}
\ch:\widetilde A\longrightarrow \Ahat\otimes\Q.
\end{equation}
Cet homomorphisme est défini sans ambiguïté par la condition d'être un homomorphisme additif, de transformer $1$ en $1$, d'être tel que les composantes de $\ch(x)$ en degré $>0$ soient données par des polynômes universels isobares à coefficients rationnels en les composantes de $x$, et que
\begin{equation}\tag{1.27}
\ch([1,1+x^1])=\exp x^1=\sum_{n\geq0}(x^1)^n/n!.
\end{equation}
On obtient alors aussitôt l'expression explicite de $\ch$ : soit $\eta$ l'endomorphisme additif de $\Ahat\otimes\Q$ défini pour $x^i$ de degré $i$ par
\begin{equation}\tag{1.28}
\eta(x^i)=1/((-1)^{i-1}(i-1)!)\,x^i.
\end{equation}
On a alors
\begin{equation}\tag{1.29}
\ch\left([n,1+\sum_{i\geq1}x^i]\right)=n+\eta\left(\log\left(1+\sum_{i\geq1}x^i\right)\right).
\end{equation}
Comme on peut ``inverser'' cette formule, du moins si $A$ n'a qu'un nombre fini de composantes homogènes, on voit que dans ce cas, on a un isomorphisme $\ch$ de l'anneau $\widetilde A\otimes\Q$ sur l'anneau $\Ahat\otimes\Q$ (ce qui élucide complètement la structure de $\widetilde A\otimes\Q$).

Rappelons que d'après Hirzebruch [1], une série formelle $f\in\Q[[t]]$ définit par des formules polynômiales universelles un homomorphisme additif
\[
1+\Ahp\longrightarrow 1+(\Ahat\otimes\Q)^+,
\]
noté $\mathfrak C_f$ ; lorsque $f(t)=t/(1-\exp(-t))$, on pose simplement $\mathfrak C_f=\mathfrak C$ ($\mathfrak C$ est l'initiale de Todd). On étend la définition de $\mathfrak C(x)$ à $x\in\widetilde A$. Ceci dit, on a le résultat suivant :

\textbf{Proposition 1.3.} Soit $N\in\widetilde A$ de la forme
\[
N=[q,1+N^1+\cdots+N^q]
\]
(on a $N^i=0$ pour $i>q$). Si l'on pose :
\[
\lambda_{-1}(N)=\sum_{i\geq0}(-1)^i\lambda^i(N)=\sum_{i=0}^{n}(-1)^i\lambda^i(N),
\]
on aura :
\begin{equation}\tag{1.30}
\ch(\lambda_{-1}(N))=(-1)^q N^q\mathfrak C(\check N)^{-1},
\end{equation}
ce qui équivaut à :
\begin{equation}\tag{1.30 bis}
\ch(\lambda_{-1}(\check N))=N^q\mathfrak C(N)^{-1}.
\end{equation}

\subsection*{\S 4. Les opérations $\lambda^p(N,x)$}

Soit $A$ l'anneau des séries formelles, à coefficients dans $\Z$, en deux suites infinies d'indéterminées, notées $\lambda^iN$ et $\lambda^i x$ ($i\geq1$). On pose $\lambda^1N=N$ et $\lambda^1x=x$, et l'on introduit sur $A$ l'unique structure de $\lambda$-anneau spécial pour lequel on ait $\lambda^i(N)=\lambda^iN$ et $\lambda^i(x)=\lambda^ix$ et pour lequel les opérations $\lambda^i$ soient ``continues'' (si l'on se bornait aux polynômes en des indéterminées $\lambda^iN$ et $\lambda^ix$, on aurait le ``$\lambda$-anneau libre engendré par $N$ et $x$'' ; notre anneau $A$ en est un complété). On peut alors former
\begin{equation}\tag{1.31}
\lambda_{-1}(N)=\sum_{i\geq0}(-1)^i\lambda^iN,
\end{equation}
qui est une série formelle de terme constant $1$, donc inversible ; on peut donc définir sans ambiguïté un élément $\lambda^p(N,x)$ de $A$ en posant :
\begin{equation}\tag{1.32}
\lambda^p(N,x)\lambda_{-1}(N)=\lambda^p(x\lambda_{-1}(N)).
\end{equation}

\textbf{Théorème 1.4.} Soit $J_q$ l'idéal fermé de $A$ engendré par les $\lambda^iN$ avec $i>q$ ; identifions $A/J_q$ à l'anneau des séries formelles en les $\lambda^i x$ et les $\lambda^jN$ ($1\leq j\leq q$). Alors $\lambda^p(N,x)$ réduit modulo $J_q$ est un polynôme, isobare de poids $p$ par rapport aux variables $\lambda^i x$.

Soient $N$ et $F$ deux espaces vectoriels de dimensions finies $q,q'$ sur un corps $k$ algébriquement clos de caractéristique $0$ ; soit $N^{(p)}$ le noyau de l'application
\[
(a_1,\ldots,a_p)\longmapsto a_1+\cdots+a_p
\]
de $N^p$ dans $N$. On fait opérer le groupe symétrique $\mathfrak S_p$ sur $N^p$, donc sur $N^{(p)}$, donc sur $\Lambda(N^{(p)})$, puis sur
\[
\Lambda(N^{(p)})\otimes\bigotimes^p F.
\]
Considérons cet espace vectoriel comme un module gradué sur $\mathfrak S_p\times\Gl(N)\times\Gl(F)$ et prenons sa ``partie alternée'' qui est un module gradué sur $\Gl(N)\times\Gl(F)=G$. On a alors
\begin{equation}\tag{1.33}
\lambda^p(N,F)=E_G\left(\left(\Lambda(N^{(p)})\otimes\bigotimes^p F\right)^{\mathrm{alt}}\right).
\end{equation}

N.B. Si $M$ est un complexe de $G$-modules, on note $E_G(M)$ la somme alternée dans $K(G)$ des classes des composantes de $M$ ; c'est la caractéristique d'Euler--Poincaré de $M$ ; dans la formule (1.33), le premier membre $\lambda^p(N,F)$ désigne l'élément de l'anneau $K(G)$ obtenu en substituant aux variables $\lambda^i_N$ et $\lambda^i_F$ les classes dans $K(G)$ des $G$-modules $\lambda^i(N)$ et $\lambda^i(F)$ respectivement.

\textbf{Démonstration :} En développant le deuxième membre de (1.32) par la deuxième formule (1.11), on voit que la première assertion est prouvée si on prouve l'assertion analogue obtenue en faisant $x=1$, i.e. si l'on prouve que dans le $\lambda$-anneau $\Z[\lambda^1N,\ldots,\lambda^qN]$ où $\lambda^iN=0$ pour $i>q$, l'élément $\lambda^p(\lambda_{-1}(N))$ est divisible par $\lambda_{-1}(N)$. A priori, le quotient
\[
\lambda^p(N,1)=\lambda^p(\lambda_{-1}(N))/\lambda_{-1}(N)
\]
est une série formelle contenue dans le corps des quotients $K$ de l'anneau des séries formelles en les $\lambda^iN$ pour $1\leq i\leq q$. D'autre part, désignons aussi par $N$ un vectoriel de dimension $q$ sur un corps algébriquement clos $k$ de caractéristique $0$, et considérons l'élément
\[
\mathscr L^p(N)=E_G((\Lambda(N^{(p)}))^{\mathrm{alt}})\in K(G),
\]
où maintenant l'on a $G=\Gl(N)$. On sait (\S 2, théorème 1.3) que $K(G)$ s'identifie au sous-anneau de $K$ engendré par les $\lambda_N^i$ pour $1\leq i\leq q$ et par l'inverse de $\lambda_N^q$. D'autre part, on vérifie facilement dans cet anneau que l'on a
\[
\mathscr L^p(N)\lambda_{-1}(N)=\lambda^p(\lambda_{-1}(N)).
\]
Comme on est dans un grand corps $K$ et que $\lambda_{-1}(N)\neq0$, on en conclut $\mathscr L^p(N)=\lambda^p(N,1)$ ; les deux membres sont dans l'intersection de $K(G)$ et de l'anneau des séries formelles en les $\lambda_N^i$, donc sont des polynômes en les $\lambda_N^i$.

Pour prouver (1.33) dans le cas général, on vérifie que le deuxième membre satisfait dans $K(G)$ à la formule analogue à (1.32), c'est-à-dire que son produit par $\lambda_{-1}(N)$ est $\lambda^p(F\lambda_{-1}(N))$. De même, le produit de $\lambda^p(N,F)\in K(G)$ par $\lambda_{-1}(N)$ est $\lambda^p(F\lambda_{-1}(N))$, comme il résulte du fait que $K(G)$ est spécial. Comme $K(G)$ est intègre (cf. \S 2) et que $\lambda_{-1}(N)\neq0$ (même référence), on a bien démontré la formule (1.33). cqfd.

\textbf{Remarques :} a) Lorsque la dimension $q'$ de $F$ est $\geq p$, alors (1.33) caractérise le polynôme $\lambda^p(N,x)$ en les $\lambda_N^i$ ($1\leq i\leq q$) et les $\lambda_x^j$ ($1\leq j\leq p$).

b) Il est plausible que (1.33) est encore vérifiée si le corps $k$ n'est pas algébriquement clos\footnote{Cela se prouve en effet par la même méthode.}. Il faudrait résoudre cette question si l'on voulait prouver le théorème de Riemann--Roch sur un corps de caractéristique $0$ non algébriquement clos par la méthode qu'on va exposer. Le passage d'un $\mathfrak S_p$-module à sa partie alternée n'a d'ailleurs de sens raisonnable que si le corps de base est de caractéristique $0$.

Soit maintenant $K$ un $\lambda$-anneau spécial quelconque et soit $N$ un élément de $K$ tel que $\lambda^iN=0$ pour $i$ assez grand. On peut alors, pour tout $x\in K$ et tout entier $p\geq1$, considérer les éléments $\lambda^p(N,x)$, qui s'expriment par des polynômes universels à coefficients entiers en les $\lambda_N^i$ et les $\lambda_x^j$. On a $\lambda^1(N,x)=x$ et l'on a la formule (1.32) dans $K$. Pour exprimer les propriétés des opérations $x\mapsto\lambda^p(N,x)$, il est commode d'introduire un nouvel $\lambda$-anneau noté $K_N$. En tant qu'anneau commutatif, $K_N$ s'obtient par adjonction d'une unité à l'anneau dont le groupe additif sous-jacent est $K$ et dont la multiplication est donnée par
\begin{equation}\tag{1.34}
x_Ny=xy\mu\qquad (\mu=\lambda_{-1}(N)).
\end{equation}
L'application $\lambda_t$ de $K_N$ dans $K_N[[t]]$ est définie par :
\begin{equation}\tag{1.35}
\lambda_t(n,x)=(1+t)^n\sum_{p\geq0}\lambda^p(N,x)t^p,
\end{equation}
où l'on a posé $\lambda^0(N,x)=(1,0)$ (élément unité de $K_N$). Je dis que, muni de ces opérations, $K_N$ est un $\lambda$-anneau spécial. Il s'agit donc de vérifier les formules
\[
\lambda_t(X+Y)=\lambda_t(X)\lambda_t(Y),\quad
\lambda_t(XY)=\lambda_t(X)\lambda_t(Y),\quad
\lambda_t(\lambda^i(X))=\lambda^i(\lambda_t(X))
\]
pour $X,Y\in K_N$. On peut se borner à vérifier ces formules lorsque $X$ et $Y$ sont dans $K$, et l'on constate qu'il suffit de les vérifier lorsque $K$ est le $\lambda$-anneau spécial libre engendré par $X,Y$ et $N$ soumis aux relations $\lambda^iN=0$ pour $i>q$. Mais l'homomorphisme de groupes
\begin{equation}\tag{1.36}
K_N\longrightarrow K,
\end{equation}
qui transforme unité en unité et sur $K$ se réduit à l'homomorphisme $x\mapsto x\mu$, est un homomorphisme compatible avec les structures d'anneaux et les applications $\lambda^i$, comme on le vérifie aussitôt. Dans le cas actuel, sa restriction à $K$ est injective ; comme $K$ est un $\lambda$-anneau spécial, il en résulte facilement que les relations ci-dessus sont bien vérifiées pour $X,Y\in K$, et ceci achève de démontrer notre assertion.

D'après la formule (1.23), il est naturel, pour $x\in K$ et $n\geq1$, de poser
\begin{equation}\tag{1.37}
\gamma^n(N,x)=\sum_{i=0}^{n-1}\binom{n-1}{i}\lambda^i(N,x),
\end{equation}
en d'autres termes, on pose $\gamma^n(N,x)=\gamma^n(x_N)$ où $x_N$ désigne $x$ considéré comme élément du $\lambda$-anneau $K_N$.

\textbf{Proposition 1.5.} Soient $A$ un anneau gradué commutatif avec unité, $N$ un élément de $\widetilde A$ de la forme
\[
N=[q,1+N^1+\cdots+N^q]
\]
et $x$ un élément quelconque de $\widetilde A$ :
\[
x=[\nu,1+\sum_{i\geq1}x^i].
\]
Alors, pour tout entier $j\geq0$, on a $\gamma^{q+j}(N,x)\in\widetilde A^j$, et la composante de degré $j$ de $\gamma^{q+j}(N,x)$ est de la forme
\begin{equation}\tag{1.38}
\left(\gamma^{q+j}(N,x)\right)^{(j)}=(-1)^{j-1}(j-1)!G_{q,j}(\nu,x^1,\ldots,x^j;N^1,\ldots,N^j),
\end{equation}
où $G_{q,j}$ est un polynôme universel à coefficients entiers caractérisé par la condition
\begin{equation}\tag{1.39}
G_{q,j}(\nu,x^1,\ldots,N^j)N^q=(-1)^q(x*\lambda_{-1}(N))^{(q+j)}.
\end{equation}

\textbf{Démonstration.} On peut supposer que $A$ est l'anneau des polynômes à coefficients entiers en des indéterminées $x^i$ et $N^j$ (avec $1\leq j\leq q$). Comme on a évidemment, combinant (1.32) et (1.37),
\begin{equation}\tag{1.40}
\gamma^n(N,x)\lambda_{-1}(N)=\gamma^n(x\lambda_{-1}(N)),
\end{equation}
et que, pour $n=q+j$, le deuxième membre est de filtration $\geq q+j$ (utiliser la formule (1.22 bis) en tenant compte de ce que $x\lambda_{-1}(N)$ est d'augmentation nulle), tandis que $\lambda_{-1}(N)$ est de filtration $q$ (formule (1.22)), il en résulte que $\gamma^n(N,x)$ est de filtration $j$, puisque en vertu de (1.18) le gradué associé à $A$ est ici intègre. Utilisant les formules précédentes, on obtient immédiatement la forme indiquée pour $(\gamma^{q+j}(N,x))^{(j)}$.
\begin{flushright}cqfd.\end{flushright}

\section*{CHAPITRE II}
\begin{center}
Classes de faisceaux algébriques cohérents\par
et classes de Chern
\end{center}

\subsection*{\S 1. La théorie de Chow}

Dans tout ce chapitre, on suppose fixé un corps de base $k$ qu'on supposera algébriquement clos, pour simplifier. Une grande partie de ce qui va suivre, et peut-être tout, est cependant valable sans cette restriction. Les espaces algébriques et les applications régulières considérés seront définis sur $k$.

Un espace algébrique est dit \emph{quasi-projectif} s'il est isomorphe à une partie localement fermée d'un espace projectif.

Soit $X$ une variété non singulière quasi-projective connexe de dimension $n$. On désigne par $A(X)$ et l'on appelle \emph{anneau de Chow} de $X$ l'anneau gradué des classes de cycles sur $X$ à équivalence rationnelle près. ($A^i(X)$ étant formé des classes de cycles de dimension $n-i$.) Le fait que $A(X)$ soit bien un anneau gradué résulte du théorème suivant de Chow :

\textbf{Théorème 2.1.}\footnote{Par suite d'un malentendu de la part de l'auteur, le théorème énonce plus qu'il n'est actuellement connu. En fait, on utilise seulement le ``moving lemma'' de Chow, disant qu'on peut bouger deux cycles dans leurs classes de sorte que leur intersection soit non excédentaire ; cf. Séminaire Chevalley 1958, Anneau de Chow et applications.} Tout cycle sur $X$ est rationnellement équivalent à un cycle non singulier. Étant donnés des éléments de $A^i(X)$ et $A^j(X)$, on peut trouver des représentants non singuliers de ces éléments, soient $Y^{n-i}$ et $Z^{n-j}$, qui se coupent partout transversalement.

Un cycle est dit non singulier si ses composantes sont des sous-variétés non singulières. On dit que deux cycles se coupent partout transversalement si toute composante de l'un coupe partout transversalement toute composante de l'autre, i.e. si les deux variétés sécantes sont non singulières aux points d'intersections et si leurs variétés tangentes y sont en position générale.

Soit $f$ un morphisme (i.e. une application régulière) de $X$ dans $Y$ ($X$ et $Y$ étant quasi-projectives et non singulières), alors $f$ définit (en utilisant le théorème précédent) un homomorphisme d'anneaux gradués :
\begin{equation}\tag{2.1}
f^*:A(Y)\longrightarrow A(X)
\end{equation}
(image réciproque de classes de cycles). Ainsi $A(X)$ devient un foncteur contravariant en $X$.

Soit $f$ un morphisme propre de $X^n$ dans $Y^m$ (i.e. dans la terminologie de Weil, le graphe de $f$ est complet au-dessus de $Y$, cf. [3]). Alors $f$ définit aussi un homomorphisme additif
\begin{equation}\tag{2.2}
f_*:A(X^n)\longrightarrow A(Y^m)
\end{equation}
conservant la dimension des cycles, donc augmentant le degré de $m-n$.




% --- Batch current continuation: source pages 47--75 ---
Ainsi, relativement aux morphismes propres, et faisant abstraction des structures multiplicatives et des graduations, $A(X)$ est un foncteur covariant en $X$. Notons la formule classique
\begin{equation}\tag{2.3}
 f_*(x\cdot f^*(y))=f_*(x)\cdot y .
\end{equation}

Nous utiliserons aussi la théorie de Chow des classes de Chern. Elle consiste à attacher, à tout fibré vectoriel $E$ sur $X$ (toujours supposée quasi-projective non singulière), des classes de Chern
\[
 c^i(E)\in A^i(X)\qquad (i\geq 1),
\]
de façon à satisfaire aux conditions suivantes. On pose $c^0(E)=1$ et
\begin{equation}\tag{2.4}
 \Ctil(E)=\left[\rang E,\sum_{i\geq 0}c^i(E)\right]\in \Atil(X)
\end{equation}
(voir chapitre I, \S 3 pour la définition de $\Atil(X)$). Les conditions à vérifier par une théorie des classes de Chern s'expriment alors ainsi :
\begin{equation}\tag{2.5}
 \Ctil(E)=\Ctil(E')+\Ctil(E'')
\end{equation}
si le fibré vectoriel $E$ est extension de $E'$ par $E''$. Si l'on pose
\[
 C(E)=\sum_i c^i(E),
\]
cette condition s'exprime par $C(E)=C(E')C(E'')$. La formule (2.5) permet d'étendre $\Ctil$ en un homomorphisme additif de $K(\xi(X))$ (cf. chapitre I, \S 2, exemple 2) dans $\Atil(X)$ :
\begin{equation}\tag{2.6}
 \Ctil:K(\xi(X))\longrightarrow \Atil(X),
\end{equation}
et l'on veut que $\Ctil$ soit un $\lambda$-homomorphisme de $\lambda$-anneaux. Cette dernière condition s'écrit aussi
\begin{equation}\tag{2.7}
 \Ctil(E\otimes F)=\Ctil(E)\Ctil(F),\qquad
 \Ctil(\lambda^iE)=\lambda^i\Ctil(E),
\end{equation}
pour deux fibrés vectoriels $E$ et $F$; bien entendu, $\Ctil$ est automatiquement compatible avec les augmentations. Soient $D$ un diviseur sur $X$ et $L(D)$ le fibré vectoriel qu'il définit; on veut
\begin{equation}\tag{2.8}
 c^1(L(D))=\ellc(D),\qquad c^i(L(D))=0\quad (i>1),
\end{equation}
où, pour un cycle quelconque $Z$, on note $\ellc(Z)$ sa classe dans $A(X)$. Enfin, on impose que les $c^i$ soient fonctoriels, i.e. que si $f$ est un morphisme de $X$ dans $Y$ et $E$ un fibré vectoriel sur $Y$, on ait
\begin{equation}\tag{2.9}
 c^i(f^!(E))=f^*(c^i(E)),
\end{equation}
où $f^!(E)$ est l'image inverse de $E$ sur $X$. On peut aussi écrire (2.9) sous la forme suivante : la notion d'image inverse de fibré définit un $\lambda$-homomorphisme de $\lambda$-anneaux augmentés
\begin{equation}\tag{2.10}
 f^!:K(Y)\longrightarrow K(X),\footnote{La notation $f^!$ utilisée ici est en conflit avec celle utilisée en théorie de la dualité (cf. par exemple R. Hartshorne, \emph{Residues and Duality}, Lecture Notes in Math. no. 20, Springer, 1966). C'est pourquoi dans ce Séminaire nous la remplaçons par la notation $f^*$.}
\end{equation}
on écrit ici $K(X)$ au lieu de $K(\xi(X))$, et (2.9) signifie que l'on a commutativité dans le diagramme
\[
\begin{tikzcd}[column sep=large,row sep=large]
K(Y) \arrow[r,"f^!"] \arrow[d,"\Ctil_Y"'] & K(X) \arrow[d,"\Ctil_X"]\\
A(Y) \arrow[r,"f^*"'] & A(X).
\end{tikzcd}
\]
On a mis en indice au symbole $\Ctil$ l'espace sur lequel on le considère.

\textbf{Théorème 2.2.} Il existe une théorie des classes de Chern satisfaisant aux conditions (2.5), (2.7), (2.8) et (2.9).

Nous admettrons aussi la proposition suivante.

\textbf{Proposition 2.3.} Soit $Y$ une partie fermée de $X$ et soit $U$ son complémentaire; si $x$ est un élément de $A(X)$ qui induit $0$ sur $U$, il existe un représentant de $x$ qui soit un cycle porté par $Y$.

\textbf{Corollaire.} Soit $p=\dim X-\dim Y$. Alors, si $i<p$, l'homomorphisme
\[
 A^i(X)\longrightarrow A^i(U)
\]
est bijectif. Tout élément du noyau de $A^p(X)\to A^p(U)$ est combinaison linéaire de composantes de dimension $n-p$ de $Y$; donc est proportionnel à $\ellc(Y)$ si $Y$ est irréductible.

\subsection*{2. Définition des classes de Chern des faisceaux algébriques cohérents}

Soit $\cC$ une catégorie abélienne, et soit $\cC_0$ une sous-classe de $\cC$. On désigne par $K(\cC_0)$ le groupe quotient du groupe libre engendré par les classes, à isomorphisme près, d'éléments de $\cC_0$, par le sous-groupe engendré par les éléments de la forme
\[
(E)-(E')-(E''),
\]
où $E$ est une extension de $E'$ par $E''$, avec $E,E',E''\in\cC_0$. On doit supposer que les classes d'éléments de $\cC_0$, pour la relation d'isomorphisme, forment un ensemble. On désigne par $\gamma_{\cC_0}(E)$ la classe d'un $E\in\cC_0$ dans $K(\cC_0)$, et l'on omet l'indice $\cC_0$ quand il n'y a pas de risque de confusion. Le groupe $K(\cC_0)$ est solution d'un problème universel relatif aux fonctions $E\mapsto f(E)$ définies sur $\cC_0$, à valeurs dans un groupe abélien arbitraire $A$, et qui sont additives en ce sens que $f(E)=f(E')+f(E'')$ chaque fois que $E$ est extension de $E'$ par $E''$.

Soient $F$ un foncteur additif de $\cC$ dans $\cC'$ et $\cC'_0$ une sous-classe de $\cC'$ satisfaisant à la même condition que $\cC_0$, tels que pour tout $A\in\cC_0$ les $R^qF(A)$ soient dans $\cC'_0$ et nuls pour $q$ assez grand; on suppose que dans $\cC$ existent des résolutions injectives, de sorte que les foncteurs dérivés $R^qF$ de $F$ existent. Alors $F$ définit un homomorphisme de $K(\cC_0)$ dans $K(\cC'_0)$, noté encore $F$, par la formule
\begin{equation}\tag{2.11}
 F(\gamma_{\cC_0}(E))=\sum_{i\geq 0}(-1)^i\gamma_{\cC'_0}(R^iF(E)).
\end{equation}
Plus généralement, si l'on a un foncteur cohomologique $(F^i)$ de $\cC$ dans $\cC'$ tel que pour tout $E\in\cC_0$ les $F^i(E)$ soient dans $\cC'_0$ et nuls sauf pour un nombre fini de valeurs de $i$, on en déduit un homomorphisme de $K(\cC_0)$ dans $K(\cC'_0)$. Ces remarques s'étendent aux multifoncteurs; de plus, les suites spectrales classiques donnent des propriétés de transitivité que nous n'expliciterons pas.

En particulier, si $F$ est un foncteur exact, la formule (2.11) se simplifie :
\begin{equation}\tag{2.11 bis}
 F(\gamma_{\cC_0}(E))=\gamma_{\cC'_0}(F(E)).
\end{equation}
Ceci s'applique en particulier si $\cC=\cC'$, si $F$ est le foncteur identique et si $\cC_0\subset\cC'_0$; on a un homomorphisme canonique de $K(\cC_0)$ dans $K(\cC'_0)$.

\textbf{Théorème 2.3.} Soient $\cC$ une catégorie abélienne et $\cC_0\subset\cC_1$ deux sous-classes telles que les classes, à isomorphisme près, d'objets de $\cC_1$ forment un ensemble. On suppose les conditions suivantes satisfaites :
\begin{enumerate}[label=(\roman*)]
\item Pour toute suite exacte $0\to A'\to A\to A''\to0$ dans $\cC$ telle que $A$ et $A''$ soient dans $\cC_0$ (resp. $\cC_1$), on a $A'\in\cC_0$ (resp. $A'\in\cC_1$).
\item Soient $u:A\to B$ et $v:C\to B$, avec $A,B,C\in\cC_1$ et $u$ surjectif; alors il existe $L\in\cC_0$ et des homomorphismes $u':L\to C$ et $v':L\to A$ tels que $u'$ soit surjectif et $vu'=uv'$.
\item Pour tout $A\in\cC_1$, il existe un entier $d=d(A)$ tel que pour toute résolution gauche $L$ de $A$ par des objets de $\cC_0$, on ait $Z_d(L)\in\cC_0$, où $Z_d$ désigne les cycles de dimension $d$.
\end{enumerate}
Dans ces conditions, l'homomorphisme canonique de $K(\cC_0)$ dans $K(\cC_1)$ est un isomorphisme.

\emph{Principe de la démonstration.} Soit $A\in\cC_1$. Il existe, en vertu de (ii) et (iii), une résolution finie $L$ de $A$ par des objets de $\cC_0$. Posons
\[
 \ell(A)=\sum_{i\geq 0}(-1)^i\gamma_{\cC_0}(L_i).
\]
On montre que l'élément de $K(\cC_0)$ ainsi construit ne dépend pas de la résolution $L$ choisie. Si l'on a une deuxième résolution $L'$, on coiffe $L$ et $L'$ par une troisième résolution $L''$, avec des homomorphismes surjectifs de $L''$ dans $L$ et $L'$; ceci est possible grâce à (i) et (ii). On est alors ramené à prouver que $E_{\cC_0}(L)=E_{\cC_0}(L'')$. Or la différence des deux est $E_{\cC_0}(N)$, où $N$ est le noyau de l'homomorphisme surjectif de $L''$ sur $L$; ce noyau est dans $\cC_0$ d'après (i). Comme $N$ est un complexe acyclique en toutes dimensions, $E_{\cC_0}(N)=0$, d'où la conclusion. Ici $E_{\cC_0}(L)$, pour un complexe fini à valeurs dans $\cC_0$, désigne la caractéristique d'Euler--Poincaré
\[
 \sum_i(-1)^i\gamma_{\cC_0}(L_i).
\]
On prouve de façon analogue que la fonction $\ell:\cC_1\to K(\cC_0)$ ainsi construite est additive, donc définit un homomorphisme de $K(\cC_1)$ dans $K(\cC_0)$, et il est immédiat que cet homomorphisme est inverse de l'homomorphisme canonique de $K(\cC_0)$ dans $K(\cC_1)$. CQFD.

\textbf{Remarque.} Pour vérifier (ii), il suffit de vérifier que $\cC_0$ est contenu dans une sous-classe $\cC_2$ de $\cC$ satisfaisant aux deux conditions :
\begin{enumerate}[label=(ii \alph*)]
\item Tout $A\in\cC_1$ est isomorphe à un quotient d'un $L\in\cC_2$.
\item Si $u:A\to B$, avec $A,B\in\cC_2$, alors $\Ker u\in\cC_1$.
\end{enumerate}

\textbf{Corollaire.} Soit $X$ un espace algébrique quasi-projectif. Soient $\cF(X)$ la catégorie des faisceaux algébriques cohérents sur $X$, $\xi(X)$ la catégorie des faisceaux algébriques cohérents localement libres (équivalente à la catégorie des fibrés vectoriels sur $X$), et $\cF_0(X)$ la catégorie des faisceaux algébriques cohérents $F$ tels que pour tout $x\in X$, $F_x$ soit un $\cO_x$-module de dimension cohomologique finie. On a donc $\xi(X)\subset\cF_0(X)\subset\cF(X)$. Alors l'homomorphisme canonique
\begin{equation}\tag{2.12}
 K(\xi(X))\longrightarrow K(\cF_0(X))
\end{equation}
est un isomorphisme.

On vérifie les conditions du théorème 2.3. Les conditions (i) et (iii) résultent de la définition de la dimension cohomologique et du fait qu'un module projectif de type fini sur l'anneau local noethérien $\cO_x$ est libre. Pour vérifier (ii), on démontre que $\cC=\cF(X)$ satisfait aux conditions (ii a) et (ii b) de la remarque. La deuxième est triviale; il reste donc à prouver le résultat suivant.

\textbf{Proposition 2.4 (Serre).} Soit $X$ un espace algébrique quasi-projectif. Alors tout faisceau algébrique cohérent sur $X$ est isomorphe au quotient d'un faisceau algébrique localement libre.

Cette proposition est vraie si $X$ est une partie fermée de l'espace projectif en vertu de [2], puisque, pour $n$ assez grand, $F(n)$ est engendré par ses sections, ce qui signifie que $F$ est un quotient de $\cO(-n)^k$. Dans le cas général, $X$ est une partie ouverte d'une partie fermée d'un espace projectif, et il suffit d'utiliser le résultat suivant, dû à Cartier et Serre dans des cas particuliers.

\textbf{Proposition 2.5.} Soit $F$ un faisceau algébrique cohérent sur une partie ouverte $U$ d'un espace algébrique $X$. Alors $F$ est isomorphe à la restriction d'un faisceau algébrique cohérent sur $X$.

En zornifiant sur les ouverts contenant $U$ sur lesquels on peut prolonger $F$, on est ramené au cas où $X$ est affine. Soit $\overline F$ le faisceau algébrique sur $X$ dont les sections sur un ouvert $V$ sont les sections de $F$ sur $U\cap V$. On constate facilement que $\overline F$ est quasi-cohérent, c'est-à-dire défini sur tout ouvert affine $V$, donc ici sur $X$ tout entier, par un module, de type fini ou non, sur l'anneau de coordonnées de l'ouvert en question. C'est immédiat si $U$ est lui-même affine, car $\overline F$ est alors défini sur l'ouvert $V$ par $\Gamma(U\cap V,F)$ considéré comme module sur $\Gamma(V,\cO_X)$. Dans le cas général, $U$ est réunion finie d'ouverts affines $U_i$, donc $\overline F$ est le noyau d'un homomorphisme évident de faisceaux quasi-cohérents
\[
 \prod_i \overline F|_{U_i}\longrightarrow \prod_{i,j}\overline F|_{U_i\cap U_j},
\]
et est lui-même quasi-cohérent. Comme $X$ est affine, $\overline F$ est donc engendré par ses sections. Comme $\overline F|_U=F$ est cohérent, il existe un nombre fini de sections de $\overline F$ qui engendrent $F|_U$. Ces sections engendrent un sous-faisceau cohérent de $\overline F$, dont la restriction à $U$ est isomorphe à $F$. CQFD.

Supposons maintenant $X$ quasi-projective non singulière. En vertu du théorème des syzygies, on a, avec les notations du corollaire au théorème 2.3,
\[
 \cF_0(X)=\cF(X).
\]
Dans ce cas, on a donc
\begin{equation}\tag{2.13}
 K(\xi(X))=K(\cF(X)).
\end{equation}
On peut exprimer ainsi cet isomorphisme.

\textbf{Corollaire 2 au théorème 2.3.} Soit $X$ une variété quasi-projective non singulière. Pour tout groupe abélien $A$, il y a correspondance biunivoque entre les fonctions additives $f$ de fibrés vectoriels sur $X$, à valeurs dans $A$, et les fonctions additives $g$ de faisceaux algébriques cohérents sur $X$, à valeurs dans $A$. À toute $f$ correspond l'unique $g$ telle que l'on ait
\begin{equation}\tag{2.14}
 f(E)=g(\cO_X(E))
\end{equation}
pour tout fibré vectoriel $E$ sur $X$, en notant $\cO_X(E)$ le faisceau des germes de sections régulières de $E$ sur $X$.

Pour calculer $g(F)$ pour un faisceau algébrique cohérent quelconque, on considère une résolution finie de $F$ par des faisceaux de la forme $\cO_X(E_i)$, et l'on a alors
\begin{equation}\tag{2.14 bis}
 g(F)=\sum_i(-1)^if(E_i).
\end{equation}

En particulier, si $X$ est connexe, on a envisagé au \S 1 une fonction additive $\Ctil(E)$ du fibré vectoriel $E$, à valeurs dans le groupe $A=\Atil(X)$. On en déduit une fonction additive de faisceaux cohérents, soit $\Ctil(F)$, et en passant aux composantes des fonctions $F\mapsto c^i(F)\in A^i(X)$. Les $c^i(F)$ sont appelées les classes de Chern du faisceau cohérent $F$; elles se calculent théoriquement à partir des classes de Chern des fibrés vectoriels par (2.14 bis), mais en écriture multiplicative au second membre. Quant à la composante de $\Ctil(F)$ suivant $\Z$, c'est le rang du faisceau $F$, qui correspond au rang d'un module sur un anneau intègre; si $X$ est affine, cas auquel on peut toujours se ramener, $F$ est défini par un module de type fini sur l'anneau de coordonnées de $X$, et l'on prend le rang de ce module.

\subsection*{3. Généralités fonctorielles sur $K(X)$}

Si $X$ est un espace algébrique, on pose pour abréger $K(X)=K(\cF(X))$; lorsque $X$ est quasi-projectif non singulier, on identifie ce groupe à $K(\xi(X))$ en vertu du corollaire 2 au théorème 2.3. Comme $K(\xi(X))$ est non seulement un groupe abélien mais un $\lambda$-anneau, il s'ensuit par transport de structure que $K(X)$ est aussi un $\lambda$-anneau, donc muni d'une structure multiplicative et d'applications $\lambda^i$. Il importe de définir ces structures directement, sans passer aux fibrés vectoriels.

Si $E$ est un fibré vectoriel sur $X$, on désigne par $\gamma(E)$ sa classe dans $K(X)$; si $F$ est un faisceau algébrique cohérent sur $X$, $\gamma(F)$ désigne de même sa classe dans $K(X)$. Si $Y$ est une sous-variété irréductible de $X$, on pose $\gamma(Y)=\gamma(\cO_Y)$, où $\cO_Y$ désigne le faisceau des anneaux locaux de $Y$, considéré comme faisceau algébrique cohérent sur $X$. On définit par linéarité le symbole $\gamma(Y)$ lorsque $Y$ est un cycle quelconque sur $X$. Quand il y aura des risques de confusion, on écrira $\gamma_X$ pour $\gamma$.

Soit $X$ un espace algébrique quelconque. On introduit sur $K(X)$ une filtration, en désignant par $K(X)^i$ le sous-groupe de $K(X)$ engendré par les $\gamma(F)$ avec
\[
 \dim\Supp F\leq n-i,
\]
où $n$ est la dimension de $X$. Il résulte du lemme de dévissage de Serre, sous la forme donnée dans [3], le résultat suivant.

\textbf{Proposition 2.6.} Le groupe $K(X)^i$ est engendré par les $\gamma(Y)$, où $Y$ parcourt l'ensemble des sous-variétés de dimension $n-i$ de $X$. Si $X$ est irréductible, on a $\Gr^0(K(X))=\Z$.

Ce dernier isomorphisme est donné par l'homomorphisme de $K(X)$ dans $\Z$ défini à l'aide de la notion de rang d'un faisceau algébrique cohérent.

Supposons maintenant $X$ non singulière. S'inspirant des développements du début du \S 2, on définit une multiplication dans $K(X)$ en posant
\begin{equation}\tag{2.15}
 \gamma(F)\gamma(G)=\sum_{i\geq0}(-1)^i\gamma(\Tor_i(F,G)),
\end{equation}
les $\Tor_i$ étant pris, bien entendu, par rapport aux anneaux locaux de $X$. La suite exacte des $\Tor_i$ montre que (2.15) définit bien une application biadditive de $K(X)\times K(X)$ dans $K(X)$; l'associativité résulte de la suite spectrale des $\Tor$ multiples. On vérifie aussi de cette façon que l'on peut calculer le produit de plusieurs facteurs par la formule
\begin{equation}\tag{2.16}
 \gamma(F_1)\cdots\gamma(F_n)=\sum_{i\geq0}(-1)^i\gamma(\Tor_i(F_1,\ldots,F_n)),
\end{equation}
où, dans le deuxième membre, les $\Tor_i$ sont des $\Tor$ simultanés, calculés par résolution projective simultanée des arguments $F_i$.

Si l'on suppose dans (2.15) que $G$ est localement libre, on trouve
\begin{equation}\tag{2.15 bis}
 \gamma(F)\gamma(G)=\gamma(F\otimes G),
\end{equation}
ce qui montre en particulier que $\gamma(\cO_X)$ est élément unité de $K(X)$, et que l'homomorphisme canonique de $K(\xi(X))$ dans $K(X)$ est un homomorphisme d'anneaux.

La formule (2.15) est évidemment suggérée par la formule des $\Tor$ de Serre en théorie des intersections. Cette formule prouve la deuxième partie de la proposition suivante, la première résultant immédiatement de (2.15) et d'un calcul de $\Tor$ locaux.

\textbf{Proposition 2.7.} Soient $X$ une variété non singulière, $Y$ et $Z$ deux cycles dans $X$. Si $Y$ et $Z$ se coupent partout transversalement, on a
\begin{equation}\tag{2.17}
 \gamma(Y)\gamma(Z)=\gamma(Y\cdot Z).
\end{equation}
Si $Y$ et $Z$ sont homogènes et de dimension respective $n-i$ et $n-j$, et s'ils se coupent sans composante excédentaire, on a
\begin{equation}\tag{2.18}
 \gamma_X(Y)\gamma_X(Z)=\gamma_X(Y\cdot Z)\mod K(X)^{i+j+1}.
\end{equation}
On fera attention à ce qu'on n'a pas en général égalité dans (2.18). On verra au \S 4 que, si $X$ est quasi-projective non singulière, la structure d'anneau de $K(X)$ est compatible avec sa filtration et que l'anneau gradué associé $\Gr(K(X))$ s'identifie à un quotient de $A(X)$ par un sous-groupe de torsion, peut-être toujours réduit à zéro.\footnote{Non; pour un contre-exemple cf. ce Séminaire XIV 4.7.}

Je ne sais définir directement les $\lambda^i(F)$ en termes de faisceaux que si le corps $k$ est de caractéristique $0$. C'est la raison qui nous a empêché de démontrer le théorème de Riemann--Roch lorsque $k$ n'est pas de caractéristique $0$. Soit $F$ un faisceau algébrique cohérent sur $X$ sur lequel opère le groupe $\mathfrak S_p$ des permutations de $p$ lettres; on note $F^{\mathrm{alt}}$ l'ensemble des $x\in F$ tels que $s\cdot x=\varepsilon(s)x$ pour tout $s\in\mathfrak S_p$, où $\varepsilon(s)$ est la signature de $s$. C'est évidemment un sous-faisceau algébrique cohérent de $F$. Ceci dit, on a :

\textbf{Théorème 2.8.} Soit $X$ une variété quasi-projective non singulière sur le corps $k$ de caractéristique $0$. Soit $F$ un faisceau algébrique cohérent sur $X$; on a alors
\begin{equation}\tag{2.19}
 \lambda^p(\gamma(F))=\sum_{i\geq0}(-1)^i\gamma\left((\Tor_i(F,F,\ldots,F))^{\mathrm{alt}}\right).
\end{equation}
Il est clair que le groupe $\mathfrak S_p$ agit sur le faisceau $\Tor_i(F,\ldots,F)$.

\emph{Principe de la démonstration.} On considère une résolution finie $L$ de $F$ par des faisceaux localement libres $L_i$, et l'on utilise le lemme général suivant.

\textbf{Lemme 2.9.} Soit $L=(L_i)$ un faisceau algébrique cohérent localement libre gradué, ayant un nombre fini de composantes non nulles. On a dans $K(X)$ la formule
\begin{equation}\tag{2.20}
 E\left((\otimes^pL)^{\mathrm{alt}}\right)=\lambda^p(E(L)).
\end{equation}
Dans cette formule, on considère $\otimes^pL$ comme un faisceau où le groupe $\mathfrak S_p$ opère en tenant compte des graduations, qui introduisent donc des signes de façon bien connue, cf. Cartan--Eilenberg. De plus, si $M$ est un faisceau algébrique cohérent gradué sur $X$, on pose
\[
 E(M)=\sum_i(-1)^i\gamma(M_i).
\]

Pour prouver (2.20), on est ramené à la formule analogue dans $K(G)$, $G$ étant un groupe algébrique, par exemple un produit de groupes $\operatorname{Gl}(n,k)$, et $L$ un $G$-module de dimension finie sur $k$ (cf. chapitre I, \S 2, exemple 2). Dans ce cas, cette formule a déjà été utilisée dans le théorème 1.4, pour montrer que le deuxième membre de (1.33) satisfait la même équation fonctionnelle que le premier membre; la vérification en est élémentaire. CQFD.

Soit $f$ un morphisme de $X$ dans $Y$, $Y$ étant un espace algébrique. La notion d'image inverse de fibré vectoriel définit un homomorphisme $f^!:K(\xi(Y))\to K(\xi(X))$ de $\lambda$-anneaux. En identifiant les fibrés vectoriels à des faisceaux algébriques localement libres, cette notion correspond à l'image réciproque de faisceau algébrique. Dans le cas où $Y$ est non singulière, on définit plus généralement un homomorphisme de groupes
\begin{equation}\tag{2.21}
 f^!:K(Y)\longrightarrow K(X)
\end{equation}
par une formule faisant intervenir une somme alternée de $\Tor$, inspirée par les développements du début du \S 2, et que le lecteur explicitera. On a alors commutativité dans le diagramme
\[
\begin{tikzcd}[column sep=large,row sep=large]
K(\xi(Y)) \arrow[r,"f^!"] \arrow[d] & K(\xi(X)) \arrow[d]\\
K(Y) \arrow[r,"f^!"'] & K(X).
\end{tikzcd}
\]
Ainsi, lorsque $X$ et $Y$ sont toutes deux non singulières et quasi-projectives, les deux définitions de $f^!$ coïncident moyennant l'identification de $K(X)$ à $K(\xi(X))$ et de $K(Y)$ à $K(\xi(Y))$.

Un autre cas où l'on peut définir (2.21), sans que $Y$ soit nécessairement non singulière, est celui où le foncteur $F\mapsto f^!(F)$ de $\cF(Y)$ dans $\cF(X)$ est exact, i.e. lorsque les anneaux locaux de $X$ sont des modules plats sur les anneaux locaux de $Y$. Il suffit alors de poser
\[
 f^!(\gamma_Y(F))=\gamma_X(f^!(F)).
\]
Ce cas se rencontre en particulier lorsque $X$ est un espace fibré localement trivial sur $Y$; pour le voir, on est ramené au cas d'un produit $X=Y\times T$. Notons qu'on vérifie aussi dans ce cas que $f^!(\cO_Z)=\cO_{f^{-1}(Z)}$ si $Z$ est une sous-variété irréductible de $Y$, d'où
\begin{equation}\tag{2.22}
 f^!(\gamma_Y(Z))=\gamma_X(f^{-1}(Z))
\end{equation}
si $X$ est plat sur $Y$. Sans cette dernière condition, (2.22) devient inexacte en général, mais on a l'analogue suivant de la proposition 2.7.

\textbf{Proposition 2.8.} Soit $f:X\to Y$ un morphisme de variétés non singulières, et soit $Z$ un cycle sur $Y$. Si $Z$ est partout transversal à $f$, la formule (2.22) est exacte; si $Z$ est homogène de codimension $i$ et si $f^{-1}(Z)$ n'a pas de composante excédentaire, on a
\begin{equation}\tag{2.23}
 f^!(\gamma_Y(Z))=\gamma_X(f^{-1}(Z))\mod K(X)^{i+1}.
\end{equation}
Démonstration analogue à celle de la proposition 2.7.

Signalons enfin que, lorsque $X$ et $Y$ sont quasi-projectives et non singulières, $f^!$ est compatible avec les filtrations de $K(X)$ et de $K(Y)$.

Soient de nouveau $X$ et $Y$ deux espaces algébriques quelconques, et $f$ un morphisme propre de $X$ dans $Y$. Pour tout faisceau algébrique cohérent $F$ sur $X$, le faisceau image directe $f_*(F)$ est algébrique cohérent, et il en est de même des faisceaux $R^if_*(F)$. Par définition, $R^if_*(F)$ est le faisceau sur $Y$ associé au préfaisceau qui à l'ouvert $U$ fait correspondre $H^i(f^{-1}(U),F)$; on a en particulier
\begin{equation}\tag{2.24}
 R^0f_*(F)=f_*(F),\qquad R^qf_*(F)=0\quad(q>\dim X),
\end{equation}
et
\begin{equation}\tag{2.25}
 \Gamma(U,R^qf_*(F))=H^q(f^{-1}(U),F)
\end{equation}
si l'ouvert $U$ est affine. D'après le \S 2, on définit un homomorphisme de $K(X)$ dans $K(Y)$, noté encore $f_*$, par la formule
\begin{equation}\tag{2.26}
 f_*(\gamma(F))=\sum_{q=0}^{\infty}(-1)^q\gamma(R^qf_*(F)).
\end{equation}
Ainsi $K(X)$ devient un foncteur covariant en $X$, relativement aux morphismes propres. La relation $(fg)_*=f_*g_*$, tout comme la relation $(fg)^!=g^!f^!$ oubliée en son lieu, est un cas particulier des relations de transitivité auxquelles nous avions fait allusion au \S 2.

Notons la formule suivante, valable pour $Y$ non singulière :
\begin{equation}\tag{2.27}
 f_*(x\cdot f^!(y))=f_*(x)\cdot y\qquad(x\in K(X),\ y\in K(Y)).
\end{equation}
Dans le cas particulier où les images réciproques des points de $Y$ par $f$ sont finies et contenues dans un ouvert affine,\footnote{Cette dernière condition est évidemment conséquence de la première, chose que l'auteur ignorait encore en 1957 !} le foncteur $F\mapsto f_*(F)$ est exact et $R^qf_*(F)=0$ pour $q>0$; on a donc $f_*(\gamma_X(F))=\gamma_Y(f_*(F))$. En particulier, si $X$ est un sous-espace algébrique de $Y$ et si $i$ est l'injection canonique de $X$ dans $Y$, on a
\begin{equation}\tag{2.28}
 i_*(\gamma_X(F))=\gamma_Y(F),
\end{equation}
où, au second membre, $F$ est considéré comme un faisceau sur $Y$ nul en dehors de $X$.

\subsection*{4. Quelques résultats techniques}

Soit $X$ un espace algébrique. La projection $p$ de $X\times k$ sur $X$ définit un homomorphisme
\begin{equation}\tag{2.29}
 p^!:K(X)\longrightarrow K(X\times k).
\end{equation}

\textbf{Proposition 2.9.} L'homomorphisme (2.29) est surjectif.

On raisonne par récurrence sur $n=\dim X$. L'assertion est triviale si $n<0$; supposons $n\geq0$. En vertu de la proposition 2.6, il suffit de vérifier que, pour toute partie fermée irréductible $Y$ de $X\times k$, $\gamma(Y)$ est dans l'image de $p^!$. Il suffit de continuer le raisonnement en supposant $X$ irréductible. Si $p(Y)$ n'est pas dense dans $X$, l'assertion résulte de l'hypothèse de récurrence; sinon, ou bien $Y=X\times k$ et l'assertion est triviale, ou bien $Y$ est une hypersurface dans $X\times k$. Soit alors $U$ un ouvert affine de $X$ formé de points non singuliers. Si $D$ est le diviseur sur $U\times k$ induit par $Y$, on sait que $D$ est linéairement équivalent à un diviseur de la forme $D'\times k$, où $D'$ est un diviseur sur $U$. Alors $\gamma_{U\times k}(D)-\gamma_{U\times k}(D'\times k)$ appartient à $K(U\times k)^2$, et d'après ce qui a déjà été démontré cet élément est de la forme $p^!(x')$ avec $x'\in K(U)$. Donc $\gamma_{U\times k}(D)=p^!(x'+\gamma_U(D'))$. Mais $x'$ est restriction d'un élément $y$ de $K(X)$, et la restriction de $\gamma_{X\times k}(D)-p^!(y)$ à $U\times k$ est nulle. D'après la proposition 2.10, cet élément provient de $K((X-U)\times k)$, et l'on achève par l'hypothèse de récurrence. CQFD.

\textbf{Corollaire.} Soient $X$ un espace algébrique et $p$ la projection de $X\times k^n$ sur $X$; alors $p^!:K(X)\to K(X\times k^n)$ est surjectif. Il est même bijectif si $X$ est non singulière.

Récurrence sur $n$ pour la première assertion. Pour la deuxième, on pose $f(x)=(x,0)$, d'où $pf=1$ et $f^!p^!=1$.

\textbf{Proposition 2.10.} Soient $X$ un espace algébrique, $U$ une partie ouverte de $X$ et $Y$ son complémentaire. Alors la suite
\begin{equation}\tag{2.30}
 K(Y)\longrightarrow K(X)\longrightarrow K(U)\longrightarrow 0
\end{equation}
est exacte.

N.B. On comparera avec la proposition 2.3.

La surjectivité de $K(X)\to K(U)$ résulte aussitôt de la proposition 2.5 ou de la proposition 2.6. Reste à prouver que si $x\in K(X)$ a une restriction nulle à $K(U)$, alors $x$ est dans l'image de $K(Y)$. On revient à la définition : $x=\gamma_X(F)-\gamma_X(G)$ avec $F,G$ cohérents, et $\gamma_U(F)-\gamma_U(G)=0$. Cette dernière relation signifie qu'il existe des faisceaux cohérents $P'$ et $Q'$ sur $U$, munis de filtrations dont les facteurs soient deux à deux isomorphes, et un isomorphisme sur $U$
\[
 P=F+Q'\longrightarrow Q=G+P'.
\]
Par la proposition 2.5, $P'$ et $Q'$ proviennent de faisceaux cohérents sur $X$, et leurs filtrations se prolongent par le lemme 2.11. On obtient
\[
 \gamma_X(F)-\gamma_X(G)=(\gamma_X(P)-\gamma_X(Q))+\bigl(\gamma_X(P')-\gamma_X(Q')\bigr).
\]
En introduisant les facteurs $P_i$ et $Q_i$, on se ramène au cas où $R|_U\simeq S|_U$. Le graphe $T$ de cet isomorphisme est restriction à $U$ d'un sous-faisceau cohérent de $R\times S$; alors
\[
 \gamma_X(R)-\gamma_X(S)=(\gamma_X(R)-\gamma_X(T))-(\gamma_X(S)-\gamma_X(T)),
\]
et chaque terme est dans l'image de $K(Y)$, car les noyaux et conoyaux des projections correspondantes sont portés par $Y$.

\textbf{Lemme 2.11.} Soient $X$ un espace algébrique, $F$ un faisceau algébrique cohérent sur $X$, $U$ une partie ouverte de $X$ et $G$ un sous-faisceau cohérent de $F|_U$. Alors $G$ est la restriction à $U$ d'un sous-faisceau cohérent de $F$.

Soit $\overline G$ le sous-faisceau de $F$ dont les sections sur un ouvert $V$ sont les sections de $F$ sur $V$ dont la restriction à $U\cap V$ est une section de $G$. Il suffit de prouver que $\overline G$ est cohérent. On peut supposer $X$ et $U$ affines. L'anneau de coordonnées de $U$ est un anneau de fractions $A_S$ de l'anneau de coordonnées $A$ de $X$; si $F$ est défini par le $A$-module $M$, alors $F|_U$ est défini par $M_S=A_S\otimes_AM$. Le sous-faisceau $G$ de $F|_U$ est défini par un sous-module $N'$ de $M_S$, et $\overline G$ est défini par le sous-module $N$ de $M$, image inverse de $N'$ par l'homomorphisme canonique $M\to M_S$. CQFD.

\textbf{Corollaire 1 à la proposition 2.10.} Soit $x\in K(X)$. Pour que $x$ appartienne à $K(X)^i$, il faut et il suffit qu'il existe une partie fermée $Y$ de $X$, de dimension $\leq n-i$, telle que la restriction de $x$ à $X-Y$ soit nulle.

\textbf{Corollaire 2.} Supposons $X$ connexe de dimension $n$, sans singularités en dimension $n-1$. Si $D$ et $D'$ sont deux diviseurs linéairement équivalents sur $X$, on a
\[
 \gamma(D)=\gamma(D')\mod K(X)^2.
\]
Il suffit, par le corollaire 1, de supposer $X$ non singulière; on utilise alors
\begin{equation}\tag{2.31}
 \gamma(D)=1-\gamma(L(-D))\mod K(X)^2,
\end{equation}
ou encore
\begin{equation}\tag{2.31 bis}
 \gamma(L(-D))=1-\gamma(D)\mod K(X)^2.
\end{equation}
Les deux membres sont des homomorphismes du groupe des diviseurs dans le groupe multiplicatif des éléments inversibles de $K(X)/K(X)^2$, et l'on est ramené au cas d'une hypersurface irréductible. Plus généralement, si $D$ est une hypersurface dont les composantes sont disjointes, on a l'égalité
\begin{equation}\tag{2.32}
 \gamma(D)=1-\gamma(L(-D)),
\end{equation}
comme il résulte de la suite exacte
\[
0\to\cO_X(L(-D))\to\cO_X\to\cO_D\to0.
\]

\textbf{Théorème 2.12.} Soient $X$ une variété quasi-projective non singulière, $Z$ et $Z'$ deux cycles homogènes de codimension $p$ sur $X$ linéairement équivalents. Alors
\begin{equation}\tag{2.33}
 \gamma_X(Z)=\gamma_X(Z')\mod K(X)^{p+1}.
\end{equation}
On peut trouver un cycle $Y$ de dimension $n-p+1$ sur $X\times k$ et deux points $a,a'\in k$ tels que $Z=Y\cdot(X\times a)$ et $Z'=Y\cdot(X\times a')$. Par (2.23), $\gamma_X(Z)=i_a^!(\gamma_{X\times k}(Y))$ modulo $K(X)^{p+1}$, et de même pour $a'$. Il suffit donc de prouver $i_a^!=i_{a'}^!$, ce qui résulte aussitôt de la proposition 2.9. CQFD.

\textbf{Corollaire 1.} La filtration de $K(X)$ est compatible avec sa structure d'anneau : $K(X)^iK(X)^j\subset K(X)^{i+j}$. On a de plus un homomorphisme surjectif naturel d'anneaux gradués
\begin{equation}\tag{2.34}
 \varphi:A(X)\longrightarrow\Gr(K(X))
\end{equation}
défini par $\varphi(\ellc(Z^{n-p}))=\gamma(Z^{n-p})\mod K(X)^{p+1}$. Le corollaire résulte du théorème 2.12, des propositions 2.6 et 2.7 et du théorème de Chow 2.1.

\textbf{Corollaire 2.} Soient $X$ et $Y$ deux variétés quasi-projectives non singulières et $f:X\to Y$ un morphisme. Alors $f^!:K(Y)\to K(X)$ est compatible avec les filtrations, et l'on a un diagramme commutatif
\[
\begin{tikzcd}[column sep=large,row sep=large]
A(Y) \arrow[r,"\varphi_Y"] \arrow[d,"f^*"'] & \Gr(K(Y)) \arrow[d,"\Gr(f^!)"]\\
A(X) \arrow[r,"\varphi_X"'] & \Gr(K(X)).
\end{tikzcd}
\]
On décompose $f$ au moyen de son graphe en une immersion suivie d'une projection. Pour une projection le corollaire est trivial; pour une immersion, on procède comme pour le corollaire 1, en utilisant la proposition 2.8 au lieu de 2.7.

\textbf{Proposition 2.13 (``décomposition cellulaire'').} Soit $P$ un espace algébrique, irréductible ou non, muni d'une famille croissante $P^0\subset P^1\subset\cdots\subset P^n=P$ de parties fermées, telles que pour tout $i$, $0\leq i\leq n$, les composantes connexes de $P^i-P^{i-1}$ soient des espaces algébriques $V_{i,\alpha}$ isomorphes à $k^i$; on pose $P^{-1}=\varnothing$. Alors :
\begin{enumerate}[label=\alph*)]
\item Pour tout espace algébrique $X$, l'application naturelle de $K(X)\otimes K(P)$ dans $K(X\times P)$ déduite du produit tensoriel de faisceaux sur $X$ et $P$ est surjective.
\item Les $\gamma_P(V_{i,\alpha})$ forment un système de générateurs de $K(P)$.
\end{enumerate}
La démonstration se fait par récurrence sur $n$, en utilisant le corollaire de la proposition 2.9 et la proposition 2.10. N.B. J'ignore si les $\gamma_P(V_{i,\alpha})$ forment une base de $K(P)$ et si l'homomorphisme de a) est bijectif; nous n'en aurons pas besoin.

\textbf{Théorème 2.14.} Soient $x_i$ ($1\leq i\leq q$) et $E_j$ ($1\leq j\leq r$) des lettres, et $n_j>0$. On considère les symboles $\lambda^k(x_i)$ ($k\geq1$) et $\lambda^k(E_j)$ ($1\leq k\leq n_j$) comme des indéterminées, et soit $R$ un polynôme à coefficients entiers. On suppose que, pour tout anneau gradué $A$, lorsque l'on substitue dans $R$ aux $x_i$ des éléments $\xi_i$ et aux $E_j$ des éléments $\eta_j=[n_j,1+\eta_j^{(1)}+\cdots+\eta_j^{(n_j)}]$ de $\Atil$ soumis aux restrictions $\varepsilon(\xi_i)=0$ et $\eta_j^{(k)}=0$ pour $k>n_j$, on ait
\[
 R((\lambda^k(\xi_i)),(\lambda^k(\eta_j)))\in\Atil^p.
\]
Alors, pour toute variété quasi-projective irréductible non singulière $X$, tous éléments $x'_i\in K(X)^1$ et tous fibrés vectoriels $E'_j$ de dimension $n_j$ sur $X$, on a
\begin{equation}\tag{*}
 R((\lambda^k(x'_i)),(\lambda^k(\gamma(E'_j))))\in K(X)^p.
\end{equation}

\emph{Démonstration.} En vertu du \S 2, les $x'_i$ sont des différences d'éléments de la forme $\gamma(E)$. En plongeant localement $X$ dans un espace projectif, et en utilisant que $F(n)$ est engendré par ses sections pour $n$ assez grand, on voit que les fibrés vectoriels intervenant sont des images réciproques de fibrés canoniques sur des grassmanniennes. Grâce au corollaire 2 du théorème 2.12, on se ramène au cas où $X$ est un produit de grassmanniennes. Dans ce cas, la conclusion résulte du lemme suivant.

\textbf{Lemme 2.15.} Supposons que $X$ soit un produit de grassmanniennes. Alors on peut trouver un $\lambda$-homomorphisme de $K(X)$ dans un $\lambda$-anneau sans torsion de la forme $\Atil$, où $A$ est un anneau gradué, compatible avec les filtrations, et tel que l'homomorphisme correspondant $\Gr(K(X))\to\Gr(A)$ soit injectif.

Pour prouver le lemme, nous utiliserons la théorie de Chow des classes de Chern. Nous considérons l'homomorphisme $\Ctil:K(X)\to\Atil(X)$ et posons $A=A(X)$; il faut prouver que $\psi=\Gr(\Ctil):\Gr(K(X))\to\Gr(A(X))\simeq A(X)$ est injectif. Le fait que $\Ctil$ soit compatible avec les filtrations est vrai pour toute variété quasi-projective non singulière. D'autre part, si $X$ est un produit de grassmanniennes $G_j$, on a sur $X$ des fibrés vectoriels $E_j$, images réciproques des fibrés canoniques. Nous utilisons les faits connus : a) sur $X$, l'équivalence linéaire des cycles équivaut à l'équivalence numérique et $A(X)$ est de rang fini; b) les $c^i(E_j)$ engendrent $A(X)$. De a) résulte que $\varphi:A(X)\to\Gr(K(X))$ est injective, donc bijective, par le lemme suivant.

\textbf{Lemme 2.15.} Soit $X^n$ une variété projective non singulière, et soit $Z^{n-p}$ un cycle sur $X$ tel que $\varphi(Z)\in\Gr^p(K(X))$ soit nul. Alors $Z^{n-p}$ est numériquement équivalent à zéro.

Pour le démontrer, on utilise la fonction additive $\chi(F)=\sum_{i\geq0}(-1)^i\dim H^i(X,F)$ et le fait que cette fonction vaut $1$ pour $F=\cO_{(x)}$, pour tout $x\in X$.

Ainsi $A(X)$ et $\Gr(K(X))$ ont, dans le cas envisagé, même rang fini en tout degré, et $\Gr(K(X))$ est sans torsion. De b) résulte formellement que $\Ctil:K(X)\to\Atil(X)$ est surjective modulo des groupes finis. On en déduit que $\Gr(\Ctil)$ est bijectif modulo des groupes finis en tout degré; comme $\Gr(K(X))$ est sans torsion, $\Gr(\Ctil)$ est injectif. CQFD.

N.B. L'homomorphisme $\Ctil$ n'est pas surjectif en général, même si $X$ est un espace projectif de dimension $2$.

\subsection*{5. Définition faisceautique des classes de Chern. Application à l'étude des morphismes d'injection}

\textbf{Théorème 2.16.} Soit $X$ une variété quasi-projective non singulière. Alors, pour tout $x\in K(X)$ et tout entier $i\geq0$, l'élément $\gamma^i(x-\varepsilon(x))$ est de filtration $\geq i$. Soit $c^i(x)\in\Gr^i(K(X))$ la classe de cet élément modulo $K(X)^{i+1}$; si l'on pose
\[
 \widetilde c(x)=\left[\varepsilon(x),1+\sum_{i\geq0}c^i(x)\right]\in\widetilde{\Gr(K(X))},
\]
alors $\widetilde c$ satisfait aux conditions analogues à celles du théorème 2.2.

La première assertion résulte du théorème 2.14, ainsi que le fait que $x\mapsto\widetilde c(x)$ soit un homomorphisme de $\lambda$-anneaux. La fonctorialité de $\widetilde c$ se vérifie immédiatement, et (2.8) résulte de (2.31). CQFD.

\textbf{Corollaire 1.} On a $c^i(x)=\varphi(c^i(x))$.

Ceci résulte d'un théorème d'unicité pour les classes de Chern, qui se démontre comme chez Hirzebruch par passage aux variétés de drapeaux. On doit seulement montrer que si $Y$ est une variété de drapeaux associée à un fibré vectoriel sur $X$, alors $\Gr(K(X))\to\Gr(K(Y))$ est injectif; cela se vérifie par la détermination explicite de $K(Y)$ à partir de $K(X)$.

\textbf{Corollaire 2.} Supposons que, pour tout $Z\in A(X)$, il existe un morphisme $f$ de $X$ dans un produit $X'$ de grassmanniennes tel que $Z$ soit dans l'image de $f^*$. Alors les homomorphismes canoniques
\[
 \varphi:A(X)\to\Gr(K(X)),\qquad \psi:\Gr(K(X))\to A(X)
\]
vérifient en dimension $i$ les relations
\begin{equation}\tag{2.35}
 \psi\varphi=(-1)^{i-1}(i-1)!\,1,
 \qquad
 \varphi\psi=(-1)^{i-1}(i-1)!\,1.
\end{equation}
Comme $\varphi$ est surjective, il suffit de prouver la première formule; on peut alors supposer $X$ produit de grassmanniennes. La seconde formule se ramène à
\[
 \gamma^i(x-\varepsilon(x))=(-1)^{i-1}(i-1)!\,x\mod K(X)^{i+1},
\]
pour tout $x\in K(X)^i$, puis à la relation analogue dans $A(X)$, conséquence de (1.22 bis).

\textbf{Remarque.} Le corollaire 2 rend plausible la validité de (2.35) dans tous les cas. On prouve en tout cas la deuxième formule modulo le sous-groupe de torsion de $\Gr(K(X))$, et en caractéristique $0$ elle résulte du théorème de Riemann--Roch. L'intérêt de (2.35) tient à ce qu'elle montre que $\varphi$ est injective modulo torsion; après tensorisation par $\Q$, les deux définitions des classes de Chern seraient identiques.

Soit toujours $X$ quasi-projective non singulière; soit de plus $Y$ une sous-variété non singulière de $X$ et soit $i$ l'injection de $Y$ dans $X$. Pour $i_*$ et $i^!$ on a les formules suivantes :
\begin{enumerate}[label=(\roman*)]
\item $i^!$ est un homomorphisme de $\lambda$-anneaux.
\item $i_*(x\cdot i^!(y))=i_*(x)\cdot y$, d'où $i_*i^!(y)=\xi y$, avec $\xi=i_*i^!(1)=\gamma_X(N)$.
\item $i^!i_*(x)=x\lambda_{-1}(N)$, d'où $i^!(\xi)=\lambda_{-1}(N)$.
\item $i_*(x)i_*(x')=i_*(xx'\lambda_{-1}(N))$.
\end{enumerate}
Par définition, $N$ est la classe dans $K(Y)$ du fibré conormal à $Y$ dans $X$. Les formules (i) et (ii) valent pour des morphismes quelconques; (iii) et (iv) se vérifient par (2.15) et la formule analogue pour les images inverses, en utilisant
\begin{equation}\tag{2.36}
 \Tor(\cO_Y,\cO_Y)=\Lambda(N).
\end{equation}
La formule (iv) exprime que $i_*$ est un homomorphisme d'anneaux si l'on introduit sur $K(Y)$ la multiplication $x\cdot_Nx'=xx'\lambda_{-1}(N)$. Elle se complète par
\begin{enumerate}[label=(\roman*)]
\setcounter{enumi}{4}
\item $i_*(\lambda^p(N,x))=\lambda^p(i_*(x))$ pour $p\geq1$,
\end{enumerate}
ce qui signifie que $i_*$ définit un $\lambda$-homomorphisme de $K(Y)_N$ dans $K(X)$. Cette formule, certaine dans tous les cas, n'est démontrée ici qu'en caractéristique $0$, par les théorèmes 1.4 et 2.8.

On a des formules analogues pour $i_*$ et $i^*$ :
\begin{enumerate}[label=(\roman* bis)]
\item $i^*$ est un homomorphisme d'anneaux gradués.
\item $i_*(x\cdot i^*(y))=i_*(x)y$, d'où $i_*i^*(y)=\xi'y$, avec $\xi'=i_*(1)$.
\item $i^*i_*(x)=x(-1)^qN^q$, d'où $i^*(\xi')=(-1)^qN^q$.
\item $i_*(x)i_*(x')=i_*(xx'(-1)^qN^q)$.
\item $i_*$ augmente les degrés de $q$ unités.
\end{enumerate}
Ici $q$ est la codimension de $Y$ dans $X$, et $N^i=c^i(N)$. Les formules (iii bis) et (iv bis) résultent par réduction des formules (ii) et (iv), en se plaçant dans $\Gr(K(Y))$ et $\Gr(K(X))$ et non dans $A(Y)$ et $A(X)$.

Toutes les formules universelles liant $i_*$ et $i^!$, respectivement $i_*$ et $i^*$, sont conséquences formelles des précédentes. De (v) résulte
\[
 i_*(\gamma^p(N,x))=\gamma^p(i_*(x)).
\]
Le théorème 2.14 et la proposition 1.5 montrent que $\gamma^p(N,x)$, pour $p=q+j$, est de filtration $\geq j$. Si l'on note $G_{q,j}(N,x)$ sa classe dans $\Gr^j(K(Y))$, le théorème 2.16 donne
\[
 i_*(G_{q,j}(N,x))=c^p(i_*(x)).
\]
Enfin, la proposition 1.5 jointe au théorème 2.14 permet d'exprimer $G_{q,j}(N,x)$ comme un polynôme à coefficients entiers en les $c^i(N)$, les $c^i(x)$ et $\varepsilon(x)$.

\textbf{Théorème 2.17.} Supposons le corps $k$ de caractéristique $0$ et soit $i:Y\to X$ un morphisme d'injection de variétés quasi-projectives non singulières. On a alors les formules (i) à (v), et de plus, pour tout entier $j\geq0$, la formule
\begin{equation}\tag{2.37}
 c^{q+j}(i_*(x))=i_*\left(\left(\frac{\widetilde c(x)*\lambda_{-1}(\widetilde c(N))}{(-1)^qN^q}\right)^{(j)}\right).
\end{equation}
Le deuxième membre de (2.37) désigne de façon abrégée le polynôme en les composantes de $\widetilde c(x)$ et $\widetilde c(N)$ décrit dans la proposition 1.5.

\textbf{Corollaire.} Dans les conditions précédentes, on a
\begin{equation}\tag{2.38}
 \ch(\widetilde c(i_*(x)))=i_*\left(\ch(\widetilde c(x))\mathfrak C(\check N)^{-1}\right),
\end{equation}
où $\ch$ est l'homomorphisme de Chern défini par (1.26).

Il suffit d'utiliser (2.37) écrit sous la forme
\begin{equation}\tag{2.37 bis}
 \widetilde c(i_*(x))=1+i_*\left(\frac{\widetilde c(x)*\lambda_{-1}(\widetilde c(N))-1}{(-1)^qN^q}\right),
\end{equation}
la formule explicite (1.29), la formule (iv bis), et enfin (1.30).

On peut écrire (2.38) sous la forme
\begin{equation}\tag{2.38 bis}
 \ch_X(F)=i_*\left(\ch_Y(F)\mathfrak C(T_{X/Y})^{-1}\right),
\end{equation}
où $F$ est un faisceau algébrique cohérent sur $Y$ considéré au premier membre comme un faisceau sur $X$ nul en dehors de $Y$, où $\ch_X$ et $\ch_Y$ sont les homomorphismes de Chern pris respectivement sur $X$ et $Y$, et où $T_{X/Y}$ est le fibré normal à $Y$ dans $X$.

\textbf{Remarques.} Contrairement à (2.38), la formule (2.37) est une formule sans dénominateur. Elle n'a été démontrée, même en caractéristique $0$, qu'en se plaçant dans $\Gr(K(X))$ et non dans $A(X)$, ce qui n'est peut-être pas la même chose. Elle est certainement vraie en toute caractéristique et dans $A(X)$. Les images des deux membres par $i^*$ sont les mêmes; en appliquant $i_*$, le produit de leur différence par $N^q$ est nul. Cela suggère une méthode de démonstration différente, en cherchant à obtenir $Y$ comme image réciproque transversale d'une sous-variété $Y'$ dans une variété $X'$, telle qu'en dimension $\leq n=\dim X$ la classe de $Y'$ dans $A(X')$ ne soit pas diviseur de zéro. Malheureusement on ne peut espérer que $X'$ soit une grassmannienne ou quelque chose d'approchant.

On notera que, si l'on néglige les éléments de torsion, la formule (2.38) est équivalente à (2.37).




% Batch 067 macro additions
\providecommand{\cS}{\mathcal S}
\providecommand{\cD}{\mathcal D}
\providecommand{\cE}{\mathscr E}
\providecommand{\cG}{\mathscr G}
\providecommand{\cL}{\mathscr L}
\providecommand{\cM}{\mathscr M}
\providecommand{\cP}{\mathscr P}
\providecommand{\cK}{\mathcal K}
\providecommand{\Db}{\mathrm D}
\providecommand{\Kb}{\mathrm K}
\providecommand{\Trc}{\operatorname{Tr}}
\providecommand{\ob}{\operatorname{ob}}
\providecommand{\Coh}{\operatorname{Coh}}
\providecommand{\Cone}{\operatorname{Cone}}
\providecommand{\im}{\operatorname{Im}}
\providecommand{\Id}{\operatorname{Id}}
\providecommand{\RR}{\mathrm{R.R.}}
\providecommand{\RRR}{\mathrm{R.R.R.}}
\providecommand{\GrR}{\mathrm{G}}
\providecommand{\pr}{\operatorname{pr}}


\section*{6. Le théorème de Riemann--Roch}

Son énoncé est le suivant : soit \(f\) un morphisme propre d'une variété quasi-projective non singulière \(Y\) dans une variété \(X\) de même espèce. Alors on a
\begin{equation}\tag{2.39}
 f_*\bigl(\ch_Y(x)\,\mathfrak C(Y)\bigr)
 =
 \ch_X\bigl(f_!(x)\bigr)\,\mathfrak C(X).
\end{equation}
Les deux membres sont dans \(A(X)\otimes\mathbb Q\) et \(\mathfrak C(X)=\mathfrak C(T_X)\).

Il permet de calculer, aux éléments de torsion près, les classes de Chern de \(f_!(x)\) connaissant les classes de Chern de \(x\), celles de \(T_X\) et \(T_Y\), et l'homomorphisme \(f_*\). Ou encore, il permet de déterminer l'homomorphisme
\[
  f_!:K(Y)\otimes\mathbb Q \longrightarrow K(X)\otimes\mathbb Q
\]
lorsqu'on identifie au moyen des homomorphismes \(\ch_Y\) et \(\ch_X\) ces groupes respectivement à \(A(Y)\otimes\mathbb Q\) et \(A(X)\otimes\mathbb Q\) (si la formule (2.35) est vraie), à l'aide de \(f_*\) et des classes de Chern de \(T_X\) et de \(T_Y\).

Je ne sais démontrer la formule (2.39) qu'en caractéristique \(0\), et comme égalité dans \(\GrR(K(X))\otimes\mathbb Q\). Pour ceci, \(Y\) étant localement fermé dans un espace projectif \(P\), on décompose \(f\) en le produit d'un morphisme d'injection de \(Y\) dans \(P\times X\), et d'un morphisme de projection de \(P\times X\) sur \(X\), et l'on démontre (2.39) pour chacun de ces deux morphismes.

Pour un morphisme d'injection, on voit immédiatement que les formules (2.38) et (2.39) sont équivalentes. L'hypothèse de caractéristique \(0\) est intervenue uniquement par l'intermédiaire du théorème 2.8 qui permet d'exprimer les \(\lambda^i(\gamma(F))\), pour un faisceau algébrique cohérent \(F\), par des invariants cohomologiques locaux de \(F\).

Pour un morphisme de projection \(f:P\times X\to X\), on peut, grâce à la proposition 2.13, supposer que \(x\) est de la forme \(y\otimes z\), avec \(y\in K(P)\) et \(z\in K(X)\). Alors, la formule de Künneth donne
\[
  f_!(x)=\chi(y)\,z,
\]
où, pour un faisceau algébrique cohérent \(F\) sur \(P\), l'on pose
\[
  \chi(F)=\sum_{i\geq 0}(-1)^i\dim H^i(P,F).
\]
On est alors aussitôt ramené à prouver la formule
\begin{equation}\tag{2.40}
  \chi(F)=\pi_r\bigl(\ch_p(F)\,\mathfrak C(P)\bigr),
\end{equation}
où \(\pi_r\) désigne la composante de degré \(r=\dim P\). En vertu de la proposition 2.13, on peut supposer que \(F=\mathcal O_Q\), où \(Q\) est une sous-variété linéaire de \(P\). Alors \(\chi(F)=1\), puisque c'est le genre arithmétique de \(Q\); reste à voir que le deuxième membre de la formule (2.40) vaut \(1\).

Si en effet \(V\) est un hyperplan de \(P\) et si l'on pose
\[
  \xi=\gamma_P(\mathcal O_V),\qquad L=\ell(V),
\]
on aura
\[
  1-\xi=\gamma_P(E(-L)),
\]
où \(E(-L)\) est le fibré vectoriel défini par \(-L\). D'où
\[
  C_p(1-\xi)=[1,1-L],
\]
et par suite
\[
  \ch_p(1-\xi)=e^{-L},\qquad \ch_p(\xi)=1-e^{-L},
\]
et donc
\[
  \ch_p(\xi^p)=(1-e^{-L})^p.
\]
Or si \(p\) est la dimension de \(Q\) dans \(P\), on a \(\xi^p=\gamma_P(Q)\); on en déduit
\[
  \ch_p(F)\,\mathfrak C(P)
  =
  \bigl(L/(1-e^{-L})\bigr)^{r+1}(1-e^{-L})^p.
\]
Il faut vérifier que le coefficient de \(L^r\) dans cette série formelle en \(L\) vaut \(1\), ce qui résulte de la propriété caractéristique de la série formelle \(L/(1-e^{-L})\), à savoir que le coefficient de \(L^q\) dans \(\bigl(L/(1-e^{-L})\bigr)^{q+1}\) vaut \(1\) (poser \(q=r+1-p\)).

\subsection*{Références bibliographiques}
\begin{enumerate}[label={[\arabic*]}]
\item Hirzebruch : \emph{Neue topologische Methoden in der algebraischen Geometrie}.
\item Serre : Faisceaux algébriques cohérents, \emph{Ann. of Math.}
\item Grothendieck : Sur les faisceaux algébriques cohérents, in Séminaire Cartan 1956/57.
\end{enumerate}

\section*{Démonstration du théorème de Riemann--Roch}
\begin{center}
(Grothendieck, 1er Novembre 1957)
\end{center}

\subsection*{Rappel de notations}
Ce sont celles du « rapport sur Riemann--Roch », cité \(\RRR\), dans ce qui suit.

Toutes les variétés considérées sont non singulières et quasi-projectives, c'est-à-dire isomorphes à une sous-variété localement fermée d'un espace projectif. Le corps de base est algébriquement clos, de caractéristique quelconque.

Si \(X\) est une telle variété, on note \(K(X)\) le groupe abélien engendré par les classes de faisceaux cohérents sur \(X\), modulo l'identification d'une extension à une somme; la somme alternée des Tor fait de \(K(X)\) un anneau. Si \(F\) est un faisceau, respectivement \(E\) un fibré à fibre vectorielle, respectivement \(Y\) une sous-variété de \(X\), on note \(\gamma_X(F)\), respectivement \(\gamma_X(E)\), respectivement \(\gamma_X(Y)\), l'élément de \(K(X)\) correspondant à \(F\), respectivement à \(E\), respectivement à \(\mathcal O_Y\). En fait, on écrit souvent \(F,E\) au lieu de \(\gamma_X(F)\).

On note \(A(X)\) l'anneau gradué des classes de cycles de \(X\) pour l'équivalence rationnelle (cf. Chow); \(A^i(X)\) est formé des cycles de codimension \(i\). Si \(E\) est un fibré à fibre vectorielle, on note \(C(E)=\sum C^i(E)\), \(C^i(E)\in A^i(X)\), sa classe de Chern (définie comme image réciproque des cycles de Schubert); on en déduit par linéarité les \(C^i(y)\), pour tout \(y\in K(X)\). Si
\[
  C(y)=\prod(1+a_i),
\]
on pose, suivant Hirzebruch,
\[
  \ch(y)=\sum e^{a_i},\qquad
  T(y)=\prod\frac{a_i}{1-e^{-a_i}}.
\]
Ce sont des éléments de \(A(X)\otimes\mathbb Q\). Lorsque \(y\) est le fibré tangent à \(X\), on écrit \(T(X)\) au lieu de \(T(y)\).

Si \(G\) est un fibré à fibre vectorielle, on note \(\lambda^pG\) sa puissance extérieure \(p\)-ème, et
\[
  \lambda_{-1}(G)=\sum(-1)^p\lambda^p(G);
\]
c'est un élément de \(K(X)\).

Si \(f\) est une application régulière propre d'une variété \(Y\) dans une variété \(X\), on définit des homomorphismes additifs
\[
  f^*:A(X)\longrightarrow A(Y),\qquad
  f_*:A(Y)\longrightarrow A(X).
\]
Le premier est homogène de degré \(0\) et multiplicatif. Le second est homogène de degré \(\dim X-\dim Y\). On définit également des homomorphismes
\[
  f^*:K(X)\longrightarrow K(Y),\qquad
  f_!:K(Y)\longrightarrow K(X).
\]
Si \(F\) est un faisceau cohérent sur \(X\), \(f^*(F)\) est égal, par définition, à la somme alternée des \(\Tor_p^{\mathcal O_X}(\mathcal O_Y,F)\). Pour les fibrés à fibre vectorielle, c'est simplement l'opération d'image réciproque.

Si \(G\) est un faisceau cohérent sur \(Y\), \(f_*(G)\) est cohérent sur \(X\). Si l'on note \(R^pf_*\) le \(p\)-ème foncteur dérivé du foncteur \(f_*\) ainsi défini, on a, par définition,
\[
  f_!(G)=\sum(-1)^p R^pf_*(G).
\]

\subsection*{Théorème de Riemann--Roch}
Soit \(f:Y\to X\) un morphisme, application régulière, propre d'une variété quasi-projective non singulière \(Y\) dans une autre \(X\), et soit \(y\in K(Y)\). On a :
\begin{equation}\tag{1}
  f_*\bigl(\ch(y)T(Y)\bigr)=\ch(f_!(y))T(X)
  \quad\text{dans } A(X)\otimes\mathbb Q.
\end{equation}

\emph{Démonstration} abrégée.

\textbf{Lemme 1.} Considérons des morphismes propres de variétés
\[
  Z \xrightarrow{g} Y \xrightarrow{f} X.
\]
Si le théorème \(\RR\) est vrai pour \(f\) et \(g\), il l'est pour \(fg\). S'il est vrai pour \(fg\) et \(g\), il l'est pour \(f\). Si \(f_*\) est injective, et si le théorème \(\RR\) est vrai pour \(fg\) et \(f\), il l'est pour \(g\).

Trivial. En fait, lorsqu'on suppose \(\RR\) vrai pour \(fg\) et \(g\), il faut encore supposer que \(g_!\) est surjectif pour pouvoir affirmer que \(\RR\) est vrai pour \(f\).

Un morphisme propre étant composé d'une injection et d'une projection \(X\times P\to X\), où \(P\) est un espace projectif, on est ramené à traiter séparément ces deux cas. Le deuxième, on l'a vu, se ramène au calcul standard, moyennant la proposition 2.13 (« décomposition cellulaire ») de \(\RRR\). On va donc supposer \(Y\subset X\), avec le morphisme d'injection \(i\). La formule (1) devient alors :
\begin{equation}\tag{2}
  \ch(i_!(y))=i_*\bigl(\ch(y)T(E)^{-1}\bigr),
\end{equation}
où \(E=T_{X/Y}\) est le fibré normal à \(Y\) dans \(X\).

\textbf{Lemme 2.} Supposons qu'il existe sur \(X\) un fibré \(G\) à fibre vectorielle de rang \(p=\codim_XY\), tel que
\begin{equation}\tag{3}
  \gamma_X(Y)=i_*(1)=\lambda_{-1}(G),\qquad
  i^*(G)=E,\qquad
  (-1)^pC^p(G)=Y.
\end{equation}
Si alors on a
\begin{equation}\tag{4}
  y=i^*(x),\qquad x\in K(X),
\end{equation}
la formule (2) est valable.

Démonstration immédiate, à partir du fait que \(\ch\) est un homomorphisme d'anneaux, et que l'on a
\begin{equation}\tag{5}
  \ch(\lambda_{-1}(G))=(-1)^pC^p(G)T(G)^{-1}.
\end{equation}

\textbf{Corollaire.} \(\RR\) est vrai si \(Y\) est une hypersurface et si \(y\) est de la forme \(i^*(x)\), avec \(x\in K(X)\).

\textbf{Lemme 3.} Supposons que \(X=Y\times P\), \(P\) étant un espace projectif, et que \(i\) soit l'injection \(u\mapsto (u,v_0)\), \(v_0\in P\). Alors \(i_*\) est injective et \(\RR\) est vrai.

C'est du pur calcul standard.

Supposons de nouveau \(X,Y,y\) quelconques, toujours avec \(Y\subset X\). Soit \(X'\) la variété éclatée de \(X\) à l'aide de \(Y\), et \(Y'\) l'hypersurface image réciproque de \(Y\). On a un diagramme d'applications
\[
\begin{tikzcd}
Y' \arrow[r,"j"] \arrow[d,"g"'] & X' \arrow[d,"f"]\\
Y \arrow[r,"i"'] & X .
\end{tikzcd}
\]
L'image réciproque de \(E\) dans \(Y'\) admet un sous-fibré de rang \(1\), dont le dual est noté \(L\). D'ailleurs \(L^{-1}\) est aussi la restriction à \(Y'\) du fibré \(L(Y')\) défini sur \(X'\) par le diviseur \(Y'\). On pose \(F=E/L^{-1}\). Des calculs élémentaires donnent le lemme suivant.

\textbf{Lemme 4.} On a les formules, où \(p=\codim_XY\) :
\begin{align}
  f^*i_!(y)&=j_*\bigl(\lambda_{-1}(\check F)y\bigr), \tag{6}\\
  \lambda_{-1}(\check F)&\equiv \gamma^{p-1}(E)\lambda^p(\check E)
    \pmod{1-L}, \tag{7}\\
  f_*(1)&=1,\quad\text{d'où }f_*f^*=1, \tag{8}\\
  f^*i_*(\alpha)&\equiv j_*\bigl(\alpha C^{p-1}(F)\bigr)
    \pmod{\Ker f_*}\quad\text{pour tout }\alpha\in A(Y). \tag{9}
\end{align}
Pour simplifier l'écriture, on a fait opérer \(K(Y)\) sur \(K(Y')\) et \(A(Y)\) sur \(A(Y')\) grâce aux homomorphismes \(g^*\). Dans la formule (7), il semble que \(\gamma^{p-1}(E)\) désigne \(\lambda^{p-1}(E-2)\), c'est-à-dire
\[
  \sum_{i+j=p-1}(\,\cdots\,)\lambda^i(E);
\]
la congruence a lieu dans \(K(Y')\).

La formule (6) résulte de l'égalité :
\begin{equation}\tag{6 bis}
  \Tor_i^{\mathcal O_X}\bigl(\mathcal O_Y(v),\mathcal O_Y(v')\bigr)
  =
  \lambda^i\check F(v')
  \quad (v'\in Y',\ v=g(v')).
\end{equation}
La formule (7) résulte des formules
\begin{equation}\tag{10}
  g_*(L^{-i})=0\quad\text{pour }0<i\leq p-1,\qquad g_*(1)=1,
\end{equation}
elles-mêmes conséquences de Künneth et de la cohomologie des espaces projectifs.

La formule (8) est triviale. Enfin (9) résulte de \(g_*(C^{p-1}(F))=1\), qui se déduit, par exemple, des formules immédiates :
\begin{equation}\tag{10 bis}
  g_*(\xi^i)=0\quad\text{si }0\leq i<p-1,\qquad
  g_*(\xi^{p-1})=1,
\end{equation}
avec \(\xi=C^1(L)\in A^1(Y')\). En fait, (9) est très vraisemblablement une identité, et pas seulement une congruence. Cette formule inutilement compliquée peut se remplacer par exemple par VII 4.1.

De la formule (7) résulte aussitôt que \(\lambda_{-1}(\check F)\in j^*(K(X'))\), pourvu que l'on ait
\begin{equation}\tag{11}
  y\,\gamma^{p-1}(E)\lambda^p(\check E)\in i^*(K(X)).
\end{equation}
Supposons cette condition (11) satisfaite. D'après le corollaire au lemme 2, \(\RR\) est vrai pour \((Y',X',y\lambda_{-1}(\check F))\). D'autre part, pour prouver (2), la formule (8) ramène à prouver que
\begin{equation}\tag{12}
  f^*\ch(i_!(y))
  \equiv
  f^*\bigl(i_*(\ch(y)T(E)^{-1})\bigr)
  \pmod{\Ker f_*}.
\end{equation}
Or, cette dernière formule se vérifie aussitôt, utilisant (6), (9), \(\RR\) pour \((Y',X',y\lambda_{-1}(\check F))\) et la formule :
\begin{equation}\tag{13}
  \ch(\lambda_{-1}(\check F))=C^{p-1}(F)T(F)^{-1}
\end{equation}
(cf. formule (5)).

Ainsi, \(\RR\) est prouvé chaque fois que \(y\,\gamma^{p-1}(E)=0\), en particulier chaque fois que \(p-1>\dim Y\), puisque \(\gamma^{p-1}(E)\) est de filtration \(\geq p-1\) dans \(K(Y)\), filtré par la codimension des supports; c'est-à-dire chaque fois que
\[
  \dim Y < \frac{\dim X-1}{2}.
\]
Or, on se ramène à ce cas, en conjuguant le lemme 1 et le lemme 3, prenant pour \(P\) un espace projectif de dimension \(r\) assez grande, par exemple \(r>\dim X+1\). C.Q.F.D.

\textbf{N.B.} On s'est servi du fait que \(\gamma^{p-1}(E)\) était de filtration \(\geq p-1\). On peut prouver ce fait directement, sans passer par les grassmanniennes comme dans \(\RRR\), 2.16, par l'étude du fibré projectif associé à \(E\). Cette méthode donne plus généralement le théorème 2.14 de \(\RRR\).

\subsection*{Extrait d'une lettre de Grothendieck}
« moralement », la nouvelle méthode repose sur la détermination de \(K(X)\) et de \(A(X)\) lorsque \(X\) est, soit un fibré projectif associé à un fibré vectoriel, soit une variété obtenue par éclatement d'une sous-variété non singulière. Dans ce deuxième cas, je n'ai pas fait la détermination complète; il manque un petit quelque chose. Avec les notations du petit papier, il faut prouver la formule (9), mais comme égalité :
\begin{equation}\tag{9}
  f^*(i_*(\alpha))=j_*(\alpha C^{p-1}(F)),\qquad \alpha\in A(Y).
\end{equation}
On constate que les deux membres ont même image par \(f_*\), et même image par \(j^*\), de sorte que la question est équivalente à :
\begin{equation}\tag{9 bis}
  \Ker f_*\cap\Ker j^*=0.
\end{equation}
En tout cas, (9) est vraie en réduisant à \(G(K(X'))\), ce qui détruit un peu de torsion, tout au plus, car elle résulte alors de (6); il n'y a donc guère de doute qu'elle ne soit vraie.

Je te signale d'ailleurs que l'égalité (9) résulte très facilement de (6) sans passer par l'ennuyeuse (6 bis).

\begin{flushleft}
P.C.C.J-P.S.
\end{flushleft}

\newpage
\section*{Exposé I. Généralités sur les conditions de finitude dans les catégories dérivées}
\begin{center}
par L. Illusie
\end{center}

\subsection*{0. Introduction}
L'objet des Exposés I à IV est de développer, avec la généralité qui convient dans ce Séminaire, le formalisme des conditions de finitude dans les catégories dérivées des topos annelés.

Comme il a été dit dans l'Exposé 0, c'est le besoin de définir sur des schémas arbitraires des « groupes de Grothendieck » possédant de bonnes propriétés de variance qui a obligé à généraliser et à assouplir les notions de finitude utilisées jusqu'à présent.

Une notion classique comme celle de faisceau cohérent sur un espace annelé \((X,\mathcal O_X)\) devient sans intérêt dès que \(\mathcal O_X\) n'est plus cohérent. On pourrait songer à la remplacer par la notion de présentation finie, mais celle-ci présente l'inconvénient que le noyau d'un épimorphisme de modules de présentation finie n'est plus en général de présentation finie. On arrive à une notion satisfaisante en remarquant que, si \(\mathcal O_X\) est cohérent, un faisceau cohérent \(F\) est non seulement de présentation finie mais de \(n\)-présentation finie pour tout \(n\in\mathbb N\), « \(n\)-présentation finie » voulant dire qu'il existe localement une suite exacte
\[
  L_n\longrightarrow L_{n-1}\longrightarrow \cdots \longrightarrow L_0\longrightarrow F\longrightarrow 0,
\]
où les \(L_i\) sont libres de type fini. Le faisceau \(\mathcal O_X\) n'étant plus supposé cohérent, disons qu'un \(\mathcal O_X\)-module \(F\) est \emph{pseudo-cohérent} s'il vérifie la condition précédente, c'est-à-dire s'il est de \(n\)-présentation finie pour tout \(n\). Alors la notion de pseudo-cohérence possède, vis-à-vis des suites exactes courtes, la même propriété de stabilité que la notion de cohérence : si deux des termes d'une suite exacte courte de \(\mathcal O_X\)-modules sont pseudo-cohérents, le troisième l'est aussi.

On généralise aisément la notion précédente aux complexes de faisceaux. Un complexe \(F\) de \(\mathcal O_X\)-modules est dit \emph{pseudo-cohérent} s'il est \(n\)-pseudo-cohérent pour tout \(n\in\mathbb Z\), « \(n\)-pseudo-cohérent » signifiant qu'il existe localement un quasi-isomorphisme \(L\to F\), où \(L\) est un complexe à degrés bornés supérieurement et à composantes libres de type fini en degré \(>n\). Pour un complexe \(F\) réduit au degré \(0\), il revient au même de dire que \(F\) est pseudo-cohérent en tant que complexe ou en tant que module.

La notion de complexe pseudo-cohérent a d'excellentes vertus de stabilité, décrites dans l'Exposé I §§1 et 2. D'abord c'est une notion stable par isomorphisme dans la catégorie dérivée \(\Db(X)\) de la catégorie des \(\mathcal O_X\)-modules; qui mieux est, si deux des objets d'un triangle distingué de \(\Db(X)\) sont pseudo-cohérents, le troisième l'est aussi : en d'autres termes, la sous-catégorie pleine \(\Db(X)_{\mathrm{coh}}\) de \(\Db(X)\) formée des complexes pseudo-cohérents est une sous-catégorie triangulée. En outre, la pseudo-cohérence se conserve par image inverse et produit tensoriel dérivés. Dans le cas où \(\mathcal O_X\) est cohérent, dire qu'un complexe \(F\) est pseudo-cohérent signifie simplement que \(F\) est à cohomologie localement bornée supérieurement et que tous ses \(H^i(F)\) sont cohérents.

Si \(\mathcal O_X\) est un faisceau d'anneaux locaux réguliers, tout \(\mathcal O_X\)-module cohérent \(F\) admet localement une résolution finie à gauche par des modules libres de type fini, c'est-à-dire qu'il existe localement une suite exacte
\[
  0\longrightarrow L_n\longrightarrow \cdots \longrightarrow L_0\longrightarrow F\longrightarrow 0,
\]
où les \(L_i\) sont libres de type fini. Il est donc naturel, quand \(\mathcal O_X\) est un faisceau d'anneaux locaux quelconque, de s'intéresser aux modules qui jouissent de la propriété précédente. De tels modules sont appelés parfaits. Plus généralement, on dit qu'un complexe \(F\) de \(\mathcal O_X\)-modules est \emph{parfait} s'il existe localement un quasi-isomorphisme \(L\to F\), où \(L\) est un complexe à degrés bornés et à composantes libres de type fini. On montre dans l'Exposé I §3 que la sous-catégorie pleine \(\Db(X)_{\mathrm{parf}}\) de \(\Db(X)\) formée des complexes parfaits est, de même que \(\Db(X)_{\mathrm{coh}}\), une sous-catégorie triangulée, « stable » par image inverse et produit tensoriel dérivés. On a évidemment
\[
  \Db(X)_{\mathrm{parf}}\subset\Db(X)_{\mathrm{coh}},
\]
mais l'inclusion est stricte en général. La pseudo-cohérence peut se redéfinir « par passage à la limite » à partir de la perfection : un complexe \(F\) est pseudo-cohérent si et seulement si, pour tout \(n\in\mathbb Z\), \(F\) peut être « localement approché à l'ordre \(n\) par un complexe parfait », par quoi l'on entend qu'il existe localement un triangle distingué
\[
\begin{tikzcd}
L \arrow[rr] && F \arrow[dl]\\
& R \arrow[ul] &
\end{tikzcd}
\]
où \(L\) est parfait et \(R\) acyclique en degré \(>n\). Inversement, la perfection se récupère à partir de la pseudo-cohérence grâce à une condition de régularité supplémentaire : un complexe est parfait si et seulement s'il est pseudo-cohérent et localement de tor-dimension finie (Exposé I §5). On retrouve ainsi le fait que, si les anneaux locaux sont réguliers, tout faisceau cohérent est parfait, plus généralement tout complexe pseudo-cohérent à cohomologie localement bornée est parfait.

Pseudo-cohérence et perfection sont les deux notions de finitude fondamentales avec lesquelles nous travaillerons dans ce Séminaire. Les §§1 à 5 de l'Exposé I sont consacrés à leur définition et à l'établissement de leurs propriétés élémentaires de stabilité. Pour ceci, on utilise seulement deux ou trois propriétés locales de la catégorie des modules libres de type fini vis-à-vis de la catégorie de tous les modules. Aussi s'avère-t-il commode et utile, notamment dans l'Exposé II, d'axiomatiser la situation, en introduisant une notion de pseudo-cohérence (respectivement de perfection) dans une catégorie fibrée sur un site relativement à une sous-catégorie fibrée vérifiant des conditions convenables.

Les §§6 à 8 de l'Exposé I généralisent aux complexes parfaits un certain nombre de notions bien connues pour les faisceaux localement libres de type fini : rang, dualité, trace d'un endomorphisme.

Dans l'Exposé II, on examine le problème de l'existence de « résolutions globales » : à quelles conditions, sur un topos annelé \((X,\mathcal O_X)\), un complexe pseudo-cohérent (respectivement parfait) est-il globalement isomorphe dans \(\Db(X)\) à un complexe de modules localement libres de type fini? L'importance de cette question réside essentiellement dans le fait qu'on ne sait pas, pour le moment, généraliser aux complexes parfaits certaines constructions usuelles sur les modules localement libres, comme puissances extérieures et symétriques, classes de Chern. C'est pourquoi il est commode d'avoir à sa disposition des critères suffisants maniables pour l'existence de résolutions globales, ce qui permet de ramener certaines questions sur les complexes pseudo-cohérents (respectivement parfaits) à des questions analogues sur les faisceaux localement libres de type fini. De tels critères sont donnés dans l'Exposé II.

On montre en particulier qu'il existe des résolutions globales dans les cas suivants : \(X\) est le topos zariskien d'un schéma affine, ou plus généralement d'un schéma quasi-compact possédant un faisceau inversible ample, voire une famille ample de faisceaux inversibles au sens de Kleiman (par exemple un schéma régulier); ou encore \(X\) est le topos des faisceaux d'ensembles sur un espace topologique compact, \(\mathcal O_X\) le faisceau des fonctions continues à valeurs complexes. À titre d'illustration, et pour mettre en évidence la souplesse des notions introduites, on donne en appendice une définition « purement faisceautique » de l'indice d'une famille d'opérateurs elliptiques.

L'Exposé III étudie la stabilité des conditions de finitude par image directe dérivée. Pour obtenir des énoncés raisonnables, il y a lieu de « relativiser » les notions de pseudo-cohérence et de perfection. On se place ici dans le cadre des schémas ordinaires, ce qui suffit pour le Séminaire, mais il ne fait pas de doute qu'il faudra tôt ou tard développer une théorie analogue pour les schémas ou espaces analytiques relatifs. Soit \(p:X\to S\) un \(S\)-schéma localement de type fini et soit \(F\) un complexe de \(\mathcal O_X\)-modules. On dit que \(F\) est pseudo-cohérent (respectivement parfait) relativement à \(p\) si l'on peut localement plonger, par une immersion fermée, \(X\) dans un \(S\)-schéma lisse \(X'\) de manière que \(F\) prolongé par zéro sur \(X'\) soit pseudo-cohérent (respectivement parfait). La notion de pseudo-cohérence relativement à \(p\) est surtout intéressante dans le cas où \(S\) n'est pas localement noethérien, car dans le cas contraire elle coïncide avec la notion de pseudo-cohérence ordinaire. D'autre part, comme on peut s'y attendre, dire que \(F\) est parfait relativement à \(p\) équivaut à dire que \(F\) est pseudo-cohérent relativement à \(p\) et localement de tor-dimension finie relativement à \(p\), c'est-à-dire relativement au faisceau d'anneaux \(p^{-1}(\mathcal O_S)\).

Le théorème central de l'Exposé III est le théorème de finitude qui affirme, sous une forme un peu plus précise, que, si \(f:X\to Y\) est un morphisme propre de \(S\)-schémas localement de type fini, le foncteur \(Rf_*\) transforme complexes pseudo-cohérents relativement à \(S\) en complexes pseudo-cohérents relativement à \(S\); ce théorème est en réalité une conjecture, mais est prouvé néanmoins dans les deux cas particuliers suivants : \(S\) est localement noethérien; \(f\) est projectif. Il semble malheureusement que l'extension au cas général soit du même ordre de difficulté que le théorème analogue en géométrie analytique, le théorème de Grauert. Combinant le théorème de finitude avec une formule essentiellement triviale dite formule de projection, on obtient des critères maniables pour la stabilité de la perfection relative par image directe. On retrouve, en corollaire, les théorèmes de continuité et de semi-continuité de Grauert (EGA III 7.6).

L'Exposé IV traduit en langage de « groupes de Grothendieck » les résultats des Exposés I à III. Sur un topos annelé \((X,\mathcal O_X)\), on note \(K^*(X)\) (respectivement \(K_\bullet(X)\)) le groupe de Grothendieck de la catégorie des complexes parfaits et de tor-dimension finie (respectivement pseudo-cohérents et à cohomologie bornée). Le groupe \(K^*(X)\) est un anneau, et \(K_\bullet(X)\) un module sur \(K^*(X)\). \(K^*\) est un foncteur contravariant, comme le suggère le point en haut, tandis que \(K_\bullet\) est un foncteur covariant pour les morphismes propres et pseudo-cohérents de schémas ou d'espaces analytiques, et sous réserve que l'on dispose du théorème de finitude. Il y a aussi des variantes inhabituelles : covariance de \(K^*\), contravariance de \(K_\bullet\), pour des morphismes satisfaisant des hypothèses de régularité convenables. Enfin, les critères de globalisation de l'Exposé II permettent de faire le lien avec les « groupes de Grothendieck » définis naïvement à partir des faisceaux cohérents ou des faisceaux localement libres de type fini.

\section*{1. Définitions préliminaires}

\subsection*{1.1. Catégories fibrées à fibres additives (resp. abéliennes, resp. triangulées)}

\textbf{1.1.1.} Soit \(S\) une catégorie, et soit \(\mathcal C\) une \(S\)-catégorie fibrée (SGA 1 VI 6.1) à fibres des catégories additives (resp. triangulées). Nous dirons que \(\mathcal C\) est une \(S\)-catégorie additive (resp. triangulée), ou que \(\mathcal C\) est additive (resp. triangulée) sur \(S\), si, pour toute flèche \(f:X\to Y\) de \(S\), le foncteur image inverse
\[
  f^*:\mathcal C_Y\longrightarrow \mathcal C_X
\]
déterminé à isomorphisme unique près est additif (resp. exact). Nous dirons que \(\mathcal C\) est une \(S\)-catégorie abélienne si \(\mathcal C\) est une \(S\)-catégorie additive à fibres abéliennes, et que \(\mathcal C\) est une \(S\)-catégorie abélienne plate si en outre les foncteurs image inverse sont exacts.

\textbf{1.1.2.} Soit \(\mathcal C\) une \(S\)-catégorie triangulée. On définit, par exemple à l'aide d'un clivage de \(\mathcal C\) (SGA 1 VI 7.1), une \(S\)-catégorie fibrée \(\Trc(\mathcal C)\) dont la fibre en chaque \(X\in \ob(S)\) est la catégorie \(\Trc(\mathcal C_X)\) des triangles distingués de \(\mathcal C_X\).

\textbf{1.1.3.} Soient \(\mathcal C\) et \(\mathcal C'\) des \(S\)-catégories additives (resp. triangulées) et \(F:\mathcal C\to\mathcal C'\) un foncteur. On dit que \(F\) est un \(S\)-foncteur additif (resp. exact) si \(F\) est un \(S\)-foncteur cartésien additif (resp. exact) sur chaque fibre.



% Batch 068 macro additions
\providecommand{\cC}{\mathcal C}
\providecommand{\cD}{\mathcal D}
\providecommand{\cK}{\mathcal K}
\providecommand{\cO}{\mathcal O}
\providecommand{\cS}{\mathcal S}
\providecommand{\Cpx}{\mathrm C}
\providecommand{\Kpx}{\mathrm K}
\providecommand{\Dpx}{\mathrm D}
\providecommand{\Ho}{\mathrm H}
\providecommand{\Ob}{\operatorname{ob}}
\providecommand{\Coker}{\operatorname{Coker}}
\providecommand{\Ker}{\operatorname{Ker}}
\providecommand{\Im}{\operatorname{Im}}
\providecommand{\Hom}{\operatorname{Hom}}
\providecommand{\Cone}{\operatorname{Cone}}
\providecommand{\id}{\operatorname{Id}}
\providecommand{\qis}{\xrightarrow{\sim}}


\textbf{1.1.3 (suite).} Si \(\mathcal C\) et \(\mathcal C'\) sont des \(S\)-catégories triangulées et si \(F\) est un \(S\)-foncteur exact, alors \(F\) définit un \(S\)-foncteur cartésien
\[
  \Trc(\mathcal C)\longrightarrow \Trc(\mathcal C').
\]

\textbf{1.1.4.} Soit \(\mathcal C\) une \(S\)-catégorie triangulée, et soit \(\mathcal C'\) une sous-\(S\)-catégorie fibrée pleine de \(\mathcal C\). On dit que \(\mathcal C'\) est une \emph{sous-\(S\)-catégorie triangulée} de \(\mathcal C\) si, pour tout \(X\in\Ob S\), \(\mathcal C'_X\) est une sous-catégorie triangulée de \(\mathcal C_X\) \([V]\), chap. I, §1, 2.3.

\textbf{1.1.5.} Soit \(\mathcal C\) une \(S\)-catégorie additive. On définit (SGA 4 XVII 2.2, 2.4) une \(S\)-catégorie additive (resp. triangulée)
\[
  \Cpx(\mathcal C)\qquad\text{(resp. }\Kpx(\mathcal C)\text{)}
\]
dont la fibre en \(X\in\Ob S\) est la catégorie \(\Cpx(\mathcal C_X)\) (resp. \(\Kpx(\mathcal C_X)\)) des complexes à différentielle de degré \(+1\) (resp. complexes « à homotopie près ») de \(\mathcal C_X\). Si \(\mathcal C\) est une \(S\)-catégorie abélienne plate, on définit (loc. cit.) une \(S\)-catégorie triangulée \(\Dpx(\mathcal C)\), dont la fibre en \(X\in\Ob S\) est la catégorie dérivée \(\Dpx(\mathcal C_X)\); si \(f:X\to Y\) est une flèche de \(S\), nous noterons
\[
  f^*:\Dpx(\mathcal C_Y)\longrightarrow \Dpx(\mathcal C_X)
\]
le foncteur image inverse, qui se déduit de
\[
  f^*:\Kpx(\mathcal C_Y)\longrightarrow \Kpx(\mathcal C_X)
\]
par simple passage au quotient. On définit des sous-\(S\)-catégories triangulées pleines
\[
  \Kpx^*(\mathcal C)\qquad\text{(resp. }\Dpx^*(\mathcal C)\text{)}
  \quad (*=+,-,b)
\]
par les conditions habituelles de degré.

\subsection*{1.2. Terminologie}

Soient \(S\) un site, \(\mathcal C\) une \(S\)-catégorie abélienne plate (1.1), \(\mathcal C_0\) une sous-\(S\)-catégorie fibrée strictement pleine.

Soit \(F\in\Ob\mathcal C_X\), pour \(X\in\Ob S\). L'ensemble des \(f:U\to X\) de \(S/X\) tels qu'il existe un épimorphisme
\[
  E\longrightarrow f^*F,\qquad E\in\Ob\mathcal C_{0,U},
\]
est un crible de \(X\). Si celui-ci est couvrant, nous dirons que \(F\) est de \(\mathcal C_0\)-\emph{type fini} (ou simplement de type fini s'il n'y a pas de confusion à redouter).

Nous dirons que \(\mathcal C_0\) est \emph{localement stable par noyau d'épimorphisme} si la condition suivante est réalisée : pour tout \(X\in\Ob S\) et tout épimorphisme
\[
  u:E\longrightarrow F
\]
de \(\mathcal C_X\), avec \(E\) et \(F\in\Ob\mathcal C_{0,X}\), le noyau de \(u\) est localement dans \(\mathcal C_0\), c'est-à-dire que l'ensemble des \(f:U\to X\) de \(S/X\) tels que
\[
  f^*\Ker(u)\in\Ob\mathcal C_{0,U}
\]
est un raffinement de \(X\).

Par le procédé standard de décomposition d'une suite exacte longue en suites exactes courtes, on voit que la condition précédente équivaut à la condition, en apparence plus forte, que voici : pour tout \(X\in\Ob S\) et toute suite exacte de \(\mathcal C_X\)
\begin{equation}\tag{*}
  0\longrightarrow E^0\longrightarrow E^1\longrightarrow\cdots\longrightarrow E^n\longrightarrow 0
\end{equation}
avec \(E^i\in\Ob\mathcal C_{0,X}\) pour \(1\leq i\leq n\), \(E^0\) est localement dans \(\mathcal C_0\).

Nous dirons que \(\mathcal C_0\) est \emph{localement quasi-stable par noyau d'épimorphismes} si, pour toute suite exacte telle que \((*)\), \(E^0\) est de \(\mathcal C_0\)-type fini.

D'autre part, nous dirons que \(\mathcal C_0\) est \emph{localement relevable dans} \(\mathcal C\) si, pour tout \(X\in\Ob S\), tout épimorphisme
\[
  u:E\longrightarrow F
\]
de \(\mathcal C_X\), avec \(F\in\Ob\mathcal C_{0,X}\), admet localement une section. Pour que \(\mathcal C_0\) soit localement relevable dans \(\mathcal C\), il faut et il suffit que, pour tout diagramme de \(\mathcal C_X\) (\(X\in\Ob S\))
\[
\begin{tikzcd}
& H \arrow[d]\\
E \arrow[r,"u"'] & F
\end{tikzcd}
\]
où \(u\) est un épimorphisme et \(H\in\Ob\mathcal C_0\), l'ensemble des \(U\) au-dessus de \(X\) tel qu'il existe un diagramme commutatif
\[
\begin{tikzcd}
& H_U \arrow[d]\\
E_U \arrow[r,"u_U"'] \arrow[ur] & F_U
\end{tikzcd}
\]
soit un raffinement de \(X\).

Nous dirons que \(\mathcal C_0\) est \emph{localement quasi-relevable dans} \(\mathcal C\) si la condition plus faible suivante est réalisée : pour tout \(X\in\Ob S\), et tout épimorphisme
\[
  u:E\longrightarrow F
\]
de \(\mathcal C_X\), avec \(F\in\Ob\mathcal C_{0,X}\), l'ensemble des \(f:U\to X\) de \(S/X\) tels qu'il existe
\[
  v:F'\longrightarrow f^*E,\qquad F'\in\Ob\mathcal C_{0,U},
\]
tel que \(f^*(u)v\) soit un épimorphisme, est un raffinement de \(X\). Cette condition équivaut à la suivante : pour tout diagramme de \(\mathcal C_X\)
\[
\begin{tikzcd}
& H \arrow[d]\\
E \arrow[r,"u"'] & F ,
\end{tikzcd}
\]
où \(u\) est un épimorphisme, et \(H\in\Ob\mathcal C_0\), l'ensemble des \(U\) au-dessus de \(X\) tels qu'il existe un diagramme commutatif
\[
\begin{tikzcd}
G \arrow[r,"v"] \arrow[d] & H_U \arrow[d]\\
E_U \arrow[r,"u_U"'] & F_U
\end{tikzcd}
\]
avec \(G\in\Ob\mathcal C_0\) et \(v\) un épimorphisme, est un raffinement de \(X\).

Si \(S\) est un site chaotique, par exemple le site ponctuel (un objet, un morphisme), nous omettrons l'adverbe « localement » dans les définitions précédentes.

\subsection*{1.3. Exemples}

\textbf{a)} Soient \(S\) un site ponctuel, \(A\) un anneau, \(\mathcal C\) la catégorie des \(A\)-modules à gauche. La sous-catégorie pleine \(\mathcal C_0\) de \(\mathcal C\) formée des \(A\)-modules à gauche projectifs de type fini est stable par noyau d'épimorphisme et relevable dans \(\mathcal C\); un \(A\)-module est de \(\mathcal C_0\)-type fini si et seulement s'il est de type fini au sens ordinaire. La sous-catégorie pleine \(\mathcal C_1\) de \(\mathcal C\) formée des \(A\)-modules de type fini est quasi-relevable dans \(\mathcal C\), mais non relevable en général; si \(A\) est noethérien, \(\mathcal C_1\) est stable par noyau d'épimorphisme. La sous-catégorie pleine \(\mathcal C'_0\) de \(\mathcal C\) formée des \(A\)-modules libres de type fini est relevable dans \(\mathcal C\) et quasi-stable, non stable en général, par noyau d'épimorphisme.

\textbf{b)} Soient \((S,\mathcal O_S)\) un espace annelé en anneaux locaux, \(S\) le site des ouverts de \(S\), \(\mathcal C\) la \(S\)-catégorie fibrée, et même scindée (SGA 1 VI 9), définie par le foncteur
\[
  U\longmapsto\{\text{catégorie des }\mathcal O_U\text{-modules à gauche}\};
\]
\(\mathcal C\) est une \(S\)-catégorie abélienne plate (1.1). Soit \(\mathcal C_0\) la sous-\(S\)-catégorie additive strictement pleine de \(\mathcal C\) définie par
\[
  U\longmapsto\{\text{catégorie des }\mathcal O_U\text{-modules libres de type fini}\}.
\]
Un \(\mathcal O_S\)-module \(M\) est localement de \(\mathcal C_0\)-type fini si et seulement si \(M\) est de type fini au sens ordinaire (EGA \(0_{\mathrm I}\), 5.2). De EGA \(0_{\mathrm I}\), 5.2.7 et 5.4.9, résulte que \(\mathcal C_0\) est localement stable par noyau d'épimorphisme et localement relevable dans \(\mathcal C\).

\textbf{c)} Soient \(G\) un groupe, \(S\) le topos des \(G\)-ensembles (SGA 4 V 3.11), \(k\) un corps, considéré comme anneau constant de \(S\), \(\mathcal C\) la \(S\)-catégorie scindée définie par le foncteur
\[
  U\longmapsto\{\text{catégorie des }k_U\text{-modules à gauche}\};
\]
\(\mathcal C\) est une \(S\)-catégorie abélienne plate. Soit \(\mathcal C_0\) la sous-\(S\)-catégorie additive strictement pleine de \(\mathcal C\) définie par
\[
  U\longmapsto\{\text{catégorie des }k_U\text{-modules libres de type fini}\}.
\]
Comme la localisation dans \(S\) équivaut à l'oubli des opérations de \(G\), il est clair qu'un \(G\)-\(k\)-module \(M\) est localement de \(\mathcal C_0\)-type fini si et seulement si \(M\) est de dimension finie comme espace vectoriel sur \(k\), et que \(\mathcal C_0\) est localement stable par noyau d'épimorphisme et localement relevable dans \(\mathcal C\).

On laisse au lecteur le soin de regarder la variante « galoisienne » de c), où \(G\) est un groupe profini et \(S\) le topos des ensembles où \(G\) opère continûment.

En fait, les exemples b) et c) sont des cas particuliers de l'exemple suivant.

\textbf{d)} Soient \((S,\mathcal O_S)\) un topos annelé (cf. SGA 4 IV 2.10), et \(\mathcal C\) la \(S\)-catégorie scindée définie par
\[
  U\longmapsto\{\text{catégorie des }\mathcal O_U\text{-modules à gauche}\};
\]
\(\mathcal C\) est une \(S\)-catégorie abélienne plate. Soit \(\mathcal C_0\) la sous-\(S\)-catégorie additive strictement pleine de \(\mathcal C\) définie par
\[
  U\longmapsto\{\text{catégorie des }\mathcal O_U\text{-modules libres de type fini}\}.
\]
Un \(\mathcal O_S\)-module localement de \(\mathcal C_0\)-type fini sera dit simplement de \emph{type fini}. La catégorie \(\mathcal C_0\) est localement quasi-stable par noyau d'épimorphisme et localement relevable dans \(\mathcal C\), comme il résulte du

\textbf{Lemme 1.3.1.} Soit \(P\) un \(\mathcal O_S\)-module. Conditions équivalentes :
\begin{enumerate}[label=(\roman*)]
\item \(P\) est de type fini et le foncteur \(\Hom(P,\cdot)\) est exact;
\item \(P\) est de type fini et tout épimorphisme \(M\to P\) admet localement une section;
\item \(P\) est localement facteur direct d'un module libre de type fini.
\end{enumerate}
La preuve est immédiate et laissée en exercice au lecteur.

\subsection*{1.4. Un lemme de « prolongement » (cf. EGA \(0_{\mathrm{III}}\), 11.9.1)}

\textbf{Proposition 1.4.1.} Soient \(S\) un site, \(\mathcal C\) une \(S\)-catégorie abélienne plate (1.1), \(\mathcal C_0\) une sous-\(S\)-catégorie additive strictement pleine de \(\mathcal C\); on suppose \(\mathcal C_0\) localement quasi-relevable dans \(\mathcal C\) (1.2). Soient \(S\in\Ob S\),
\[
  u:E\longrightarrow F
\]
un morphisme (de degré \(0\)) de complexes de \(\mathcal C_S\) (différentielles de degré \(+1\)), \(G\) le cône de \(u\) (cf. [V], chap. I, §1, 2.1), \(n\in\mathbb Z\); on suppose que \(\Ho^n(G)\) est de \(\mathcal C_0\)-type fini (1.2). Alors l'ensemble des \(X\in\Ob S/S\) vérifiant la propriété \((P)\) ci-dessous est un raffinement de \(S\).

\((P)\) : il existe un objet \(M\) de \(\mathcal C_{0,X}\) et des flèches de \(\mathcal C_X\)
\[
  r:M\longrightarrow Z^{n+1}(E_X)=
  \Ker\!\left(E_X^{n+1}\xrightarrow{d^{n+1}}E_X^{n+2}\right),
  \qquad
  s:M\longrightarrow F_X^n
\]
tels que le diagramme suivant de \(\mathcal C_X\)
\[
\begin{tikzcd}[column sep=large]
\cdots \arrow[r] & E_X^{n-1}
 \arrow[r,"{\binom{d_E^{n-1}}{0}}"]
 \arrow[d,"u_X^{n-1}"'] &
E_X^n\oplus M
 \arrow[r,"{(d_E^n,r)}"]
 \arrow[d,"{(u_X^n,s)}"'] &
E_X^{n+1} \arrow[r] \arrow[d,"u_X^{n+1}"] &
\cdots\\
\cdots \arrow[r] & F_X^{n-1} \arrow[r,"d_F^{n-1}"'] &
F_X^n \arrow[r,"d_F^n"'] &
F_X^{n+1} \arrow[r] &
\cdots
\end{tikzcd}
\]
soit un morphisme de complexes \(u':E'\to F_X\), c'est-à-dire que \(d_F^n s=u_X^{n+1}r\), et que
\begin{enumerate}[label=(\roman*)]
\item \(\Ho^{n+1}(u')\) soit un monomorphisme, avec
\[
  \Im\Ho^{n+1}(u')=\Im\Ho^{n+1}(u_X);
\]
\item \(\Ho^n(u')\) soit un épimorphisme.
\end{enumerate}

De manière imagée, on peut dire que, quitte à modifier localement \(E\) et \(u\), en degré \(n\) seulement, par l'addition à \(E^n\) d'un objet de \(\mathcal C_0\), on peut localement rendre \(\Ho^{n+1}(u)\) injectif, sans changer son image, et \(\Ho^n(u)\) surjectif.

\emph{Démonstration.} La marche est la suivante : on montre d'abord que l'on peut choisir localement \(r\) et \(s\) de manière à rendre \(\Ho^{n+1}(u)\) injectif; on vérifie que l'hypothèse de finitude sur \(\Ho^n(G)\) est conservée; une nouvelle correction locale en degré \(n\) permet alors d'obtenir \(\Ho^n(u)\) surjectif.

La suite exacte de cohomologie
\[
  \cdots\longrightarrow \Ho^n(G)\longrightarrow
  \Ho^{n+1}(E)\longrightarrow
  \Ho^{n+1}(F)\longrightarrow \cdots
\]
montre que le noyau de \(\Ho^{n+1}(u)\) est de \(\mathcal C_0\)-type fini. En d'autres termes, l'ensemble des \(X\) au-dessus de \(S\) tels qu'il existe une suite exacte
\[
  M\longrightarrow \Ho^{n+1}(E_X)
  \xrightarrow{\Ho^{n+1}(u)}
  \Ho^{n+1}(F_X),
  \qquad M\in\Ob\mathcal C_{0,X},
\]
est un raffinement de \(S\). Comme \(\mathcal C_0\) est localement quasi-relevable dans \(\mathcal C\), on peut d'ailleurs imposer que la flèche
\[
  M\longrightarrow \Ho^{n+1}(E_X)
\]
se factorise à travers \(Z^{n+1}(E_X)\) en
\[
  M\xrightarrow{r}Z^{n+1}(E_X)\longrightarrow \Ho^{n+1}(E_X).
\]
La flèche composée
\[
  M\xrightarrow{u_X^{n+1}r}Z^{n+1}(F_X)\longrightarrow \Ho^{n+1}(F_X)
\]
étant nulle par construction, \(u_X^{n+1}r\) se factorise à travers
\[
  B^{n+1}(F_X)=\Im\!\left(F_X^n\xrightarrow{d_F^n}F_X^{n+1}\right).
\]
Appliquant à nouveau la quasi-relevabilité, on peut imposer que \(u_X^{n+1}r\) se factorise à travers \(F_X^n\) en
\[
  M\xrightarrow{s}F_X^n\longrightarrow B^{n+1}(F_X)\longrightarrow Z^{n+1}(F_X).
\]
On vérifie alors trivialement que \(u'\), défini comme en (1.4.1), est un morphisme de complexes satisfaisant à (i).

Soit \(G'\) le cône de \(u':E'\to F_X\). Montrons que \(\Ho^n(G')\) est de \(\mathcal C_0\)-type fini. On a, par construction, une suite exacte
\[
  0\longrightarrow E_X\xrightarrow{v}E'\longrightarrow M[-n]\longrightarrow 0,
\]
avec \(u'v=u_X\). Pour comparer \(G_X\) et \(G'\), on peut, par exemple, former « l'octaèdre » de \(\Kpx(\mathcal C_X)\)
\[
\begin{tikzcd}[column sep=small,row sep=small]
& G_X \arrow[rr] \arrow[dr] && M[-n+1] \arrow[dl]\\
E_X \arrow[ur,"+1"] \arrow[rr,"v"'] && E' \arrow[ur] \arrow[rr,"u'"] && F_X .
\end{tikzcd}
\]
La suite exacte de cohomologie du triangle distingué marqué montre que
\[
  \Ho^n(G_X)\longrightarrow \Ho^n(G')
\]
est un épimorphisme, donc que \(\Ho^n(G')\) est bien de \(\mathcal C_0\)-type fini.

La correction précédente nous ramène donc au cas où \(\Ho^{n+1}(u)\) est un monomorphisme, ce que nous supposerons maintenant. La suite exacte de cohomologie de \(E\to F\to G\) donne alors une suite exacte
\[
  \Ho^n(E)\xrightarrow{\Ho^n(u)}
  \Ho^n(F)\longrightarrow \Ho^n(G)\longrightarrow 0.
\]
Comme \(\Ho^n(G)\) est de \(\mathcal C_0\)-type fini, l'ensemble des \(X\) au-dessus de \(S\) tels qu'il existe un épimorphisme
\[
  M\longrightarrow \Ho^n(G_X),\qquad M\in\Ob\mathcal C_{0,X},
\]
est, par définition, un raffinement de \(S\). Utilisant la quasi-relevabilité de \(\mathcal C_0\) dans \(\mathcal C\), on peut imposer que cet épimorphisme se factorise à travers \(Z^n(F_X)\) en
\[
  M\xrightarrow{s}Z^n(F_X)\longrightarrow \Ho^n(F_X)\longrightarrow \Ho^n(G_X).
\]
Il est clair que \(u'\), défini comme en (1.4.1) avec \(r=0\), est un morphisme de complexes et que
\[
  \Ho^n(u'):\Ho^n(E_X)\oplus M\longrightarrow \Ho^n(F_X)
\]
est un épimorphisme. Cela termine la démonstration de (1.4.1).

\textbf{Lemme 1.4.2.} Soient \(\mathcal C\) une catégorie abélienne, \(u:E\to F\) un morphisme de complexes de \(\mathcal C\) (resp. une flèche de \(\Dpx(\mathcal C)\)), \(G\) le cône de \(u\). Conditions équivalentes pour \(n\in\mathbb Z\) donné :
\begin{enumerate}[label=(\roman*)]
\item \(\Ho^i(G)=0\) pour \(i\geq n\);
\item \(\Ho^i(u)\) est un isomorphisme pour \(i>n\) et un épimorphisme pour \(i=n\).
\end{enumerate}
La preuve est immédiate à partir de la suite exacte de cohomologie.

\textbf{Définition 1.4.3.} Nous dirons que \(u\) est un \(n\)-quasi-isomorphisme (resp. un \(n\)-isomorphisme) si \(u\) vérifie les conditions équivalentes de (1.4.2).

Cela étant, on déduit aussitôt de (1.4.1) le

\textbf{Corollaire 1.4.4.} Soient \(S,S,\mathcal C,\mathcal C_0\) comme en (1.4.1). Soit \(u:E\to F\) un \((n+1)\)-quasi-isomorphisme de complexes de \(\mathcal C_S\) (\(n\in\mathbb Z\) donné) tel que \(\Ho^n(G)\) soit de \(\mathcal C_0\)-type fini. Alors l'ensemble des \(X\) au-dessus de \(S\) tels qu'il existe \((M,r,s)\) comme en (1.4.1) tels que \(u'\) soit un \(n\)-quasi-isomorphisme est un raffinement de \(S\).

De manière imagée, on peut, moyennant la condition de finitude sur \(\Ho^n(G)\), localement « gagner un cran », c'est-à-dire modifier localement \(u\) en degré \(n\) par l'addition à \(E^n\) d'un objet de \(\mathcal C_0\) de manière à faire de \(u\) un \(n\)-quasi-isomorphisme.

\textbf{Corollaire 1.4.5.} Soient \(\mathcal C\) une catégorie abélienne, \(n\in\mathbb Z\), et \(u:E\to F\) un \(n\)-quasi-isomorphisme de complexes de \(\mathcal C\). Il existe un diagramme commutatif de \(\Cpx(\mathcal C)\)
\[
\begin{tikzcd}
E \arrow[r,"v"] \arrow[d,"u"'] & E' \arrow[dl,"u'"]\\
F &
\end{tikzcd}
\]
où \(u'\) est un quasi-isomorphisme et \(v\) un monomorphisme induisant l'identité en degré \(\geq n\), c'est-à-dire tel que
\[
  E'^i=E^i,\qquad v^i=\id(E^i)\quad\text{pour }i\geq n.
\]

\emph{Preuve.} Appliquant (1.4.4) dans le cas où \(S\) est le site ponctuel et \(\mathcal C_0=\mathcal C\), on obtient un triangle commutatif de \(\Cpx(\mathcal C)\)
\[
\begin{tikzcd}
E \arrow[r,"v_{n-1}"] \arrow[d,"u"'] & E'_{\,n-1} \arrow[dl,"u'_{n-1}"]\\
F &
\end{tikzcd}
\]
où \(u'_{n-1}\) est un \((n-1)\)-quasi-isomorphisme et \(v_{n-1}\) un monomorphisme induisant l'identité en degré \(\geq n\). Itérant le procédé, on trouve, pour \(i\leq n-1\), une suite de triangles commutatifs de \(\Cpx(\mathcal C)\)
\[
\begin{tikzcd}
E'_i \arrow[r,"v_{i-1}"] \arrow[d,"u'_i"'] & E'_{\,i-1} \arrow[dl,"u'_{i-1}"]\\
F &
\end{tikzcd}
\]
où \(u'_i\) est un \(i\)-quasi-isomorphisme et \(v_{i-1}\) un monomorphisme induisant l'identité en degré \(\geq i\). Le corollaire s'en déduit par passage à la limite. De manière précise, notant \(d'_i\) la différentielle de \(E'_i\), on définit \(E'\) par
\[
  E'^i=E^i,\quad d'^i=d^i\quad (i\geq n),
\]
et
\[
  E'^i=(E'_i)^i,\quad d'^i=(d'_i)^i\quad (i<n),
\]
puis \(v\) et \(u'\) par
\[
  v^j=(v_i)^j,\qquad u'^j=(u'_i)^j\qquad (i\leq j),
\]
et l'on vérifie aussitôt que \(E'\), \(v\), \(u'\) ainsi définis répondent à la question.

Autrement dit, en abrégé, on peut prolonger un \(n\)-quasi-isomorphisme en un quasi-isomorphisme « sans rien changer à ce qui est donné en degré \(>n\) ».

\section*{2. Complexes pseudo-cohérents}

\textbf{2.0.} Dans ce numéro, \(S\) désigne un site, \(S\) un objet de \(S\), \(\mathcal C\) une \(S\)-catégorie abélienne plate (1.1), \(\mathcal C_0\) une sous-\(S\)-catégorie additive strictement pleine de \(\mathcal C\). On suppose que \(\mathcal C_0\) est localement quasi-stable par noyau d'épimorphisme (1.2) et localement quasi-relevable dans \(\mathcal C\) (1.2).

\textbf{Définition 2.1.} Soient \(X\in\Ob S\), \(E\) un complexe de \(\mathcal C_X\), et \(n\in\mathbb Z\). Nous dirons que \(E\) est strictement \(\mathcal C_0\)-\(n\)-pseudo-cohérent (resp. \(\mathcal C_0\)-pseudo-cohérent, resp. \(\mathcal C_0\)-parfait) si \(E^i=0\) pour \(i\) assez grand et \(E^i\in\Ob\mathcal C_{0,S}\) pour \(i\geq n\) (resp. pour tout \(i\), resp. pour tout \(i\) et \(E^i=0\) pour presque tout \(i\)).

Quand il n'y aura pas de confusion à craindre, nous dirons parfois « strictement \(n\)-pseudo-cohérent » (resp. ...) au lieu de « strictement \(\mathcal C_0\)-\(n\)-pseudo-cohérent ».

\textbf{Lemme 2.1.1.} Soient \(E\in\Ob\Cpx(\mathcal C_S)\) et \(n\in\mathbb Z\). Si \(E\) est strictement \(n\)-pseudo-cohérent et si \(\Ho^i(E)=0\) pour \(i>n\), alors \(\Ho^n(E)\) est de type fini.

\emph{Preuve.} La suite
\[
  0\longrightarrow Z^n(E)\longrightarrow E^n\longrightarrow E^{n+1}\longrightarrow\cdots\longrightarrow 0
\]
est exacte et \(E^i\in\Ob\mathcal C_0\) pour \(i\geq n\). Comme \(\mathcal C_0\) est localement quasi-stable par noyau d'épimorphisme, il s'ensuit que \(Z^n(E)\) est de type fini, donc aussi \(\Ho^n(E)\), qui en est un quotient.

\textbf{Proposition 2.2.} Soient \(F\in\Ob\Dpx(\mathcal C_S)\) et \(n\in\mathbb Z\). Alors :

\textbf{a)} Si \(F\) est strictement \(n\)-pseudo-cohérent (relativement à \(\mathcal C_0\)) et si \(u:E\to F\) est un quasi-isomorphisme, l'ensemble des \(X\) au-dessus de \(S\) tels qu'il existe un quasi-isomorphisme
\[
  v:F'\longrightarrow E_X
\]
avec \(F'\) strictement \(n\)-pseudo-cohérent est un raffinement de \(S\).

\textbf{b)} Les conditions suivantes sont équivalentes :
\begin{enumerate}[label=(\roman*)]
\item l'ensemble des \(X\) au-dessus de \(S\) tels qu'il existe un quasi-isomorphisme
\[
  u:E\longrightarrow F_X
\]
de \(\Cpx(\mathcal C_X)\), avec \(E\) strictement \(n\)-pseudo-cohérent, est un raffinement de \(S\);
\item[\((\mathrm{i}')\)] même condition que (i), mais \(u\) étant un isomorphisme de \(\Dpx(\mathcal C_X)\);
\item l'ensemble des \(X\) au-dessus de \(S\) tels qu'il existe un \(n\)-quasi-isomorphisme
\[
  u:E\longrightarrow F_X
\]
de \(\Cpx(\mathcal C_X)\), avec \(E\) strictement parfait, est un raffinement de \(S\);
\item[\((\mathrm{ii}')\)] même condition que (ii), mais \(u\) étant un \(n\)-isomorphisme de \(\Dpx(\mathcal C_X)\).
\end{enumerate}

\emph{Démonstration.} a) Il existe, pour un entier \(i\) assez grand, un \(i\)-quasi-isomorphisme
\[
  v_i:F'_i\longrightarrow E
\]
avec \(F'_i\) strictement \(n\)-pseudo-cohérent (prendre \(F'_i=0\)). Montrons que, si \(i>n\), il existe, quitte à raffiner \(S\), un \((i-1)\)-quasi-isomorphisme
\[
  v_{i-1}:F'_{i-1}\longrightarrow E
\]
avec \(F'_{i-1}\) strictement \(n\)-pseudo-cohérent. D'après (1.4.4), il suffit de montrer que \(\Ho^{i-1}(M')\), où \(M'\) est le cône de \(v_i\), est de \(\mathcal C_0\)-type fini. Or, si \(M''\) désigne le cône de \(uv_i\), l'hypothèse que \(u\) est un quasi-isomorphisme et (1.4.2) impliquent que l'on a
\[
  \Ho^j(M')\qis \Ho^j(M'')\qquad\text{pour tout }j,
\]
et
\[
  \Ho^j(M')=0\qquad\text{pour }j\geq i.
\]
La formule du cône
\[
  M''{}^j=F'_i{}^{\,j+1}\oplus F^j
\]
montre d'autre part que \(M''\) est strictement \((i-1)\)-pseudo-cohérent. Donc, en vertu de (2.1.1), \(\Ho^{i-1}(M'')=\Ho^{i-1}(M')\) est de type fini, d'où notre assertion. On en déduit, par itération, et quitte à raffiner \(S\), un \(n\)-quasi-isomorphisme
\[
  v_n:F'_n\longrightarrow E,
\]
où \(F'_n\) est strictement \(n\)-pseudo-cohérent. Mais, d'après (1.4.5), on peut prolonger \(v_n\) en un quasi-isomorphisme
\[
  v:F'\longrightarrow E
\]
sans rien changer en degré \(\geq n\), donc en particulier avec \(F'\) strictement \(n\)-pseudo-cohérent, ce qui prouve a).

b) L'équivalence des conditions (i) et \((\mathrm{i}')\) résulte aussitôt de a) et de la définition des isomorphismes de \(\Dpx^-(\mathcal C_X)\). Si \(E\in\Ob\Dpx^-(\mathcal C_S)\) est strictement \(n\)-pseudo-cohérent, il existe une suite exacte
\[
  0\longrightarrow E'\longrightarrow E\longrightarrow E''\longrightarrow 0,
\]
où \(E'\) est strictement parfait, et \(E''{}^i=0\) pour \(i\geq n\); par suite \(E'\to E\) est un \(n\)-quasi-isomorphisme, d'où l'implication \((\mathrm{i})\Rightarrow(\mathrm{ii})\). L'implication \((\mathrm{ii})\Rightarrow(\mathrm{ii}')\) étant triviale, il reste à prouver \((\mathrm{ii}')\Rightarrow(\mathrm{i})\). Soit donc
\[
  u:E\longrightarrow F
\]
un \(n\)-isomorphisme de \(\Dpx^-(\mathcal C_S)\), avec \(E\) strictement parfait; \(u\) est donc défini par un diagramme de \(\Cpx^-(\mathcal C_S)\)
\[
\begin{tikzcd}
& E' \arrow[dl,"s"'] \arrow[dr,"u'"]\\
E && F ,
\end{tikzcd}
\]
où \(s\) est un quasi-isomorphisme et \(u'\) un \(n\)-quasi-isomorphisme. En vertu de a), il existe, quitte à raffiner \(S\), un quasi-isomorphisme
\[
  t:E''\longrightarrow E',
\]
avec \(E''\) strictement \(n\)-pseudo-cohérent. On applique alors (1.4.5) pour prolonger \(u't:E''\to F\) en un quasi-isomorphisme
\[
  E'''\longrightarrow F
\]
avec \(E'''\) strictement \(n\)-pseudo-cohérent, ce qui prouve \((\mathrm{ii}')\Rightarrow(\mathrm{i})\) et achève la démonstration de (2.2).

\textbf{Définition 2.3.} Sous les hypothèses de (2.2), nous dirons que \(F\in\Ob\Dpx^-(\mathcal C_S)\) est \(n\)-\(\mathcal C_0\)-pseudo-cohérent si \(F\) vérifie les conditions équivalentes b) de (2.2). Nous dirons que \(F\) est \(\mathcal C_0\)-pseudo-cohérent si \(F\) est \(n\)-\(\mathcal C_0\)-pseudo-cohérent pour tout \(n\).

Quand cela ne pourra prêter à confusion, nous nous permettrons parfois de dire \(n\)-pseudo-cohérent, pseudo-cohérent, au lieu de \(n\)-\(\mathcal C_0\)-pseudo-cohérent, \(\mathcal C_0\)-pseudo-cohérent.

\textbf{Remarque 2.3.1.} Soit \(F\in\Ob\Dpx(\mathcal C_S)\) un complexe pseudo-cohérent. L'ensemble des \(X\) au-dessus de \(S\) tels qu'il existe un quasi-isomorphisme
\[
  E\longrightarrow F_X
\]
avec \(E\) strictement pseudo-cohérent (2.1) n'est pas en général un raffinement de \(S\). Toutefois, nous verrons plus tard (2.7 a) et Exp. II) des cas importants où il en est ainsi.

La démonstration de (2.2 a) donne le résultat suivant :

\textbf{Proposition 2.4.} Soient \(n,i\in\mathbb Z\), et soit \(u:E\to F\) un \(i\)-quasi-isomorphisme de \(\Cpx(\mathcal C_S)\), avec \(F\) \(n\)-pseudo-cohérent et \(E\) strictement \(n\)-pseudo-cohérent. Alors, quitte à raffiner \(S\), il existe un triangle commutatif de \(\Cpx(\mathcal C_S)\)
\[
\begin{tikzcd}
E \arrow[r,"v"] \arrow[d,"u"'] & E' \arrow[dl,"u'"]\\
F &
\end{tikzcd}
\]
où \(E'\) est strictement \(n\)-pseudo-cohérent, \(u'\) un quasi-isomorphisme et \(v\) un monomorphisme induisant l'identité en degré \(\geq i\).

\textbf{Proposition 2.5.} \textbf{a)} Si \(F\in\Ob\Dpx(\mathcal C_S)\) est \(n\)-pseudo-cohérent (resp. pseudo-cohérent), \(F[m]\) est \((n-m)\)-pseudo-cohérent (resp. pseudo-cohérent).

\textbf{b)} Soit
\[
\begin{tikzcd}
E \arrow[rr] && F \arrow[dl]\\
& G \arrow[ul,"+1"] &
\end{tikzcd}
\tag{*}
\]
un triangle distingué de \(\Dpx(\mathcal C_S)\). Si \(E\) et \(G\) sont \(n\)-pseudo-cohérents (resp. pseudo-cohérents), alors \(F\) est \(n\)-pseudo-cohérent (resp. pseudo-cohérent).

\emph{Preuve.} a) Trivial à partir des définitions.

b) Il suffit bien entendu de démontrer l'assertion non respée. Le triangle \((*)\) peut encore s'écrire
\[
\begin{tikzcd}
G[-1] \arrow[rr] && E \arrow[dl]\\
& F \arrow[ul,"+1"] &
\end{tikzcd}
\]
D'après a), \(G[-1]\) est \((n+1)\)-pseudo-cohérent. On est donc ramené à prouver que, si \(E\) est \((n+1)\)-pseudo-cohérent et \(F\) \(n\)-pseudo-cohérent, alors \(G\) est \(n\)-pseudo-cohérent. Quitte à se localiser, on peut supposer \(E\) (resp. \(F\)) strictement \((n+1)\)-pseudo-cohérent (resp. \(n\)-pseudo-cohérent). La flèche \(E\xrightarrow{u}F\) de \(\Dpx^-(\mathcal C_S)\) est définie par un couple de flèches de \(\Cpx^-(\mathcal C_S)\)
\[
\begin{tikzcd}
& E' \arrow[dl,"s"'] \arrow[dr,"u'"]\\
E && F ,
\end{tikzcd}
\]
où \(s\) est un quasi-isomorphisme. D'après (2.2 a), il existe, quitte à localiser, un quasi-isomorphisme \(t:E''\to E'\), où \(E''\) est strictement \((n+1)\)-pseudo-cohérent. On est donc réduit au cas où \(u\) est définie par une vraie flèche de \(\Cpx^-(\mathcal C_S)\). L'assertion résulte alors des définitions et de la formule
\[
  C(u)^i=E^{i+1}\oplus F^i,
\]
où \(C(u)\) désigne le cône de \(u\).

\textbf{Scholie 2.6.} Par « rotation » des triangles, on voit que des propriétés de pseudo-cohérence sur deux des objets du triangle \((*)\) en impliquent une autre sur le troisième. On peut résumer les différents cas dans le tableau ci-dessous (les petites lettres désignent des degrés de pseudo-cohérence et chaque ligne correspond à une implication dont le but est souligné) :
\[
\begin{array}{ccc}
E&F&G\\[0.2em]
n+1&n&\underline n\\
n&\underline n&n\\
\underline n&n&n-1 .
\end{array}
\]

Nous noterons
\[
  \Dpx(\mathcal C)_{\mathcal C_0\text{-}n\mathrm{coh}}
  \quad\text{(resp. }\Dpx(\mathcal C)_{\mathcal C_0\text{-}\mathrm{coh}}\text{)}
\]
la sous-catégorie fibrée strictement pleine de \(\Dpx(\mathcal C)\) dont la fibre en chaque \(X\in\Ob S\) est formée des objets de \(\Dpx(\mathcal C_X)\) qui sont \(n\)-pseudo-cohérents (resp. pseudo-cohérents) relativement à \(\mathcal C_0\). La catégorie
\[
  \Dpx(\mathcal C)_{\mathcal C_0\text{-}\mathrm{coh}}
\]
est, en vertu de (2.5), une sous-\(S\)-catégorie triangulée de \(\Dpx(\mathcal C)\) (1.1).

\textbf{Proposition 2.7.} Supposons que \(S\) soit un site ponctuel, donc
\[
  \mathcal C=\mathcal C_S,\qquad \mathcal C_0=\mathcal C_{0,S}.
\]
Alors :
\begin{enumerate}[label=\alph*)]
\item Si \(F\in\Ob\Dpx(\mathcal C)\) est pseudo-cohérent, il existe un quasi-isomorphisme
\[
  E\longrightarrow F,
\]
avec \(E\) strictement pseudo-cohérent (2.1).
\item Désignons par \(\Kpx^{-,\mathrm{ac}}(\mathcal C_0)\) la sous-catégorie triangulée pleine de \(\Kpx^-(\mathcal C_0)\) formée des complexes acycliques. Le foncteur canonique
\[
  \Kpx^-(\mathcal C_0)/\Kpx^{-,\mathrm{ac}}(\mathcal C_0)
  \longrightarrow \Dpx(\mathcal C)
\]
est pleinement fidèle et a pour image essentielle \(\Dpx(\mathcal C)_{\mathcal C_0\text{-}\mathrm{coh}}\).
\end{enumerate}

\emph{Preuve.} a) Soit \(F\in\Ob\Dpx(\mathcal C)_{\mathcal C_0\text{-}\mathrm{coh}}\), et supposons construit un \(i\)-quasi-isomorphisme de \(\Cpx(\mathcal C)\)
\[
  u_i:E_i\longrightarrow F
\]
avec \(E_i\) strictement pseudo-cohérent; c'est possible pour \(i\) assez grand avec \(E_i=0\). Soit \(G_i\) le cône de \(u_i\). Il résulte de (1.4.2) et (2.1.1) que \(\Ho^{i-1}(G_i)\) est de \(\mathcal C_0\)-type fini. Donc, en vertu de (1.4.4), il existe un triangle commutatif de \(\Cpx(\mathcal C)\)
\[
\begin{tikzcd}
E_i \arrow[r,"v_{i-1}"] \arrow[d,"u_i"'] & E_{i-1} \arrow[dl,"u_{i-1}"]\\
F &
\end{tikzcd}
\]
où \(E_{i-1}\) est strictement pseudo-cohérent, \(u_{i-1}\) un \((i-1)\)-quasi-isomorphisme, et \(v_{i-1}\) un monomorphisme induisant l'identité en degré \(\geq i\). On en déduit par passage à la limite (cf. preuve de (1.4.5)) un quasi-isomorphisme \(u:E\to F\) avec \(E\) strictement pseudo-cohérent.

b) Appliquer le critère \([V]\), chap. I, §2, 4.2 b)(i), et a).

Noter que la preuve de a) donne le résultat plus précis suivant, qui généralise (1.4.5) :

\textbf{Proposition 2.7.1.} Sous les hypothèses de (2.7), soit \(u:E\to F\) un \(n\)-quasi-isomorphisme de \(\Cpx(\mathcal C)\), avec \(F\) pseudo-cohérent et \(E\) strictement pseudo-cohérent. Il existe alors un triangle commutatif de \(\Cpx(\mathcal C)\)
\[
\begin{tikzcd}
E \arrow[r,"v"] \arrow[d,"u"'] & E' \arrow[dl,"u'"]\\
F &
\end{tikzcd}
\]
où \(E'\) est strictement pseudo-cohérent, \(u'\) un quasi-isomorphisme, et \(v\) un monomorphisme induisant l'identité en degré \(\geq n\).

\textbf{Définition 2.8.} Soient \(n\in\mathbb N\) et \(F\in\Ob\mathcal C_S\). Nous dirons que \(F\) est de \(n\)-\(\mathcal C_0\)-présentation finie si l'ensemble des \(X\) au-dessus de \(S\) tels qu'il existe une suite exacte
\[
  E^{-n}\longrightarrow\cdots\longrightarrow E^0\longrightarrow F_X\longrightarrow 0,
\qquad E^i\in\Ob\mathcal C_{0,X}\quad(-n\leq i\leq 0),
\]
est un raffinement de \(S\). Nous dirons que \(F\) est d'\(\infty\)-\(\mathcal C_0\)-présentation finie si \(F\) est de \(n\)-\(\mathcal C_0\)-présentation finie pour tout \(n\in\mathbb N\).

Bien entendu, nous dirons parfois \(n\)-présentation finie tout court quand il n'y aura pas d'ambiguïté possible sur \(\mathcal C_0\). Nous dirons « de présentation finie » pour « de 1-présentation finie », et « de type fini » pour « de 0-présentation finie », cette dernière notion coïncidant avec celle déjà introduite en (1.2).

\textbf{Proposition 2.9.} Soient \(n\in\mathbb N\) et \(F\in\Ob\mathcal C_S\). Conditions équivalentes :
\begin{enumerate}[label=(\roman*)]
\item \(F\) est de \(n\)-présentation finie (resp. d'\(\infty\)-présentation finie);
\item \(F\) est \((-n)\)-pseudo-cohérent (resp. pseudo-cohérent).
\end{enumerate}
Dans le cas où \(S\) est un site ponctuel, \(F\) est pseudo-cohérent si et seulement s'il existe une suite exacte
\[
  \cdots\longrightarrow E^i\longrightarrow\cdots\longrightarrow E^0\longrightarrow F\longrightarrow 0,
\qquad E^i\in\Ob\mathcal C_0\quad\text{pour tout }i.
\]

\emph{Preuve.} Pour l'équivalence \((\mathrm{i})\Leftrightarrow(\mathrm{ii})\), il suffit évidemment de traiter le cas non respé. L'implication \((\mathrm{i})\Rightarrow(\mathrm{ii})\) est immédiate sur les définitions, et l'implication \((\mathrm{ii})\Rightarrow(\mathrm{i})\) résulte aussitôt de (2.4) : la flèche \(0\to F\) est un 1-quasi-isomorphisme, on en déduit, quitte à localiser, un quasi-isomorphisme
\[
  E\longrightarrow F
\]
avec \(E\) strictement \((-n)\)-pseudo-cohérent et \(E^i=0\) pour \(i\geq 1\), c'est-à-dire une \(n\)-présentation finie de \(F\). La deuxième assertion se démontre de manière analogue (utiliser (2.7.1)).

\textbf{Proposition 2.10.} \textbf{a)} Soient \(m,n\in\mathbb Z\), avec \(n\leq m\). Si \(F\in\Ob\Dpx^-(\mathcal C_S)\) est \(n\)-pseudo-cohérent et acyclique en degré \(\geq m\), il existe localement un quasi-isomorphisme
\[
  E\longrightarrow F,
\]
avec \(E\) strictement \(n\)-pseudo-cohérent et \(E^i=0\) pour \(i\geq m\). Si \(S\) est ponctuel, et \(F\) pseudo-cohérent et acyclique en degré \(\geq m\), il existe un quasi-isomorphisme
\[
  E\longrightarrow F
\]
avec \(E\) strictement pseudo-cohérent et \(E^i=0\) pour \(i\geq m\).

\textbf{b)} Soit \(F\in\Ob\Dpx^-(\mathcal C_S)\) acyclique en degré \(>n\). Pour que \(F\) soit \(n\)-pseudo-cohérent, il faut et il suffit que \(\Ho^n(F)\) soit de type fini.

\emph{Preuve.} a) Comme \(F\) est acyclique en degré \(\geq m\), la flèche \(0\to F\) est un \(m\)-quasi-isomorphisme. L'assertion résulte de (2.4) et (2.7.1).

b) Il existe un triangle distingué de \(\Dpx^-(\mathcal C_S)\) :
\[
\begin{tikzcd}
F \arrow[rr] && F' \arrow[dl]\\
& \Ho^n(F)[-n] \arrow[ul,"+1"] &
\end{tikzcd}
\]
(pour le voir, commencer par tronquer \(F\) en remplaçant \(F^n\) par \(Z^n(F)\) et \(F^i\) par \(0\) pour \(i>n\)). La suite exacte de cohomologie implique \(\Ho^i(F')=0\) pour \(i\geq n-1\), donc \(F'\) est \((n-1)\)-pseudo-cohérent. On conclut grâce à (2.5) et (2.6).

\end{document}
